\documentclass[11pt, reqno]{amsart}

% The abridged version of main.tex: statements and sketches of proofs, with pointers to the long version.
% It reuses preamble.tex, macros.tex and refs.bib of the long version without modifying them.
% The section files are in summary/; summary/longrefs.tex is generated from main.aux of the long version
% and holds the printed numbers of its labels, which \Lr and \LL turn into references "[L, Theorem 3.8]".
% The preamble: packages, page layout, draft notes, theorem environments, cross-reference names and generic letters.
% Loaded by the main file before the notation. Only the packages loaded here are available to the sections.

% GENUS workshop style: page geometry, fonts, hyperref, fancyhdr and the claim-status tags (genus-collab/meta)
\usepackage{genus}
% text width of the original note (170 mm), so that its displayed formulas still fit the line
\geometry{margin=2cm}
\usepackage[utf8]{inputenc}
\usepackage[T1]{fontenc}

%-------------------------------------------------------------------------------
% Packages
\usepackage[dvipsnames]{xcolor}
\usepackage{amssymb}
\usepackage[mathscr]{eucal}
\usepackage{amsthm}
\usepackage{amsmath}
\usepackage{mathtools}
\definecolor{highlight}{HTML}{3465a4}
% hyperref is loaded by genus.sty, with the GENUS link colours
\hypersetup{hypertexnames=false, hyperfootnotes=false, linktocpage=true, bookmarksdepth=2}
\usepackage{cleveref}
\usepackage{url}
\usepackage{array}

%-------------------------------------------------------------------------------
% Typefaces and layout
% (mathpazo and courier dropped: genus.sty sets Latin Modern)
\usepackage{microtype}
\linespread{1.06}
\setlength{\parskip}{4pt}
\setlength{\parindent}{0pt}
\setcounter{secnumdepth}{3}
% enumitem is not assumed: first-level enumerate is labelled (i), (ii), ...; for named parts use
% description or \item[(a)] with an explicit label.
\renewcommand{\labelenumi}{(\roman{enumi})}
\renewcommand{\theenumi}{(\roman{enumi})}
\numberwithin{equation}{section}

%-------------------------------------------------------------------------------
% Draft notes: visible, coloured, usable in text and in math. Set \showdraftnotesfalse to suppress them.
\newif\ifshowdraftnotes
\showdraftnotestrue
\newcommand{\draftnotetext}[1]{\textcolor{red}{\footnotesize[\textbf{Draft note:} #1]}}
\newcommand{\draftnote}[1]{\ifshowdraftnotes\ifmmode\text{\draftnotetext{#1}}\else\draftnotetext{#1}\fi\fi}

%-------------------------------------------------------------------------------
% Theorem environments
\theoremstyle{plain}
\newtheorem{thm}{Theorem}[section]
\newtheorem{lem}[thm]{Lemma}
\newtheorem{prop}[thm]{Proposition}
\newtheorem{cor}[thm]{Corollary}
\theoremstyle{definition}
\newtheorem{defn}[thm]{Definition}
\newtheorem{rmk}[thm]{Remark}
\newtheorem{notation}[thm]{Notation}
% genus.sty defines an example environment on its own counter; redefine it on the counter of thm
\let\example\undefined\let\endexample\undefined\expandafter\let\csname c@example\endcsname\undefined
\newtheorem{example}[thm]{Example}
\newtheorem{conv}[thm]{Convention}

\crefname{equation}{Eq.}{Eqs.}
\crefname{thm}{Theorem}{Theorems}
\crefname{lem}{Lemma}{Lemmas}
\crefname{prop}{Proposition}{Propositions}
\crefname{cor}{Corollary}{Corollaries}
\crefname{defn}{Definition}{Definitions}
\crefname{rmk}{Remark}{Remarks}
\crefname{notation}{Notation}{Notations}
\crefname{example}{Example}{Examples}
\crefname{conv}{Convention}{Conventions}
\crefname{section}{Section}{Sections}
\crefname{subsection}{Section}{Sections}
\crefname{appendix}{Appendix}{Appendices}
\crefname{table}{Table}{Tables}

%-------------------------------------------------------------------------------
% Generic letters
\newcommand{\bbA}{\mathbb{A}}
\newcommand{\bbC}{\mathbb{C}}
\newcommand{\bbD}{\mathbb{D}}
\newcommand{\bbE}{\mathbb{E}}
\newcommand{\bbF}{\mathbb{F}}
\newcommand{\bbK}{\mathbb{K}}
\newcommand{\bbL}{\mathbb{L}}
\newcommand{\bbP}{\mathbb{P}}
\newcommand{\bbQ}{\mathbb{Q}}
\newcommand{\bbR}{\mathbb{R}}
\newcommand{\bbT}{\mathbb{T}}
\newcommand{\bbZ}{\mathbb{Z}}
\newcommand{\cA}{\mathcal{A}}
\newcommand{\cC}{\mathcal{C}}
\newcommand{\cE}{\mathcal{E}}
\newcommand{\cF}{\mathcal{F}}
\newcommand{\cG}{\mathcal{G}}
\newcommand{\cL}{\mathcal{L}}
\newcommand{\cN}{\mathcal{N}}
\newcommand{\cO}{\mathcal{O}}
\newcommand{\cP}{\mathcal{P}}
\newcommand{\cQ}{\mathcal{Q}}
\newcommand{\cW}{\mathcal{W}}
\newcommand{\cZ}{\mathcal{Z}}
\newcommand{\mgp}[1]{\mathsf{#1}}

%-------------------------------------------------------------------------------
% General
\newcommand{\dd}{\mathrm{d}}
\newcommand{\ri}{\mathrm{i}}
\newcommand{\re}{\mathrm{e}}
\newcommand{\de}{\partial}
% \llbracket: from a font package if available, otherwise a substitute (stmaryrd is not assumed)
\providecommand{\llbracket}{[\![}
\providecommand{\rrbracket}{]\!]}
\newcommand{\fps}[1]{\llbracket #1 \rrbracket}
\newcommand{\lser}[1]{(\!(#1)\!)}
\newcommand{\Frac}{\mathrm{Frac}}
\DeclareMathOperator{\Tr}{Tr}
\DeclareMathOperator{\tr}{tr}
\DeclareMathOperator{\End}{End}
\DeclareMathOperator{\Aut}{Aut}
\DeclareMathOperator*{\Res}{Res}
\DeclareMathOperator{\val}{val}
\DeclareMathOperator{\ev}{ev}
\DeclareMathOperator{\diag}{diag}
\DeclareMathOperator{\ord}{ord}
\newcommand{\tensor}{\otimes}
\newcommand{\coef}[1]{[#1]}                    % coefficient extraction, e.g. \coef{z^n}

% The notation of the note. Loaded by the main file after the preamble.
% A section file defines a new macro only if it is unavoidable, at its top, with \providecommand and the section prefix.

%-------------------------------------------------------------------------------
% Variables, points, rings
\newcommand{\Dq}{D}                            % q d/dq = -(1-u) d/du
\newcommand{\zh}{\hat{z}}                      % z/s^4 = L^4 z
\newcommand{\tc}{t_{\mathrm{c}}}               % varpi_1/varpi_0
\newcommand{\Cstar}{\mathbb{C}_{\ast}}         % C minus {0, 1, zeta, ..., zeta^4}
\newcommand{\AL}{\mathbb{A}_{L}}               % Q[L, L^{-1}, (L^5 - 1)^{-1}]
\newcommand{\bbH}{\mathbb{H}}                  % the state space, flat basis H^a
\newcommand{\Qn}{\mathscr{Q}}                  % the degree variable of F_g
\newcommand{\logq}{\mathfrak{l}}               % the formal log q
\newcommand{\Nring}{\mathscr{N}}               % the presented differential ring
\newcommand{\Cexp}{\mathscr{C}}                % the conifold expansion \Cexp_{\alpha,\beta}
\DeclareMathOperator{\Wr}{Wr}                  % Wronskian
\DeclareMathOperator{\Gau}{Gau}                % solution module of the gauge equation
\DeclareMathOperator{\wt}{wt}                  % weight
\DeclareMathOperator{\Hess}{Hess}              % Hessian

%-------------------------------------------------------------------------------
% The formal quintic, rows, (QRR)
\newcommand{\PFfour}{\mathrm{PF}_{4}}
\newcommand{\PFfive}{\mathrm{PF}_{5}}
\newcommand{\TFT}{\boldsymbol{\omega}}         % TFT (\TFT, \hat\TFT, \TFT^{\mathrm m}_v)
\newcommand{\Qr}{\mathcal{Q}}                  % rows Q^{(i)}_a
\newcommand{\Phiact}{\Phi^{\mathrm{act}}}
\newcommand{\Phican}{\Phi^{\mathrm{can}}}
\newcommand{\PhiSP}{\Phi^{\mathrm{SP}}}
\newcommand{\bk}{\mathsf{b}}                   % \bk_k = B_{2k}/(2k(2k-1))
\newcommand{\bw}{\mathbf{w}}                   % \bw_i
\newcommand{\hw}{h_{\bw}}                      % h_{\bw}
\newcommand{\ttr}{\mathfrak{t}}                % PPZ translation \ttr_R(z) = z1 - zR^{-1}(z)1

%-------------------------------------------------------------------------------
% Cyclic systems, gauges, normalisation
\newcommand{\Bprod}{\mathcal{B}}               % prod_a b_a
\newcommand{\Pc}{P_{\mathrm{can}}}             % the canonical gauge
\newcommand{\Dn}{\mathsf{D}}                   % normalisation series
\newcommand{\cn}{\mathsf{c}}                   % freedom P -> P c(z)
\newcommand{\Ghat}{\hat{G}}
\newcommand{\Rhat}{\hat{R}}
\newcommand{\YC}{Y^{\mathrm{C}}}
\newcommand{\YCQ}{Y^{\mathrm{C}}_{\mathrm{Q}}}
\newcommand{\IC}{I^{\mathrm{C}}}
\newcommand{\Conc}{\mathcal{C}}                % concomitant
\newcommand{\Jet}{\mathcal{J}}                 % jet matrix
\newcommand{\Mcyc}{M_{\mathrm{cyc}}}
\newcommand{\sh}{\mathrm{sh}}                  % the automorphism ell -> ell + 1
\newcommand{\fc}{\mathfrak{f}}                 % free coefficients \fc_m
\newcommand{\Fg}[1]{(\mathrm{F}_{#1})}         % the window (F_g)

%-------------------------------------------------------------------------------
% The model, transport, topological recursion
\newcommand{\Wv}{W_{v}}                        % p^4/4 + v/p
\newcommand{\Mo}{M_{0}}                        % the model R-matrix in Teleman form
\newcommand{\Fs}{\mathsf{F}}                   % model row series \Fs_b(t)
\newcommand{\Fm}{F^{\mathrm{m}}}               % model free energy
\newcommand{\FTR}{F^{\mathrm{TR}}}             % TR free energy
\newcommand{\cH}{\mathcal{H}}                  % loop spaces
\newcommand{\cR}{\mathcal{R}}                  % weighted rings (appendix A)
\newcommand{\cV}{\mathcal{V}}                  % the residue (Section 8)
\newcommand{\cX}{\mathcal{X}}                  % the variation field (Section 6)
\newcommand{\sK}{\mathcal{K}}                  % the slice K
\newcommand{\sKo}{\mathcal{K}_{0}}             % its kernel part K_0
\newcommand{\fr}{\mathrm{frame}}               % \fr(k,a)
\newcommand{\se}{\mathsf{e}}                   % Euler-dilaton vector z1 + vH
\newcommand{\ML}[1]{\mathrm{ML}_{\mathcal{K}}(#1)}  % the Model Lemma on K at genus #1

%-------------------------------------------------------------------------------
% Connection problem, Gamma identity, formal stationary phase
\newcommand{\fa}{\mathfrak{a}}                 % amplitude
\newcommand{\SPn}{\mathrm{SP}^{\circ}}         % normalised stationary phase
\newcommand{\SPh}{\widehat{\mathrm{SP}}}       % divided by its constant term
\newcommand{\Gst}{\Gamma_{\mathrm{St}}}        % formal Stirling form
\newcommand{\Ginf}{G_{\infty}}                 % the orbifold value

\showdraftnotesfalse

%-------------------------------------------------------------------------------
% References to the long version, denoted [L]
\makeatletter
\newcommand{\LongLabel}[2]{\expandafter\gdef\csname LongRef@#1\endcsname{#2}}
\newcommand{\LL}[1]{\@ifundefined{LongRef@#1}{\PackageError{summary}{Unknown label #1 of the long version}{}\textbf{??}}{\csname LongRef@#1\endcsname}}
\newcommand{\Lr}[1]{[L\@for\Lr@tmp:=#1\do{, \expandafter\LL\expandafter{\Lr@tmp}}]}
\makeatother
% Generated from main.aux of the long version (label -> printed number). Do not edit by hand.
% Each entry: \LongLabel{label of the long version}{its printed name and number}.
\LongLabel{sec:int}{Section~1}\LongLabel{sec:int-setting}{Section~1.1}
\LongLabel{sec:int-gap}{Section~1.2}\LongLabel{sec:int-results}{Section~1.3}
\LongLabel{thm:int-main}{Theorem~1.1}\LongLabel{cor:int-gap}{Corollary~1.2}
\LongLabel{thm:int-D}{Theorem~1.3}\LongLabel{thm:int-gamma}{Theorem~1.4}
\LongLabel{rmk:int-constant}{Remark~1.5}\LongLabel{sec:int-inputs}{Section~1.4}
\LongLabel{sec:int-notclaimed}{Section~1.5}\LongLabel{sec:int-strategy}{Section~1.6}
\LongLabel{sec:int-plan}{Section~1.7}\LongLabel{sec:int-ack}{Section~1.7}
\LongLabel{sec:set}{Section~2}\LongLabel{sec:set-fq}{Section~2.1}
\LongLabel{def:set-fq}{Definition~2.1}\LongLabel{def:set-genus0}{Definition~2.2}
\LongLabel{eq:set-PF5}{(2.1)}\LongLabel{eq:set-PF4}{(2.2)}
\LongLabel{lem:set-pf5}{Lemma~2.3}\LongLabel{thm:set-mirror}{Theorem~2.4}
\LongLabel{thm:set-G0}{Theorem~2.5}\LongLabel{prop:set-cyclicLR}{Proposition~2.6}
\LongLabel{cor:set-cyclicPF5}{Corollary~2.7}\LongLabel{def:set-points}{Definition~2.8}
\LongLabel{def:set-frob}{Definition~2.9}\LongLabel{sec:set-rows}{Section~2.2}
\LongLabel{def:set-rows}{Definition~2.10}\LongLabel{thm:set-qrr}{Theorem~2.11}
\LongLabel{eq:set-hw}{(2.3)}\LongLabel{eq:set-qrr}{(2.4)}
\LongLabel{rmk:set-qrrsign}{Remark~2.12}\LongLabel{lem:set-hrot}{Lemma~2.13}
\LongLabel{lem:set-Runique}{Lemma~2.14}\LongLabel{sec:set-graph}{Section~2.3}
\LongLabel{def:set-graphsum}{Definition~2.15}\LongLabel{lem:set-pushins}{Lemma~2.16}
\LongLabel{thm:set-GT}{Theorem~2.17}\LongLabel{sec:set-frames}{Section~2.4}
\LongLabel{lem:set-frobbasis}{Lemma~2.18}\LongLabel{def:set-frames}{Definition~2.19}
\LongLabel{def:set-confexp}{Definition~2.20}\LongLabel{def:set-tau}{Definition~2.21}
\LongLabel{lem:set-canon}{Lemma~2.22}\LongLabel{sec:set-model}{Section~2.5}
\LongLabel{def:set-model}{Definition~2.23}\LongLabel{prop:set-modrows}{Proposition~2.24}
\LongLabel{def:set-M0}{Definition~2.25}\LongLabel{lem:set-M0sympl}{Lemma~2.26}
\LongLabel{def:set-modcohft}{Definition~2.27}\LongLabel{lem:set-modhom}{Lemma~2.28}
\LongLabel{lem:set-dilaton}{Lemma~2.29}\LongLabel{sec:set-hat}{Section~2.6}
\LongLabel{def:set-hat}{Definition~2.30}\LongLabel{def:set-outer}{Definition~2.31}
\LongLabel{sec:set-N}{Section~2.7}\LongLabel{def:set-N}{Definition~2.32}
\LongLabel{eq:set-cg}{(2.5)}\LongLabel{rmk:set-cg}{Remark~2.33}
\LongLabel{sec:gau}{Section~3}\LongLabel{sec:gau-cyc}{Section~3.1}
\LongLabel{def:gau-conc}{Definition~3.1}\LongLabel{eq:gau-Mcyc}{(3.1)}
\LongLabel{lem:gau-confexp}{Lemma~3.2}\LongLabel{thm:gau-cyc}{Theorem~3.3}
\LongLabel{eq:gau-cyc}{(3.2)}\LongLabel{cor:gau-cycsys}{Corollary~3.4}
\LongLabel{eq:gau-system}{(3.3)}\LongLabel{sec:gau-systems}{Section~3.2}
\LongLabel{def:gau-onepole}{Definition~3.5}\LongLabel{prop:gau-quintic}{Proposition~3.6}
\LongLabel{prop:gau-equalprod}{Proposition~3.7}\LongLabel{sec:gau-U}{Section~3.3}
\LongLabel{thm:gau-U}{Theorem~3.8}\LongLabel{lem:gau-comm}{Lemma~3.9}
\LongLabel{lem:gau-lift}{Lemma~3.10}\LongLabel{lem:gau-scal}{Lemma~3.11}
\LongLabel{eq:gau-reduced}{(3.4)}\LongLabel{lem:gau-lead}{Lemma~3.12}
\LongLabel{lem:gau-degree}{Lemma~3.13}\LongLabel{prop:gau-graded}{Proposition~3.14}
\LongLabel{sec:gau-van}{Section~3.4}\LongLabel{lem:gau-R1}{Lemma~3.15}
\LongLabel{lem:gau-diag}{Lemma~3.16}\LongLabel{def:gau-van}{Definition~3.17}
\LongLabel{lem:gau-vanmod}{Lemma~3.18}\LongLabel{lem:gau-vanshape}{Lemma~3.19}
\LongLabel{lem:gau-uniq}{Lemma~3.20}\LongLabel{lem:gau-vanq}{Lemma~3.21}
\LongLabel{thm:gau-G}{Theorem~3.22}\LongLabel{eq:gau-G}{(3.5)}
\LongLabel{def:gau-can}{Definition~3.23}\LongLabel{prop:gau-symp}{Proposition~3.24}
\LongLabel{eq:gau-symp}{(3.6)}\LongLabel{sec:gau-R0}{Section~3.5}
\LongLabel{lem:gau-HPF}{Lemma~3.25}\LongLabel{eq:gau-HPF}{(3.7)}
\LongLabel{lem:gau-poly}{Lemma~3.26}\LongLabel{lem:gau-wkb}{Lemma~3.27}
\LongLabel{lem:gau-colmod}{Lemma~3.28}\LongLabel{lem:gau-colrows}{Lemma~3.29}
\LongLabel{thm:gau-R0}{Theorem~3.30}\LongLabel{eq:gau-R0zero}{(3.8)}
\LongLabel{cor:gau-rowsNN}{Corollary~3.31}\LongLabel{cor:gau-colours}{Corollary~3.32}
\LongLabel{eq:gau-R0}{(3.9)}\LongLabel{def:gau-Fg}{Definition~3.33}
\LongLabel{sec:tra}{Section~4}\LongLabel{sec:tra-tower}{Section~4.1}
\LongLabel{conv:tra-forms}{Convention~4.1}\LongLabel{lem:tra-sum}{Lemma~4.2}
\LongLabel{def:tra-tower}{Definition~4.3}\LongLabel{eq:tra-tower}{(4.1)}
\LongLabel{eq:tra-shift}{(4.2)}\LongLabel{eq:tra-values}{(4.3)}
\LongLabel{eq:tra-pairing}{(4.4)}\LongLabel{lem:tra-theta}{Lemma~4.4}
\LongLabel{def:tra-slice}{Definition~4.5}\LongLabel{eq:tra-slice}{(4.5)}
\LongLabel{lem:tra-frames}{Lemma~4.6}\LongLabel{eq:tra-decomp}{(4.6)}
\LongLabel{cor:tra-slice}{Corollary~4.7}\LongLabel{lem:tra-coeff}{Lemma~4.8}
\LongLabel{lem:tra-twoslot}{Lemma~4.9}\LongLabel{sec:tra-dress}{Section~4.2}
\LongLabel{def:tra-dressing}{Definition~4.10}\LongLabel{def:tra-edge}{Definition~4.11}
\LongLabel{eq:tra-leaf}{(4.7)}\LongLabel{sec:tra-lemmas}{Section~4.3}
\LongLabel{lem:tra-E}{Lemma~4.12}\LongLabel{eq:tra-lemE}{(4.8)}
\LongLabel{cor:tra-edges}{Corollary~4.13}\LongLabel{lem:tra-T}{Lemma~4.14}
\LongLabel{lem:tra-F}{Lemma~4.15}\LongLabel{prop:tra-can}{Proposition~4.16}
\LongLabel{sec:red}{Section~5}\LongLabel{sec:red-sum}{Section~5.1}
\LongLabel{def:red-flatsum}{Definition~5.1}\LongLabel{eq:red-vertex}{(5.1)}
\LongLabel{def:red-ML}{Definition~5.2}\LongLabel{lem:red-genus}{Lemma~5.3}
\LongLabel{sec:red-thm}{Section~5.2}\LongLabel{thm:red-main}{Theorem~5.4}
\LongLabel{cor:red-can}{Corollary~5.5}\LongLabel{sec:mod}{Section~6}
\LongLabel{sec:mod-tr}{Section~6.1}\LongLabel{conv:mod-kernel}{Convention~6.1}
\LongLabel{eq:mod-kernel}{(6.1)}\LongLabel{def:mod-curve}{Definition~6.2}
\LongLabel{def:mod-pair}{Definition~6.3}\LongLabel{def:mod-xi}{Definition~6.4}
\LongLabel{def:mod-ancvec}{Definition~6.5}\LongLabel{thm:mod-tr}{Theorem~6.6}
\LongLabel{eq:mod-tr}{(6.2)}\LongLabel{lem:mod-laplace}{Lemma~6.7}
\LongLabel{eq:mod-laplace}{(6.3)}\LongLabel{lem:mod-scale}{Lemma~6.8}
\LongLabel{prop:mod-Rtr}{Proposition~6.9}\LongLabel{eq:mod-Rtr}{(6.4)}
\LongLabel{eq:mod-edge}{(6.5)}\LongLabel{eq:mod-Rtr-pf}{(6.6)}
\LongLabel{eq:mod-235}{(6.7)}\LongLabel{lem:mod-xfun}{Lemma~6.10}
\LongLabel{sec:mod-legs}{Section~6.2}\LongLabel{def:mod-legs}{Definition~6.11}
\LongLabel{prop:mod-LD}{Proposition~6.12}\LongLabel{eq:mod-LD}{(6.8)}
\LongLabel{cor:mod-trF}{Corollary~6.13}\LongLabel{def:mod-weights}{Definition~6.14}
\LongLabel{lem:mod-closure}{Lemma~6.15}\LongLabel{lem:mod-genfn}{Lemma~6.16}
\LongLabel{lem:mod-resid}{Lemma~6.17}\LongLabel{prop:mod-dii}{Proposition~6.18}
\LongLabel{eq:mod-dii}{(6.9)}\LongLabel{sec:mod-lattice}{Section~6.3}
\LongLabel{lem:mod-count}{Lemma~6.19}\LongLabel{lem:mod-classL}{Lemma~6.20}
\LongLabel{lem:mod-modv}{Lemma~6.21}\LongLabel{thm:mod-lattice}{Theorem~6.22}
\LongLabel{cor:mod-split}{Corollary~6.23}\LongLabel{cor:mod-Ksplit}{Corollary~6.24}
\LongLabel{sec:mod-Lreg}{Section~6.4}\LongLabel{eq:mod-abdks}{(6.10)}
\LongLabel{eq:mod-710}{(6.11)}\LongLabel{lem:mod-abdks}{Lemma~6.25}
\LongLabel{thm:mod-Lreg}{Theorem~6.26}\LongLabel{prop:mod-period}{Proposition~6.27}
\LongLabel{thm:mod-gap}{Theorem~6.28}\LongLabel{rmk:mod-hock}{Remark~6.29}
\LongLabel{sec:mod-V}{Section~6.5}\LongLabel{lem:mod-fixedx}{Lemma~6.30}
\LongLabel{lem:mod-ibp}{Lemma~6.31}\LongLabel{thm:mod-var}{Theorem~6.32}
\LongLabel{lem:mod-genus0}{Lemma~6.33}\LongLabel{prop:mod-V1}{Proposition~6.34}
\LongLabel{eq:mod-V1}{(6.12)}\LongLabel{eq:mod-loc}{(6.13)}
\LongLabel{lem:mod-w11}{Lemma~6.35}\LongLabel{thm:mod-V}{Theorem~6.36}
\LongLabel{sec:mod-MLK}{Section~6.6}\LongLabel{thm:mod-MLK}{Theorem~6.37}
\LongLabel{cor:mod-onepoint}{Corollary~6.38}\LongLabel{rmk:mod-offK}{Remark~6.39}
\LongLabel{sec:flh}{Section~7}\LongLabel{sec:flh-lemmas}{Section~7.1}
\LongLabel{def:flh-outer}{Definition~7.1}\LongLabel{lem:flh-frame}{Lemma~7.2}
\LongLabel{lem:flh-dil}{Lemma~7.3}\LongLabel{lem:flh-comm}{Lemma~7.4}
\LongLabel{lem:flh-hat}{Lemma~7.5}\LongLabel{lem:flh-inputs}{Lemma~7.6}
\LongLabel{sec:flh-thm}{Section~7.2}\LongLabel{thm:flh-main}{Theorem~7.7}
\LongLabel{eq:flh-main}{(7.1)}\LongLabel{rmk:flh-fail}{Remark~7.8}
\LongLabel{sec:con}{Section~8}\LongLabel{sec:con-problem}{Section~8.1}
\LongLabel{eq:con-chain}{(8.1)}\LongLabel{rmk:con-genus}{Remark~8.1}
\LongLabel{sec:con-form}{Section~8.2}\LongLabel{def:gam-c}{Definition~8.2}
\LongLabel{prop:gam-rational}{Proposition~8.3}\LongLabel{lem:con-form}{Lemma~8.4}
\LongLabel{def:con-PhiSP}{Definition~8.5}\LongLabel{sec:con-P}{Section~8.3}
\LongLabel{thm:con-Pa}{Theorem~8.6}\LongLabel{eq:con-shiftD}{(8.2)}
\LongLabel{eq:con-bz}{(8.3)}\LongLabel{eq:con-factor}{(8.4)}
\LongLabel{lem:con-cfu}{Lemma~8.7}\LongLabel{lem:con-red}{Lemma~8.8}
\LongLabel{lem:con-S2}{Lemma~8.9}\LongLabel{lem:con-id}{Lemma~8.10}
\LongLabel{lem:con-relab}{Lemma~8.11}\LongLabel{lem:con-key}{Lemma~8.12}
\LongLabel{prop:con-res}{Proposition~8.13}\LongLabel{thm:con-Pb}{Theorem~8.14}
\LongLabel{sec:con-glue}{Section~8.4}\LongLabel{thm:con-glue}{Theorem~8.15}
\LongLabel{eq:con-glue}{(8.5)}\LongLabel{cor:con-A2}{Corollary~8.16}
\LongLabel{sec:gam}{Section~9}\LongLabel{sec:gam-series}{Section~9.1}
\LongLabel{sec:gam-lmi}{Section~9.2}\LongLabel{lem:gam-deltasrc}{Lemma~9.1}
\LongLabel{eq:gam-deltasrc}{(9.1)}\LongLabel{def:gam-lmidata}{Definition~9.2}
\LongLabel{thm:gam-lmi}{Theorem~9.3}\LongLabel{eq:gam-lmi}{(9.2)}
\LongLabel{eq:gam-step1}{(9.3)}\LongLabel{eq:gam-chart}{(9.4)}
\LongLabel{eq:gam-atQ0}{(9.5)}\LongLabel{sec:gam-thm}{Section~9.3}
\LongLabel{thm:gam-main}{Theorem~9.4}\LongLabel{prop:gam-hatflat}{Proposition~9.5}
\LongLabel{cor:gam-equiv}{Corollary~9.6}\LongLabel{thm:gam-D}{Theorem~9.7}
\LongLabel{thm:gam-ginf}{Theorem~9.8}\LongLabel{eq:gam-ginf}{(9.6)}
\LongLabel{cor:gam-poly}{Corollary~9.9}\LongLabel{cor:gam-free}{Corollary~9.10}
\LongLabel{sec:gam-rmk}{Section~9.4}\LongLabel{rmk:gam-prior}{Remark~9.11}
\LongLabel{sec:mai}{Section~10}\LongLabel{sec:mai-thm}{Section~10.1}
\LongLabel{thm:mai-main}{Theorem~10.1}\LongLabel{lem:mai-Psi}{Lemma~10.2}
\LongLabel{def:mai-Fcf}{Definition~10.3}\LongLabel{cor:mai-gap}{Corollary~10.4}
\LongLabel{rmk:mai-constant}{Remark~10.5}\LongLabel{rmk:mai-graphlocal}{Remark~10.6}
\LongLabel{rmk:mai-necessity}{Remark~10.7}\LongLabel{sec:cmp}{Section~11}
\LongLabel{sec:mai-nat}{Section~11.1}\LongLabel{lem:mai-Ghat}{Lemma~11.1}
\LongLabel{cor:mai-natfact}{Corollary~11.2}\LongLabel{cor:mai-actual}{Corollary~11.3}
\LongLabel{sec:mai-univ}{Section~11.2}\LongLabel{lem:mai-pole}{Lemma~11.4}
\LongLabel{cor:mai-univ}{Corollary~11.5}\LongLabel{rmk:mai-annih}{Remark~11.6}
\LongLabel{rmk:mai-stored}{Remark~11.7}\LongLabel{sec:set-lrproof}{Appendix~A}
\LongLabel{lem:set-grading}{Lemma~A.1}\LongLabel{lem:set-unitfrob}{Lemma~A.2}
\LongLabel{lem:set-periods}{Lemma~A.3}\LongLabel{lem:set-birkhoff}{Lemma~A.4}
\LongLabel{lem:set-wronskian}{Lemma~A.5}\LongLabel{lem:set-order3}{Lemma~A.6}
\LongLabel{eq:set-McycLR}{(A.1)}\LongLabel{sec:gau-cycproof}{Appendix~B}
\LongLabel{lem:gau-const}{Lemma~B.1}\LongLabel{lem:gau-adj}{Lemma~B.2}
\LongLabel{lem:gau-conc}{Lemma~B.3}\LongLabel{lem:gau-transv}{Lemma~B.4}
\LongLabel{lem:gau-iso}{Lemma~B.5}\LongLabel{lem:gau-sg}{Lemma~B.6}
\LongLabel{sec:tra-convproof}{Appendix~C}\LongLabel{lem:tra-conv}{Lemma~C.1}
\LongLabel{sec:fsp}{Appendix~D}\LongLabel{sec:fsp-def}{Appendix~D.1}
\LongLabel{def:fsp-ring}{Definition~D.1}\LongLabel{lem:fsp-ring}{Lemma~D.2}
\LongLabel{def:fsp-wick}{Definition~D.3}\LongLabel{lem:fsp-univ}{Lemma~D.4}
\LongLabel{eq:fsp-univ}{(D.1)}\LongLabel{lem:fsp-const}{Lemma~D.5}
\LongLabel{lem:fsp-homog}{Lemma~D.6}\LongLabel{sec:fsp-basic}{Appendix~D.2}
\LongLabel{lem:fsp-ibp}{Lemma~D.7}\LongLabel{eq:fsp-gibp}{(D.2)}
\LongLabel{lem:fsp-pg}{Lemma~D.8}\LongLabel{lem:fsp-prod}{Lemma~D.9}
\LongLabel{lem:fsp-ift}{Lemma~D.10}\LongLabel{lem:fsp-transl}{Lemma~D.11}
\LongLabel{cor:fsp-movebase}{Corollary~D.12}\LongLabel{lem:fsp-param}{Lemma~D.13}
\LongLabel{lem:fsp-dq}{Lemma~D.14}\LongLabel{lem:fsp-generic}{Lemma~D.15}
\LongLabel{eq:fsp-generic}{(D.3)}\LongLabel{sec:fsp-change}{Appendix~D.3}
\LongLabel{lem:fsp-piola}{Lemma~D.16}\LongLabel{lem:fsp-cov}{Lemma~D.17}
\LongLabel{lem:fsp-cfu}{Lemma~D.18}\LongLabel{lem:fsp-fubini}{Lemma~D.19}
\LongLabel{sec:fsp-stirling}{Appendix~D.4}\LongLabel{def:fsp-stirling}{Definition~D.20}
\LongLabel{lem:fsp-stirfe}{Lemma~D.21}\LongLabel{lem:fsp-stirling}{Lemma~D.22}
\LongLabel{eq:fsp-stirling}{(D.4)}\LongLabel{lem:fsp-dlmf5118}{Lemma~D.23}
\LongLabel{eq:fsp-5118}{(D.5)}\LongLabel{lem:fsp-refl}{Lemma~D.24}
\LongLabel{sec:fsp-res}{Appendix~D.5}\LongLabel{lem:fsp-res}{Lemma~D.25}
\LongLabel{sec:sgn}{Appendix~E}\LongLabel{sec:sgn-R}{Appendix~E.1}
\LongLabel{sec:sgn-var}{Appendix~E.2}\LongLabel{sec:sgn-model}{Appendix~E.3}
\LongLabel{sec:sgn-tr}{Appendix~E.4}\LongLabel{tab:sgn-model}{Table~1}
\LongLabel{sec:sgn-legs}{Appendix~E.5}\LongLabel{sec:sgn-tower}{Appendix~E.6}
\LongLabel{tab:sgn-tr}{Table~2}\LongLabel{sec:sgn-norm}{Appendix~E.7}
\LongLabel{sec:sgn-sp}{Appendix~E.8}\LongLabel{sec:sgn-lit}{Appendix~E.9}
\LongLabel{tab:sgn-lit}{Table~3}\LongLabel{sec:evi}{Appendix~F}
\LongLabel{sec:evi-graphs}{Appendix~F.1}\LongLabel{tab:evi-g2}{Table~4}
\LongLabel{tab:evi-g3}{Table~5}\LongLabel{tab:evi-g4}{Table~6}
\LongLabel{tab:evi-g5}{Table~7}\LongLabel{sec:evi-normal}{Appendix~F.2}
\LongLabel{tab:evi-windows}{Table~8}\LongLabel{sec:evi-model}{Appendix~F.3}
\LongLabel{tab:evi-model}{Table~9}


% Proof sketches
\newenvironment{sketch}[1][Sketch of the proof]{\begin{proof}[#1]}{\end{proof}}

%-------------------------------------------------------------------------------
% Title

\title[The formal quintic at the conifold point (abridged version)]{Conifold regularity of the formal quintic, graph by graph\\ (abridged version)}
\author{The GENUS collaboration}
\genusproject{solutions/conifold-gap-formal-quintic (abridged)}

\emergencystretch=2em
\begin{document}

\begin{abstract}
The formal quintic of Lho and Pandharipande is the $\cO(5)$-twisted equivariant Gromov--Witten theory of $\bbP^4$ specialised at the fifth roots of unity. Its genus-$g$ potential is a Givental--Teleman sum over stable graphs of genus $g$, with the $R$-matrix normalised at $q=0$ by quantum Riemann--Roch. We expand the theory at the conifold point $u=1-5^5q=0$ in a conifold frame: the graph sum is written in a differential ring generated by the periods, and a regular period and the vanishing period at $u=0$ are substituted for the two periods at $q=0$. This is a formal change of frame; no analytic continuation is involved. We factorise the expanded $R$-matrix as $R=R'\,\Mo(\tau(u)\zh)$, where $\Mo$ is the $R$-matrix of the Landau--Ginzburg model $p^4/4+v/p$ with primitive form $\dd p$ and $\tau(u)\zh$ is the natural variable defined by the canonical coordinate. For every $g\geq 2$ and every conifold frame of a two-parameter rational family, in the graph sum of $R'$ on the theory obtained from the topological part by the action of $\Mo(\tau(u)\zh)$, the term of every leg-free stable graph of genus $g$ with at least one edge is regular at the conifold point, and the single-vertex graph contributes $c_g\tc^{2-2g}$ plus a regular series, where $\tc$ is the flat coordinate of the vanishing period and $c_g=(-1)^{g-1}B_{2g}/\bigl(2g(2g-2)\,5^{g-1}\bigr)$. Summing over graphs gives the conifold gap for this expansion of the genus-$g$ potential of the formal quintic, for every $g\geq2$. The statement is reduced, graph by graph, to a lemma about the model, proved by topological recursion on the model curve, and to the equality of two normalisations of the $R$-matrix, the canonical one at the conifold point and the quantum Riemann--Roch one at $q=0$, which follows from an identity of formal stationary phase for a Mellin--Barnes integral. The results hold modulo inputs from the literature: results on topological recursion, applied only to the model curve; results of Pandharipande, Pixton and Zvonkine on the action of $R$-matrices, including the formula for the pull-back of $\psi$ classes along the map forgetting a point; the fact that the $R$-matrix is the asymptotic solution of the quantum differential equation, and its value at $q=0$ given by quantum Riemann--Roch; the genus-$0$ mirror theorem and the genus-$0$ axioms of the twisted theory; Teleman's reconstruction, with the axioms of the twisted theory in all genera; and a push-forward formula of Arbarello and Cornalba for products of $\psi$ classes; every other step is proved in the long version. We do not treat the quintic threefold itself. This is the abridged version of the paper: it gives the statements and the key points of the argument, and the complete proofs are in the long version.
\end{abstract}

\maketitle
\genusdisclaimer

% Abridged version, Section 1, part 1 of 2: the setting, the conifold gap, the results. Continued in part 2.
\section{Introduction}
\label{sec:int}

This paper is the abridged version of a longer one with the same title. Throughout, [L] denotes \emph{the long version of this paper}, and a reference such as \Lr{thm:gau-U} is to the numbering of that version. We give the definitions needed to state the results, the precise statements, and the key points of the argument. Proofs are sketches unless they are called proofs, and the complete proofs are in [L]. Every statement is a statement of [L]; where a statement is shortened by omitting auxiliary parts, it is marked as abridged.

%-------------------------------------------------------------------------------
\subsection{The setting}
\label{sec:int-setting}

Let the torus $(\bbC^*)^5$ act on $\bbP^4$ with weights $\lambda_0,\dots,\lambda_4$ on $H^0(\bbP^4,\cO(1))$. The \emph{formal quintic} of Lho and Pandharipande \cite[\S0.2, (5)]{LP18}, \cite[\S0.2, (5)]{LP20} is the $\cO(5)$-twisted equivariant Gromov--Witten theory of $\bbP^4$ at the specialisation $\lambda_i=\zeta^i$, $\zeta=\re^{2\pi\ri/5}$: its invariants are $N_{g,d}=\int_{[\overline{\mathcal M}_g(\bbP^4,d)]^{\mathrm{vir}}}e(R\pi_*f^*\cO(5))$, defined by virtual localisation \cite{MR1666787} before the weights are specialised, and $F_g=\sum_{d\geq0}N_{g,d}\Qn^d$ (\cref{def:set-fq}). By the genus-zero mirror theorem the $J$-function at $t=T(q)H$ is $I(q,z)/C_0(q)$, with $C_0=\sum_d\frac{(5d)!}{(d!)^5}q^d$ and the mirror map $T=I_1/C_0$, and $\Qn=\re^{T(q)}$. With $u:=1-5^5q$ and $L:=u^{-1/5}=1/s$, the quantum product is cyclic (\cref{prop:set-cyclicLR}); $q=0$ is $u=1$, and $u=0$ is the \emph{conifold point}.

By Teleman's reconstruction \cite[Theorem~(2) in (1.15)]{MR2917177} (see \cite[\S4.1, Theorem~4.1]{GJR18}) $F_g$ is the Givental--Teleman sum over the stable graphs of genus $g$ of Givental's symplectic $R$-matrix $R(z)=1+O(z)$ acting on the topological field theory $\TFT$, in the convention of \cite[Defs.~2.2, 2.13]{PPZ15} (\cref{def:set-rows,def:set-graphsum,thm:set-GT}); at $q=0$ quantum Riemann--Roch \cite[(17)]{GJR18} gives
\[
R(q=0)=\diag\bigl(h_{\bw_0},\dots,h_{\bw_4}\bigr),\qquad h_{\bw}(z)=\exp\Bigl(-\sum_{k\geq1}\bk_k\,z^{2k-1}\sum_{a\in\bw}a^{1-2k}\Bigr),\qquad \bk_k=\frac{B_{2k}}{2k(2k-1)},
\]
with $\bw_i=\{\zeta^i-\zeta^j\}_{j\neq i}\cup\{-5\zeta^i\}$, the input (QRR) (\cref{thm:set-qrr}); $h:=h_{\bw_0}$. In the conifold expansion defined below, where $L=s^{-1}$, the lowered rows $\Qr^{(i)}_a(z)=\eta(H^a,R(z)\tilde e_i)$ grow with $L$, and the natural variable is $\zh:=L^4z=z/s^4$.

\emph{Conifold frames.} Let $f_0,f_1,f_2$ be the solutions $f_k=u^k+O(u^3)$ at $u=0$ of the Picard--Fuchs operator $\PFfour$ of the quintic, and $\varpi_1=f_1+\frac7{10}f_2=u+O(u^2)$ the vanishing period (\cref{lem:set-frobbasis}). A \emph{conifold frame} is a period $\varpi_0=f_0+\alpha\varpi_1+\beta f_2$, $(\alpha,\beta)\in\bbQ^2$, and $\tc:=\varpi_1/\varpi_0=u+O(u^2)$ (\cref{def:set-frames}). The row coefficients lie in $\Nring\tensor\bbQ(\zeta)$, $\Nring$ a differential ring generated by the periods, $L^{\pm1}$ and a formal $\log q$ (\cref{def:set-confexp,cor:gau-rowsNN}); the \emph{conifold expansion} $\Cexp_{\alpha,\beta}\colon\Nring\to\bbQ\lser{s}$ expands at $u=0$ in the frame, and an expanded quantity is \emph{regular at the conifold point} if it lies in $\bbQ\fps{u}$.

\emph{The model and the natural form.} The model is the Landau--Ginzburg potential $\Wv(p)=p^4/4+v/p$ with primitive form $\dd p$, whose five critical points $p^5=v$ coalesce into the pole $p=0$ as $v\to0$ (\cref{def:set-model}). Its rows, the normalised stationary-phase series of $\int p^{\nu_b}\re^{-\Wv(p)/z}\dd p$, $\nu=(0,-1,3,2,1)$, form the model $R$-matrix $\Mo(t)$, $t=z/\varsigma^4$, $\varsigma=v^{1/5}$ (\cref{prop:set-modrows,def:set-M0}). Let $\tau(u):=1/{}_2F_1(1,\frac45;\frac95;u)=\bigl(\sum_{n\geq0}\frac{4u^n}{5n+4}\bigr)^{-1}$ and $v(u):=u\,\tau(u)^{-5/4}$ (\cref{def:set-tau}). At $v=v(u)$ the canonical coordinates of the model are those of the formal quintic measured from the conifold point, and $z/\varsigma^4=\tau(u)\zh=-\frac54\zeta^iz/(u_i-u_i^{\mathrm c})$ (\cref{lem:set-canon,prop:gau-equalprod}); so $M:=\Mo(\tau(u)\zh)$ is the $R$-matrix of the model at $v(u)$ in the variable of the formal quintic. Writing the expanded rows as $\Qr=R'M$, $R':=\Qr M^{-1}$, gives the \emph{natural form} of the graph sum: the outer graph sum $R'.(M.\TFT)$ has the vertices $M.\TFT$ and the edges and leaves of $R'$ (\cref{def:set-outer}).

%-------------------------------------------------------------------------------
\subsection{The conifold gap}
\label{sec:int-gap}

The conifold gap is a prediction of the topological B-model: near a conifold point the genus-$g$ potential, expanded in the flat coordinate of the vanishing period, has the universal leading pole of the $c=1$ string at the self-dual radius and no other pole \cite{Bershadsky:1993cx} (see \cite[\S1]{HKQ09}), \cite{Ghoshal:1995wm}, \cite[\S2.3]{HK07}, \cite[(1.2)]{HKQ09}. For the quintic threefold it was left open after \cite{GJR18,CGL21,CGL26} (see \cite[Remark~1.6]{GJR18}), was an input of high-genus computations \cite{HKQ09,LR22}, and is proved in \cite[Theorem~2]{CGYZ26} by master-space localisation. It holds for the resolved conifold \cite[\S1.1, footnote~1]{Br25}, for local $\bbP^2$ in all genera \cite{Br25}, and for the topological recursion of toric mirror curves at a generic one-node degeneration \cite{FCCZ26}.

The potentials of the formal quintic satisfy holomorphic anomaly equations \cite[Theorems~1 and~2, \S0.4]{LP20}. By \cite[Remark~1.9]{CGYZ26}, the approach of \cite{GJR18} reduces the gap of the quintic to two key statements, the conifold gap for the formal quintic and the conifold regularity of its correlators with special insertions, and the authors of \cite{CGYZ26} expect their method to prove the first as well. We prove, conditionally on the inputs of \cref{sec:int-inputs}, a graph-local statement for the formal quintic (Theorem~A); summing over the graphs gives its conifold gap in every genus $g\geq2$ (Corollary~B), with the constant of \cite[Conjecture~1.5]{CGYZ26} (\cref{rmk:int-constant}). The proof of \cite{CGYZ26} also factorises an $R$-matrix, that of a master-space theory \cite[Proposition~3.1]{CGYZ26}; here the $R$-matrix of the formal quintic itself is factorised, through a universal local model.

The normalisation at $q=0$ is carried to the conifold point by an identity of formal stationary phase (Theorem~Gamma), through a change of chart as in Givental's treatment of projective spaces \cite[\S10]{MR1901075}. The analogous identity for the Landau--Ginzburg integral of Guo, Janda and Ruan is stated without proof in \cite[Remark~7.2]{GJR18}.

%-------------------------------------------------------------------------------
\subsection{Results}
\label{sec:int-results}

For a leg-free stable graph $\Gamma$, $\mathrm{Cont}_\Gamma[E_X,q_X;\Omega]$ denotes the contraction of the vertex tensors of a cohomological field theory $\Omega$ along the edges $E_X$ of a matrix series $X$, with its leaf point $q_X$, including the factor $|\Aut\Gamma|^{-1}$ (\cref{def:set-graphsum,def:set-outer}).
The \emph{vertex graph} has one vertex, of genus $g$, and no edge; its term contains the leaf decorations.

\begin{thm}[Theorem~A; {\Lr{thm:int-main,thm:mai-main}}]\genusstatus{AI-proposed, unchecked}
\label{thm:int-main}
Let $g\geq2$ and let $(\alpha,\beta)\in\bbQ^2$ define a conifold frame (\cref{def:set-frames}), with vanishing period $\varpi_1=u+O(u^2)$ and $\tc=\varpi_1/\varpi_0$. Let $M:=\Mo(\tau(u)\zh)$, exactly in $u$, and $R':=\Qr M^{-1}$, where $\Qr$ is the matrix of the lowered rows of \cref{def:set-rows} expanded in the frame (\cref{def:set-confexp}). For a leg-free stable graph $\Gamma$ of genus $g$ let
\[
T_\Gamma:=\mathrm{Cont}_\Gamma\bigl[E_{R'},q_{R'};M.\TFT\bigr]
\]
be its outer term (\cref{def:set-outer}), where $\TFT$ is the topological field theory of the formal quintic (\cref{def:set-frob}), expanded in the frame. Then:
\begin{enumerate}
\item $T_\Gamma\in\bbQ\fps{u}$ for every $\Gamma$ with at least one edge;
\item for the vertex graph, $T_\Gamma\in c_g\tc^{2-2g}+\bbQ\fps{u}$, where
\[
c_g=\frac{(-1)^{g-1}B_{2g}}{2g(2g-2)\,5^{g-1}}=\frac{|B_{2g}|}{2g(2g-2)\,5^{g-1}} .
\]
\end{enumerate}
\end{thm}

For one $g$, \cref{thm:int-main} is Property N($g$) of \cref{def:set-N}. The hatted terms are $G_\Gamma=u^{2g-2}T_\Gamma$ (\cref{lem:flh-hat}). The statement is about each graph separately; we do not know how to obtain it from a statement about the singular part of the sum \Lr{rmk:mai-graphlocal}.

\begin{cor}[Corollary~B: the conifold gap of the formal quintic; {\Lr{cor:int-gap,cor:mai-gap}}]\genusstatus{AI-proposed, unchecked}
\label{cor:int-gap}
Let $g\geq2$ and $(\alpha,\beta)\in\bbQ^2$. The conifold expansion $F_g^{\mathrm{cf}}\in\bbQ\lser{s}$ of the genus-$g$ potential of the formal quintic in the frame $(\alpha,\beta)$ (\cref{def:mai-Fcf}) satisfies
\[
F_g^{\mathrm{cf}}\in c_g\tc^{2-2g}+\bbQ\fps{u} .
\]
Since $\bbQ\fps{u}=\bbQ\fps{\tc}$, the series $F_g^{\mathrm{cf}}$ has no fractional power of $u$, no term $\tc^k$ with $3-2g\leq k\leq-1$, and leading term $c_g\tc^{2-2g}$.
\end{cor}

Here $F_g^{\mathrm{cf}}$ is the image under $\Cexp_{\alpha,\beta}$ of the genus-$g$ graph sum written in $\Nring$, whose evaluation is $F_g$ by (T) and (QRR); by the group property of the action \cite[Def.~2.13, Prop.~2.14]{PPZ15} it is the sum of the $T_\Gamma$ (\cref{cor:mai-gap}).

The proof compares two normalisations of the $R$-matrix. The cyclic systems of the formal quintic and of the model at $v(u)$ are related by a gauge regular at $u=0$, unique up to a scalar factor $\cn(z)$ (Theorem~U, \cref{thm:gau-U}); the \emph{canonical gauge} $\Pc$ is fixed at the conifold point by the gauge identity $\Pc\YC(v(u))=\YCQ$ between the vanishing solutions (\cref{thm:gau-G,def:gau-can}), the rows $\Qr$ at $q=0$ by (QRR). Both $\Qr$ and $\Pc\Qr^{\mathrm m}(v(u))$, $\Qr^{\mathrm m}$ the model rows, are formal WKB solutions of the same system with the same exponentials, so
\[
\Qr\,\Dn=\Pc\,\Qr^{\mathrm m}(v(u)),\qquad \Dn=\diag(\Dn_0,\dots,\Dn_4),\qquad \Dn_i(z)=\Dn_0(\zeta^{-i}z),\qquad \Dn_0\in\bbQ\fps{z}
\]
(\cref{thm:gau-R0}; that $\Dn_0$ has rational coefficients is part of \cref{thm:con-glue}).

\begin{thm}[Theorem~C: the normalisation; {\Lr{thm:int-D,thm:gam-D,cor:mai-actual}}]\genusstatus{AI-proposed, unchecked}
\label{thm:int-D}
In every conifold frame, $\Dn_0=1$ in $\bbQ\fps{z}$. Hence $\Dn=1$, and the window $\Fg{g}$, that is $\Dn_0\equiv1\bmod z^{3g-2}$ (\cref{def:gau-Fg}), holds for every $g\geq2$. Equivalently, $\Qr=\Pc\,\Qr^{\mathrm m}(v(u))$: the actual gauge is the canonical gauge, and the $R$-matrix of the formal quintic, normalised at $q=0$ by (QRR), is in the canonical normalisation at the conifold point.
\end{thm}

\Cref{thm:int-D} is reduced to one value at $q=0$ of a saddle series:

\begin{thm}[Theorem~Gamma; {\Lr{thm:int-gamma,thm:gam-main,thm:gam-ginf}}]\genusstatus{AI-proposed, unchecked}
\label{thm:int-gamma}
Let $c(z;L)=\sum_{n\geq0}c_n(L)z^n$, $c_0(L)=1$, be the full-Stirling saddle series of \cref{def:gam-c}: the normalised stationary-phase series, divided by its constant term ($\SPh$ of \cref{sec:con-fsp}), at $\sigma=L$ of
\[
\Gamma\Bigl(1+\frac{5\sigma}{z}\Bigr)\prod_{j=0}^{4}\Gamma\Bigl(\frac{\zeta^j-\sigma}{z}\Bigr)(-q)^{\sigma/z}\,\dd\sigma,\qquad q=q(L)=\frac{1-L^{-5}}{5^5},
\]
with every Gamma factor replaced by its formal Stirling form (\cref{sec:con-fsp}). Then every $c_n$ lies in $\AL=\bbQ[L,L^{-1},(L^5-1)^{-1}]$, and:
\begin{enumerate}
\item[(a)] no $c_n$ has a pole at $L=1$, that is, at $q=0$, and $c(z;1)=h(z)$;
\item[(b)] no $c_n$ has a pole at $L=0$, that is, at $q=\infty$, and
\[
c(z;0)=\Ginf(z):=\exp\Bigl(\sum_{l\geq1}\frac{(-1)^{l+1}B_{5l+1}(\frac15)}{l(5l+1)}\,z^{5l}\Bigr),
\]
the series given by \cite[Lemma~7.7, (37)]{GJR18} at $m=5$, $k=1$, $\lambda=1$ in the notation there.
\end{enumerate}
\end{thm}

Part (a) is proved by a local Mellin inversion (\cref{thm:gam-lmi}) from the lemmas of formal stationary phase of \Lr{sec:fsp} alone; by $c_n(L)=L^{4n}\mathrm{ph}^{\mathrm{SP}}_n(L^{-5})$ (\cref{thm:con-Pa}), Theorem~P (\cref{thm:con-Pb}) and the glue (\cref{thm:con-glue}) it is equivalent to $\Dn=1$ (\cref{cor:con-A2}). Part (b) identifies the free coefficients of the canonical colour-$0$ function \Lr{cor:gam-free} and is not used for \cref{thm:int-main}; the series on the right is obtained in \cite{GJR18} for the Landau--Ginzburg series of that paper, and we do not claim that the two series coincide.

\begin{rmk}[{\Lr{rmk:int-constant,rmk:mai-constant}}]
\label{rmk:int-constant}
The constant $c_g$ is that of \cite[Conjecture~1.5]{CGYZ26} for the quintic threefold, in the same normalisation $\varpi_1=u+O(u^2)$; the two theories have different potentials, so nothing about the quintic threefold follows. For a frame whose pair $(\varpi_0,\varpi_1)$ is admissible in the sense of \cite[Def.~1.2]{CGYZ26}, which is asserted in \cite[Ex.~1.3, Rem.~1.4]{CGYZ26} and is not proved here, $F_g^{\mathrm{cf}}$ is, under the inputs of \cref{thm:int-D}, the conifold expansion of $F_g$ in the sense of that definition \Lr{sec:sgn-lit}.
\end{rmk}

% Abridged version, Section 1, part 2 of 2: inputs, what is not claimed, strategy and plan, acknowledgements.

%-------------------------------------------------------------------------------
\subsection{Inputs}
\label{sec:int-inputs}

The results are conditional on the following inputs, the same as in \Lr{sec:int-inputs}. Statement numbers refer to the arXiv versions named in the bibliography.
\begin{description}
\item[(Q)] \emph{Quoted theorems.} Topological recursion, applied only to the model curve $x=\Wv(p)$, $y=p$ and its rescalings $x=\kappa\Wv$ (\cref{sec:mod}): \cite[Def.~4.2, (4-9), with the kernel (4-3); Theorems~4.2 and~4.6; Theorem~5.1, (5-14), with \S5.3, (5-27)--(5-30)]{Eynard:2007kz}; the theorem \cite[Theorem~4.1, with (4-1)--(4-3)]{DuninBarkowski:2012bw}, with its graph weights \cite[\S2.3, (2-22)--(2-26)]{DuninBarkowski:2012bw}, as recalled in \cite[\S2.6, (2-30)--(2-32), (2-34), (2-35)]{Dunin-Barkowski:2015caa} and read as the action of \cite[Defs.~2.2, 2.13]{PPZ15} (the conventions of this reading, and a misprint in \cite[(2-33)]{Dunin-Barkowski:2015caa}, which is not used, are specified in \Lr{sec:int-inputs}); its hypotheses, namely symplecticity, \cite[(2-32)]{Dunin-Barkowski:2015caa} and the edge condition \cite[(4-1)]{DuninBarkowski:2012bw}, are proved for the model curve in \Lr{prop:mod-Rtr}; and \cite[Theorem~7.1: (7.6), (7.7) with its factor $1/\hbar^2$, (7.9), (7.10), with the kernel (1.1) and the assumptions of \S1.1]{ABDKS25}. The $R$-matrix action: \cite[\S0.2, (1); Def.~0.3, in particular (iii); (5); Remark~2.1 and (6); Def.~2.2, with the convention stated after it (p.~20); Def.~2.5, with the convention stated after it; Prop.~2.12, with the formula (8) in its proof; Def.~2.13; Prop.~2.14]{PPZ15}; Props.~2.12 and~2.14 are used as statements over $\bbQ$ and transferred to the other coefficient rings in \Lr{def:set-graphsum}.
\item[(QRR)] \emph{The $R$-matrix along the line: flatness and the value at the large-radius point} (\cref{thm:set-qrr}). The $R$-matrix $R$ of (T) has coefficients in $\End(\bbH)\otimes\bbQ(\zeta)\fps q$ and satisfies: (\emph{flatness}) for every colour $i$ the column $R\,\tilde e_i\,\re^{u_i/z}$ solves the quantum differential equation $z\partial_Tv=H*v$ along the line (\cref{def:set-rows}~(b)); (\emph{value}) $R(q=0)=\diag(h_{\bw_0},\dots,h_{\bw_4})$ in the normalised canonical basis at $q=0$. Flatness is a theorem about the $R$-matrix of the cohomological field theory, not a convention; the convention (Givental's) is only that the legs of the graph sum carry $R^{-1}(\psi)$. Sources for flatness: \cite[Lemma~5.2, (16)]{GJR18}; \cite[\S4.2, (12)--(13)]{GJR17}; in general \cite[Section~1, Proposition (a)--(c)]{Giv01} and \cite[Props.~7.7, 8.2~(ii), 8.4, 8.7~(ii)]{MR2917177}. Sources for the value: \cite[(17), in the proof of Lemma~5.2]{GJR18}; \cite[(19)--(20), in the proof of Lemma~4.6, at $t=0$]{GJR17}; \cite[\S2.2, Theorem~1 and the Remark after it]{Giv01}; \cite[Theorem~1, Corollary~1]{MR2276766}.
\item[(M)] \emph{The genus-zero mirror theorem} \cite[Theorem~11.6]{MR1408320}, in the Calabi--Yau case $n=4$, $r=1$, $l_1=5$, $\lambda'=0$, with the set-up of \cite[\S9]{MR1408320} and \cite[Prop.~11.4]{MR1408320}: at $\lambda_i=\zeta^i$, $J(T(q)H,z)=I(q,z)/C_0(q)$ with $T=I_1/C_0$ (\cref{sec:set-fq}).
\item[(G0)] \emph{The genus-zero axioms of the twisted theory}: the genus-zero twisted classes form a cohomological field theory with unit for $\eta$ \cite[\S4.1]{GJR18}, \cite[Def.~0.3]{PPZ15}, polynomial and $S_5$-symmetric in $\lambda$, with the string, divisor and splitting axioms, the genus-zero topological recursion relation and the quantum differential equation \cite[\S5.2, the display before (14), of which (14) is the form in the $q$-direction]{GJR18}.
\item[(T)] \emph{Teleman's reconstruction, with the axioms of the twisted theory in all genera} (\cref{thm:set-GT}). (i) At $\lambda_i=\zeta^i$ the classes $\Omega_{g,n}$ of \cref{def:set-fq}, in all genera, form a cohomological field theory with flat unit $1$ for $\eta$ \cite[\S4.1]{GJR18}, \cite[Def.~0.3]{PPZ15}, and satisfy the divisor equation: for $2g-2+n>0$ the part of degree $d$ of $\pi_*\Omega_{g,n+1}(\gamma_1,\dots,\gamma_n,H)$ is $d$ times that of $\Omega_{g,n}(\gamma_1,\dots,\gamma_n)$, where $\pi$ forgets the last point; so the theory at $\Qn=\re^{T(q)}$ is the theory at the point $t=T(q)H$ with $\Qn=1$. (ii) This theory, which is semisimple (\cref{cor:set-cyclicPF5}~(a)), is the graph sum of \cref{def:set-graphsum}, $\Omega=R.\TFT$, for an $R$-matrix $R$ with coefficients in $\End(\bbH)\otimes\bbQ(\zeta)\fps q$, symplectic and with $R|_{z=0}=1$ \cite[Theorem~(2) in (1.15)]{MR2917177}, \cite[\S4.1, Theorem~4.1]{GJR18}; this is the $R$-matrix of \cref{def:set-rows}, to which (QRR) refers. In \cref{thm:int-main,thm:int-D} (T) identifies that $R$-matrix, whose rows define $T_\Gamma$ and $\Dn$; in \cref{cor:int-gap} it identifies the total with $F_g$.
\end{description}
In addition, the formula \cite[(1.7)]{AC96}, with $\kappa_0=2g-2+n$ \cite[p.~5]{AC96}, is used in the dilaton equation of \cref{lem:set-dilaton}~(ii). The definitions use \cite[\S0.2, (5)]{LP18}, \cite[\S0.2, (5), (10)]{LP20} and \cite{MR1666787}. The inputs are read with the definitions and conventions of their sources that are listed in \Lr{sec:int-inputs}: the normalisation of the $J$-function in (M), the free energies and the set-up of the variational formula of \cite{Eynard:2007kz}, the recursion \cite[(2-26)--(2-29)]{Dunin-Barkowski:2015caa}, the reading of the integrals \cite[(2-31), (2-32)]{Dunin-Barkowski:2015caa} as formal Laplace series, and the edge condition \cite[(4-1)]{DuninBarkowski:2012bw} in the coordinates of \cite[\S2.6]{Dunin-Barkowski:2015caa}.

\emph{Where the inputs enter.} (M) and (G0) enter only in the genus-zero data of \cref{sec:set} (the mirror theorem, \cref{prop:set-cyclicLR,cor:set-cyclicPF5}); (QRR) through the rows: flatness in \cref{prop:set-cyclicLR}~(iii) and \cref{cor:set-cyclicPF5}~(b), the value in the colour covariance of the rows (\cref{lem:gau-colrows}) and in the glue (\cref{thm:con-glue}); (T) in \cref{thm:set-GT} and \cref{cor:mai-gap}; topological recursion only in \cref{sec:mod}. 

\emph{Conditionality.} Theorem~A holds modulo the quoted theorems (Q) of \cite{Eynard:2007kz}, \cite{DuninBarkowski:2012bw} (as recalled in \cite{Dunin-Barkowski:2015caa}), \cite{ABDKS25} and \cite{PPZ15} with the statement numbers above, the flatness of the $R$-matrix and its value at the large-radius point (QRR), the genus-$0$ mirror theorem (M) and the genus-$0$ axioms (G0) of the twisted theory, Teleman's reconstruction with the axioms of the twisted theory in all genera (T), and the formula \cite[(1.7)]{AC96}; every other step is proved in [L]. \Cref{cor:int-gap} uses, beyond \cref{thm:int-main}, only (T), part of \cite{PPZ15} and proved statements. \Cref{thm:int-D} uses the $R$-matrix that (T) identifies only through $R|_{z=0}=1$, its coefficient ring $\bbQ(\zeta)\fps q$ and (QRR), that is, flatness and the value at $q=0$, and the genus-zero data of \cref{sec:set}, which rest on (M) and (G0). \Cref{thm:int-gamma} uses no quoted input; the identity \cite[24.4.1]{DLMF} of Bernoulli polynomials used in part (b) is derived in [L] from their generating function. No finite computation enters a proof; the computations that accompanied the work are summarised in \cref{sec:evi}. All other citations are further sources, comparisons, conventions or context, and enter no proof.

%-------------------------------------------------------------------------------
\subsection{What is not claimed}
\label{sec:int-notclaimed}

\begin{enumerate}
\item Nothing is claimed about the quintic threefold or its graph sums; the formal quintic is a different theory. The conifold gap of the quintic threefold is \cite[Theorem~2]{CGYZ26}.
\item Nothing is claimed for frames other than those of \cref{def:set-frames} and the rescalings of their two periods \Lr{sec:sgn-lit}, for $g=1$, or away from $u=0$.
\item Nothing is claimed about the analytic continuation of the series $F_g(q)$ to a neighbourhood of $u=0$. The conifold expansion $\Cexp_{\alpha,\beta}$ is a formal change of frame on the ring $\Nring$ (\cref{def:set-confexp}), applied to the graph-sum representative $F_g^{\Nring}$ of $F_g$ (\cref{def:mai-Fcf}); this is the sense of ``conifold expansion'' and of ``conifold gap'' in \cref{cor:int-gap}.
\item The second key statement of \cite[Remark~1.9]{CGYZ26} is not addressed (the Model Lemma, \cref{def:red-ML}, concerns the model), and \cref{cor:int-gap} is not claimed to be literally the first, which is named there but not stated.
\item No novelty is claimed for the formal stationary-phase lemmas of \Lr{sec:fsp}, proved there for completeness, for the formula of $\Ginf$ \cite[Lemma~7.7, (37)]{GJR18}, or for the change of chart \cite[\S10]{MR1901075}.
\item Topological recursion is not applied to a mirror curve of the quintic or of the formal quintic, and \cite{FCCZ26} is not used.
\end{enumerate}

%-------------------------------------------------------------------------------
\subsection{Strategy and plan}
\label{sec:int-strategy}

\emph{Reduction to the model (\cref{sec:gau,sec:tra}).} In a conifold frame the rows satisfy a cyclic system with one simple pole at $u=0$ (\cref{prop:gau-quintic}); the model pulled back by $v(u)$ has the same product of coefficients, which fixes $\tau$ (\cref{prop:gau-equalprod}). Theorem~U gives the regular gauge between them, and the gauge identity singles out the symplectic gauge $\Pc$ (\cref{thm:gau-G,prop:gau-symp}). Its edges lie in $\sKo\tensor\sKo$ for a slice $\sK$ defined by the vanishing solution of the model (\cref{prop:tra-can}), and Theorem~R (\cref{thm:red-main}) reduces the regularity of the edge terms of the flat sum of $\Pc.\Omega^{\mathrm{mod}}_{v(u)}$, and the vertex law, to the Model Lemma on $\sK$ (\cref{def:red-ML}).

\emph{The model and the normalisation (\cref{sec:mod,sec:flh,sec:con,sec:gam}).} The Model Lemma is proved by topological recursion on the model curve (\cref{thm:mod-tr,thm:mod-MLK}). The comparison theorem (\cref{thm:flh-main}) identifies $T_\Gamma$ with the flat term, graph by graph, if $\Dn_0\equiv1\bmod z^{3g-2}$; Theorem~P (\cref{thm:con-Pa,thm:con-Pb}) and the glue (\cref{thm:con-glue}) turn $\Dn=1$ into part (a) of \cref{thm:int-gamma} (\cref{thm:gam-D}). \Cref{sec:mai} assembles the parts, and \cref{sec:evi} describes the computations that accompanied the work.

\subsection*{Acknowledgements}

This work was supported by the UKRI grant UKRI2394 and by the Leverhulme Trust grant LIP-2023-007.


% Abridged version, Section 2, part 1 of 2: the formal quintic, genus-zero data, the rows and (QRR), the graph sum and (T),
% conifold frames and the conifold expansion. Continued in part 2 (the model, the outer graph sum, Property N(g)).
\section{The formal quintic, the model and the outer graph sum}
\label{sec:set}

This section fixes the objects and states the quoted inputs in the form in which they are used; it follows \Lr{sec:set}. Throughout, $\zeta=\re^{2\pi\ri/5}$, colours $i,j$ and flat indices $a,b$ run over $\bbZ/5$, and $\overline{\mathcal M}_{g,n}$ is the moduli space of stable curves.

%-------------------------------------------------------------------------------
\subsection{The formal quintic and its genus-zero data}
\label{sec:set-fq}

\begin{defn}[The formal quintic; {\Lr{def:set-fq}}, abridged]
\label{def:set-fq}
Let $\bbT=(\bbC^*)^5$ act on $\bbP^4$ with weights $\lambda_0,\dots,\lambda_4$ on $H^0(\bbP^4,\cO(1))$ and $H:=c_1^{\bbT}(\cO(1))$. For $g,d\geq0$ the \emph{formal quintic invariant} $N_{g,d}:=\int_{[\overline{\mathcal M}_g(\bbP^4,d)]^{\mathrm{vir}}}e\bigl(R\pi_*f^*\cO(5)\bigr)$ \cite[\S0.2, (5)]{LP18}, \cite[\S0.2, (5)]{LP20} is defined by virtual localisation \cite{MR1666787} and specialised at $\lambda_i=\zeta^i$; it is a rational number. For $g\geq2$ the \emph{genus-$g$ potential} is $F_g:=\sum_{d\geq0}N_{g,d}\,\Qn^d\in\bbQ\fps{\Qn}$; along the mirror map $\Qn=\re^{T(q)}$ it becomes a power series in $q$. For $2g-2+n>0$ the classes $\Omega_{g,n}(\gamma_1,\dots,\gamma_n):=\sum_{d\geq0}\Qn^d\,\rho_*\bigl(e(R\pi_*f^*\cO(5))\prod_{j}\mathrm{ev}_j^*\gamma_j\cap[\overline{\mathcal M}_{g,n}(\bbP^4,d)]^{\mathrm{vir}}\bigr)$, with the stabilisation map $\rho$, are defined and specialised in the same way; with the \emph{twisted pairing} $\eta(a,b):=\int_{\bbP^4}ab\,e_{\bbT}(\cO(5))$ they are the classes of the cohomological field theory of the formal quintic \cite[\S4.1]{GJR18}, and $F_g=\int_{\overline{\mathcal M}_g}\Omega_{g,0}$. At $\lambda_i=\zeta^i$ the \emph{state space} is $\bbH:=\bbQ[H]/(H^5-1)$, with the flat basis $H^a$, $H^0=1$, and $\eta(H^a,H^b)=5\,\delta_{a+b\equiv3\bmod5}$.
\end{defn}

\emph{Genus-zero data} \Lr{def:set-genus0,lem:set-pf5}. Let $\Dq:=q\,\dd/\dd q$ and let $\logq$ be a formal $\log q$, $\Dq\logq=1$. Put $C_0:=\sum_{d\geq0}\frac{(5d)!}{(d!)^5}q^d$, $\tilde I_1:=5\sum_{d\geq1}\frac{(5d)!}{(d!)^5}\bigl(\sum_{m=d+1}^{5d}\frac1m\bigr)q^d$, $I_1:=C_0\logq+\tilde I_1$, the \emph{mirror map} $T:=I_1/C_0$, $\re^T:=q\exp(\tilde I_1/C_0)$, $C_1:=\Dq T$, and $L:=(1-5^5q)^{-1/5}$, so that $\Dq L=(L^6-L)/5$. The fifth-order and fourth-order operators are
\[
\PFfive:=(z\Dq)^5-1-q\prod_{k=1}^5(5z\Dq+kz),\qquad \PFfour:=\Dq^4-5q\prod_{k=1}^4(5\Dq+k),\qquad \PFfive=z^5\,\Dq\circ\PFfour-1 .
\]
The input (M) is the identity $J(T(q)H,z)=I(q,z)/C_0(q)$ at $\lambda_i=\zeta^i$ between the $J$-function of the twisted theory and the hypergeometric $I$-function \Lr{def:set-genus0,thm:set-mirror}, and (G0) gives the quantum product $*$ at $t=T(q)H$, $\Qn=1$, which is commutative, associative, has the unit $1$ and is Frobenius for $\eta$, with the quantum differential equation \Lr{thm:set-G0}. Along the line $t=T(q)H$ functions of $t$ are functions of $q$, and $\partial_T:=C_1^{-1}\Dq$.

\begin{prop}[The large-radius cyclic structure; {\Lr{prop:set-cyclicLR}}]\genusstatus{AI-proposed, unchecked}
\label{prop:set-cyclicLR}
Assume (M) and (G0), and for (iii) the flatness part of (QRR), \cref{def:set-rows}~(b). Along the line $t=T(q)H$, at $\lambda_i=\zeta^i$ and $\Qn=1$, the following hold.
\begin{enumerate}
\item In the flat basis, $H*H^a=\kappa^{\mathrm{LR}}_aH^{a+1}$ for $a\in\bbZ/5$, where $\kappa^{\mathrm{LR}}=\bigl(1,\ L^5/(C_0^2C_1^3),\ 1,\ C_0/C_1,\ C_0/C_1\bigr)$.
\item Consequently, for every solution $v$ of the quantum differential equation $z\partial_Tv=H*v$, in any differential $\bbQ\fps q$-algebra in which $z$ is a constant, the lowered components $v_a:=\eta(H^a,v)$ satisfy $z\Dq v_a=C_1\kappa^{\mathrm{LR}}_a\,v_{a+1}$ for all $a\in\bbZ/5$.
\item For every colour $i$ the rows of \cref{def:set-rows} satisfy $(z\Dq+\zeta^iL)\Qr^{(i)}_a=C_1\kappa^{\mathrm{LR}}_a\,\Qr^{(i)}_{a+1}$ for $a\in\bbZ/5$.
\end{enumerate}
\end{prop}

\begin{sketch}[Idea of the proof {\Lr{sec:set-lrproof,def:set-rows}}]
A $\bbZ/5$-grading shows that $H*H^a$ is a multiple of $H^{a+1}$. By (M) the coefficients are expressed through the periods, the coefficients of $I/z$: their Wronskian $L^{10}$ gives $\kappa^{\mathrm{LR}}_1$, and an operator identity with $\PFfour$ gives $\kappa^{\mathrm{LR}}_3=C_0/C_1$, through $H^5=1$. Part (ii) is the Frobenius property, and (iii) is (ii) for the column $R(z)\tilde e_i\re^{u_i/z}$, by flatness.
\end{sketch}

\begin{cor}[{\Lr{cor:set-cyclicPF5}}, abridged]\genusstatus{AI-proposed, unchecked}
\label{cor:set-cyclicPF5}
Assume (M) and (G0), and for the last sentence of (b) the flatness part of (QRR).
\begin{itemize}
\item[(a)] The eigenvalues of $H*$ are $\mu_i:=\zeta^iL/C_1$, the idempotents $e_i$ satisfy $\Delta_i:=1/\eta(e_i,e_i)=C_0^2\zeta^{3i}/L^2$, and the canonical coordinates along the line are functions $u_i$ with $\Dq u_i=\zeta^iL$, determined up to additive constants.
\item[(b)] For every solution $v$ of $z\partial_Tv=H*v$ in a differential $\bbQ\fps q$-algebra in which $z$ is a constant, $C_0v_0$ solves $\PFfive$. In particular, for the column $v=R(z)\tilde e_i\re^{u_i/z}$ of \cref{def:set-rows}, the function $\re^{u_i/z}\,z\,L\,\Phi_i(\zh,u)$ solves $\PFfive$: this is the WKB form of colour $i$.
\end{itemize}
\end{cor}

\begin{defn}[Points, variables, Frobenius data; {\Lr{def:set-points,def:set-frob}}]
\label{def:set-frob}
Put $u:=1-5^5q$ and $s:=u^{1/5}$, so that $L=s^{-1}$ and $\Dq=-(1-u)\,\dd/\dd u$, and $\zh:=z/s^4=L^4z$, the hatted variable. The \emph{large-radius point} is $q=0$ ($u=1$), the \emph{conifold point} is $u=0$, and the \emph{orbifold point} is $q=\infty$ ($L=0$). The \emph{normalised canonical basis} is $\tilde e_i:=\Delta_i^{1/2}e_i$ with $\Delta_i^{1/2}:=\zeta^{4i}\,C_0/L$, so $\eta(\tilde e_i,\tilde e_j)=\delta_{ij}$. The \emph{topological field theory of the formal quintic} is the TFT $\TFT$ with $\TFT_{g,n}(\tilde e_i,\dots,\tilde e_i):=\Delta_i^{(2g-2+n)/2}$ and $\TFT_{g,n}$ vanishing on mixed colours; its unit is $1_{\TFT}:=\sum_i\Delta_i^{-1/2}\tilde e_i=1$.
\end{defn}

%-------------------------------------------------------------------------------
\subsection{The \texorpdfstring{$R$}{R}-matrix, the graph sum and the quoted inputs (QRR), (T)}
\label{sec:set-rows}

\begin{defn}[The $R$-matrix and its rows; {\Lr{def:set-rows}}]
\label{def:set-rows}
Let $R(z)=1+\sum_{k\geq1}R_kz^k$, $R_k\in\End(\bbH)\otimes\bbQ(\zeta)\fps q$, be the $R$-matrix of the semisimple cohomological field theory of the formal quintic along the line $t=T(q)H$ \cite[\S4.1]{GJR18}, that is, the $R$-matrix of the quoted input (T), \cref{thm:set-GT}. It has the following properties, of which (a) is part of (T), (b) is part of (QRR), and (c) is Givental's convention: (a) $R$ is symplectic, $R(-z)^T\eta\,R(z)=\eta$ in the flat basis; (b) for every colour $i$ the column $R(z)\tilde e_i\re^{u_i/z}$ solves the quantum differential equation $z\partial_Tv=H*v$, that is, $w:=R(z)\tilde e_i$ satisfies $z\partial_Tw+\mu_iw=H*w$; (c) in the graph sums of \cref{def:set-graphsum}, $R$ acts with $R^{-1}(\psi)$ on the legs. The \emph{rows} are the lowered components $\Qr^{(i)}_a(z):=\eta(H^a,R(z)\tilde e_i)\in\bbQ(\zeta)\fps q\fps z$; $\Qr$ is the matrix $(\Qr^{(i)}_a)$, rows indexed by $a$, columns by $i$. With $\zh=L^4z$ we define $\Phi_i(\zh,u):=C_0\,\Qr^{(i)}_0(z)/(\zeta^iL)$ and the coefficients $\mathrm{ph}_n(u)$ by $\Phi_0(\zh,u)=\sum_{n\geq0}\mathrm{ph}_n(u)\,\zh^n$; then $\mathrm{ph}_0=1$. The \emph{flat amplitudes} are $A_n:=\coef{z^n}\bigl(C_0\,\Qr^{(0)}_0\bigr)=L^{4n+1}\,\mathrm{ph}_n(L^{-5})$.
\end{defn}

\begin{thm}[Flatness of the $R$-matrix and its value at $q=0$, (QRR); quoted; {\Lr{thm:set-qrr}}]\genusstatus{AI-proposed, unchecked}
\label{thm:set-qrr}
For a finite multiset $\mathbf w$ of non-zero numbers put $\hw(z):=\exp\bigl(-\sum_{k\geq1}\bk_k\,z^{2k-1}\sum_{a\in\mathbf w}a^{1-2k}\bigr)$, $\bk_k:=B_{2k}/(2k(2k-1))$, and $\bw_i:=\{\zeta^i-\zeta^j\}_{j\neq i}\cup\{-5\zeta^i\}$. The $R$-matrix of \cref{def:set-rows} has coefficients in $\End(\bbH)\otimes\bbQ(\zeta)\fps q$ and satisfies \cref{def:set-rows}~(b); and $R(q=0)=\diag(h_{\bw_0},\dots,h_{\bw_4})$ in the basis $\tilde e_i(q=0)$, to all orders in $z$.
\end{thm}

The value is used only in the form $\Phi_i(z,1)=h_{\bw_i}(z)$, for every colour $i$; in particular, with $h:=h_{\bw_0}$ and $h_n:=\coef{z^n}h$, $\mathrm{ph}_n(1)=h_n$. Since $\bw_i=\zeta^i\bw_0$, $h_{\bw_i}(z)=h_{\bw_0}(\zeta^{-i}z)$. Assuming (M) and (G0), flatness and the value at $q=0$ determine $R$, and imply (a) \Lr{lem:set-hrot,lem:set-Runique}; so no uniqueness of the $R$-matrix in (T) is used, and (a) is not an independent input.

\begin{defn}[The Givental--Teleman graph sum; {\Lr{def:set-graphsum}}, abridged]
\label{def:set-graphsum}
Let $\Omega$ be a cohomological field theory with unit $1$ \cite[Def.~0.3]{PPZ15} on a free module $V$ of finite rank with a perfect symmetric pairing $\eta$ over a commutative $\bbQ$-algebra, and $R(z)=1+O(z)\in\End(V)\fps z$ symplectic. The ancestor correlators are $\langle\psi^{k_1}a_1\cdots\psi^{k_n}a_n\rangle_h:=\int_{\overline{\mathcal M}_{h,n}}\Omega_{h,n}(a_1,\dots,a_n)\prod_j\psi_j^{k_j}$; for a point $q\in V\fps z$ of the loop space the dilaton-shifted potential is $\overline F^\Omega_h(q):=\sum_{n:\,2h-2+n>0}\frac1{n!}\langle(q+z1)^n\rangle_h$, with $z^ka\mapsto\psi^ka$, and $\dd^n\overline F^\Omega_h(q)$ is its $n$-th derivative, a symmetric form on $V\fps z$. The \emph{leaf point} and the \emph{edge} of $R$ are $q_R(z):=R^{-1}(z)(-z1)=-z1+\ttr_R(z)$ and
\[
E_R(\psi,\psi'):=\frac{\eta^{-1}-\bigl(R^{-1}(\psi)\otimes R^{-1}(\psi')\bigr)\eta^{-1}}{\psi+\psi'}\in(V\otimes V)\fps{\psi,\psi'},
\]
where $\eta^{-1}\in V\otimes V$ is the inverse bivector and $\ttr_R(z):=z1-zR^{-1}(z)1\in z^2V\fps z$ is the translation of \cite[Def.~2.13]{PPZ15}. For a stable graph $\Gamma$ without legs, with vertices $v$ of genus $g_v$ and valence $n_v$, and a symmetric bivector series $E$, the contraction $\mathrm{Cont}_\Gamma[E,q;\Omega]$ is $|\Aut\Gamma|^{-1}$ times the contraction of the vertex tensors $\dd^{n_v}\overline F^\Omega_{g_v}(q)$, one slot for each half-edge, with one copy of $E(\psi,\psi')$ for each edge. The sum $\sum_\Gamma\mathrm{Cont}_\Gamma[E_R,q_R;\Omega]$ over the leg-free stable graphs of genus $g$ is the \emph{graph sum of genus $g$}; $R.\Omega$ denotes the unit-preserving action of \cite[Defs.~2.2, 2.5, 2.13]{PPZ15}.
\end{defn}

By the formula \cite[(8)]{PPZ15} for the pull-back of $\psi$ classes, for $g\geq2$, $\mathrm{Cont}_\Gamma[E_R,q_R;\Omega]$ is the term of $\Gamma$ in $\int_{\overline{\mathcal M}_g}(R.\Omega)_{g,0}$ \Lr{lem:set-pushins}.

\begin{thm}[The reconstruction (T); quoted; {\Lr{thm:set-GT}}]\genusstatus{AI-proposed, unchecked}
\label{thm:set-GT}
The input (T) of \cref{sec:int-inputs} holds: the classes $\Omega_{g,n}$ of \cref{def:set-fq} with $\Qn=\re^{T(q)}$ are the cohomological field theory at the point $t=T(q)H$ with $\Qn=1$, and it equals $R.\TFT$, with $\TFT$ of \cref{def:set-frob} and the symplectic $R$-matrix $R=1+O(z)$ of \cref{def:set-rows}, with coefficients in $\End(\bbH)\otimes\bbQ(\zeta)\fps q$.
In particular, for $g\geq2$, $F_g=\sum_\Gamma\mathrm{Cont}_\Gamma[E_R,q_R;\TFT]$ as power series in $q$, the sum running over the leg-free stable graphs of genus $g$.
\end{thm}

%-------------------------------------------------------------------------------
\subsection{Conifold frames and the conifold expansion}
\label{sec:set-frames}

In the variable $u$, $\PFfour=u\Dq^4-(1-u)\bigl(2\Dq^3+\frac75\Dq^2+\frac25\Dq+\frac{24}{625}\bigr)$ with $\Dq=-(1-u)\,\dd/\dd u$.

\begin{lem}[The Frobenius basis at the conifold point; {\Lr{lem:set-frobbasis}}, abridged]\genusstatus{AI-proposed, unchecked}
\label{lem:set-frobbasis}
$\PFfour$ maps $\bbQ\fps u$ into itself, its local exponents at $u=0$ are $0,1,1,2$, and a power-series solution is determined by its coefficients of $u^0,u^1,u^2$; $f_k$ ($k=0,1,2$) denotes the solution with $f_k=u^k+O(u^3)$. Let $\ell$ be a formal $\log u$. Up to a non-zero factor, the \emph{vanishing period} $\varpi_1:=f_1+\frac7{10}f_2=u+O(u^2)$ is the only power-series solution $g$ for which a solution $g\ell+h$, $h\in\bbQ\fps u$, exists.
\end{lem}

\begin{defn}[Conifold frames; {\Lr{def:set-frames}}]
\label{def:set-frames}
A \emph{conifold frame} is given by $(\alpha,\beta)\in\bbQ^2$: $\varpi_0:=f_0+\alpha\varpi_1+\beta f_2$ and $\tc:=\varpi_1/\varpi_0$. Then $\varpi_0=1+\alpha u+O(u^2)$ and $\tc=u+(\frac7{10}-\alpha)u^2+O(u^3)$, so $\varpi_0(0)=1$, $\tc'(0)=1$ and $\Dq\tc(0)=-1$.
\end{defn}

\begin{defn}[The differential ring of the periods and the conifold expansion; {\Lr{def:set-confexp}}, abridged]
\label{def:set-confexp}
Let
\[
\Nring:=\bbQ[L^{\pm1},\logq,c_0,c_1,c_2,c_3,i_0,i_1,i_2][c_0^{-1},\mathsf W^{-1}],\qquad\mathsf W:=c_0^2+c_0i_1-c_1i_0,
\]
with the derivation $\Dq$ fixed on the generators by $\Dq L:=(L^6-L)/5$, $\Dq\logq:=1$, $\Dq c_j:=c_{j+1}$, $\Dq i_j:=i_{j+1}$, and by the relations $\PFfour c_0=0$ (with $u=L^{-5}$) and $\Conc(c_0,I_1)=0$, which determine $\Dq c_3$ and $\Dq i_2$; here $I_1:=c_0\logq+i_0$ and $\Conc(f,g):=\sum_{k=1}^4\sum_{j=0}^{k-1}(-1)^j\,\Dq^j(A_kg)\,\Dq^{k-1-j}f$ is the concomitant of $\PFfour=\sum_kA_k\Dq^k$, for which $g\,\PFfour f-f\,\PFfour g=\Dq\,\Conc(f,g)$ \Lr{def:gau-conc,lem:gau-adj,lem:gau-conc}. Put $C_0:=c_0$, $T:=I_1/C_0$ and $C_1:=\Dq T=\mathsf W/c_0^2$. The \emph{evaluation} $\ev\colon\Nring\to\bbQ\fps q[\logq]$, $c_j\mapsto\Dq^jC_0$, $i_j\mapsto\Dq^j\tilde I_1$, $\logq\mapsto\logq$, $L\mapsto L$, is a homomorphism of differential rings. The \emph{conifold expansion} in the frame $(\alpha,\beta)$ is the ring homomorphism $\Cexp_{\alpha,\beta}\colon\Nring\to\bbQ\lser s$ with
\[
c_j\mapsto\Dq^j\varpi_0,\qquad i_j\mapsto\Dq^j\varpi_1-j\,\Dq^{j-1}\varpi_0+\Xi\,\Dq^j\varpi_0,\qquad\logq\mapsto\log(1-u)=-\Xi,\qquad L\mapsto s^{-1},
\]
where $\Xi:=-\log(1-u)=\sum_{n\geq1}u^n/n$. It commutes with $\Dq$ and maps $C_0$, $I_1$, $T$, $C_1$ to $\varpi_0$, $\varpi_1$, $\tc$, $\Dq\tc$ \Lr{lem:gau-confexp}; its $\bbQ(\zeta)$-linear extension is denoted by the same letter. An element of $\Nring\otimes\bbQ(\zeta)$, or a series over it, is \emph{regular at the conifold point} if its conifold expansion lies in $\bbQ\fps u$, coefficientwise.
\end{defn}

The conifold expansion substitutes the periods of a conifold frame for $C_0$ and $I_1$ in a representative in the ring $\Nring$; it is a formal change of frame, and no analytic continuation is involved. Whether $\ev$ is injective is neither needed nor proved, and identities between power series in $q$, such as the symplecticity of $R$, are not transferred to the conifold expansion.

\begin{defn}[The natural variable and the canonical coordinates; {\Lr{def:set-tau,lem:set-canon}}]
\label{def:set-tau}
Put
\[
\tau(u):=\frac1{{}_2F_1(1,\frac45;\frac95;u)}=\Bigl(\sum_{n\geq0}\frac{4u^n}{5n+4}\Bigr)^{-1}\in1+u\bbQ\fps u,\qquad v(u):=u\,\tau(u)^{-5/4}\in u+u^2\bbQ\fps u,
\]
and $\varsigma:=v^{1/5}=s\,\tau^{-1/4}$, so that $z/\varsigma^4=\tau(u)\zh$. In $\bbQ(\zeta)\lser s$ the solutions $y$ of $\Dq y=\zeta^iL$ are the series $y=c-\frac54\zeta^iu^{4/5}/\tau(u)$, $c\in\bbQ(\zeta)$; so in the conifold chart the canonical coordinates are $u_i=u_i^{\mathrm c}-\frac54\zeta^iu^{4/5}/\tau(u)$, with constants $u_i^{\mathrm c}$, and $\tau(u)\,\zh=-\frac54\,\zeta^iz/(u_i-u_i^{\mathrm c})$.
\label{lem:set-canon}
\end{defn}



% Abridged version, Section 2, part 2 of 2: the model, its R-matrix and its cohomological field theory; the outer graph sum; Property N(g).

%-------------------------------------------------------------------------------
\subsection{The model}
\label{sec:set-model}

\begin{defn}[The model; {\Lr{def:set-model}}]
\label{def:set-model}
The \emph{model} is the Landau--Ginzburg potential $\Wv(p):=p^4/4+v/p$, with the parameter $v$ and the primitive form $\dd p$. Its formal integrals $I_m(v,z):=\int p^m\re^{-\Wv(p)/z}\dd p$ ($m\in\bbZ$) are symbols subject to the two relations
\[
\text{(Mod1)}\quad z\,\partial_vI_m=-I_{m-1},\qquad\text{(Mod2)}\quad mz\,I_{m-1}-I_{m+3}+v\,I_{m-2}=0,
\]
which are realised by residues at $p=0$ (\cref{def:gau-van}) and by the formal stationary phase at a critical point (\cref{sec:con-fsp}). The charges and the signs are $\nu:=(0,-1,3,2,1)$ and $\varepsilon:=(1,-1,-1,1,1)$, with indices in $\bbZ/5$; thus $\nu_a\equiv-a\bmod5$ and $\nu_a+\nu_{3-a}=2$. The \emph{rows} of the model are $Y_b:=\varepsilon_bI_{\nu_b}$, and the dictionary of classes is $H^b\leftrightarrow\varepsilon_b\,p^{\nu_b}\dd p$. The critical points of $\Wv$ are the $p$ with $p^5=v$; colour $i$ corresponds to the critical point $p=\zeta^{-i}\varsigma$, $\varsigma^5=v$, where $\re^{-\Wv/z}=\re^{U_i/z}$ with $U_i:=-\frac54\zeta^i\varsigma^4$.
\end{defn}

\begin{prop}[The model system and the T-system; {\Lr{prop:set-modrows}}, abridged]\genusstatus{AI-proposed, unchecked}
\label{prop:set-modrows}
\begin{itemize}
\item[(a)] If the symbols $I_m$ satisfy (Mod1) and (Mod2), the rows satisfy the \emph{model system} $z\,\partial_vY_b=b^{\mathrm m}_b\,Y_{b+1}$ ($b\in\bbZ/5$), with $b^{\mathrm m}:=\bigl(1,-\frac1v,1,-1,-1\bigr)$.
\item[(b)] Fix a colour $i$ and put $t:=z/\varsigma^4$. A column $\re^{U_i/z}y$ with $y_b=\varsigma^{\nu_b-1}F_b(t)$, $F_b\in\bbQ(\zeta)\fps t$ (a \emph{homogeneous WKB column of colour $i$}) solves the model system if and only if, with $\mathcal T^{(i)}_k:=-\zeta^i+\frac t5\bigl(k-4t\frac{\dd}{\dd t}\bigr)$,
\[
F_1=\mathcal T^{(i)}_{-1}F_0,\quad F_2=-\mathcal T^{(i)}_{-2}F_1,\quad F_3=\mathcal T^{(i)}_2F_2,\quad F_4=-\mathcal T^{(i)}_1F_3,\quad F_0=-\mathcal T^{(i)}_0F_4.
\]
This T-system has a unique solution $F^{(i)}$ with $F^{(i)}_0(0)=1$, and $F^{(i)}_b(0)=\varepsilon_b\zeta^{ib}$. For colour $0$ we write $\Fs_b:=F^{(0)}_b\in\bbQ\fps t$; thus $\Fs_b(0)=\varepsilon_b$, and $\Fs_0=1-\frac3{20}t+\frac9{800}t^2+O(t^3)$.
\end{itemize}
\end{prop}

\begin{defn}[The model $R$-matrix and the model cohomological field theory; {\Lr{def:set-M0,lem:set-M0sympl,def:set-modcohft}}]
\label{def:set-M0}
Let $\Rhat^{\mathrm m}(t)$ be the $5\times5$ matrix with the entries $\Rhat^{\mathrm m,(i)}_b(t):=\zeta^{i(b+1)}\Fs_b(\zeta^{-i}t)$ (row $b$, colour column $i$), let $\Psi_0:=\Rhat^{\mathrm m}(0)$, with the entries $(\Psi_0)_{bi}=\varepsilon_b\zeta^{i(b+1)}$, and put $\Mo(t):=\Psi_0^{-1}\Rhat^{\mathrm m}(t)$, the model $R$-matrix in Teleman's colour-colour form. Then $\Psi_0^T\eta^{-1}\Psi_0=1$, $\Mo(0)=1$, the columns $\re^{U_i/z}\varsigma^{\nu_b-1}\Rhat^{\mathrm m,(i)}_b(z/\varsigma^4)$ are homogeneous WKB solutions of colour $i$ of the model system, and $\Mo$ is symplectic: $\Mo(-t)^T\Mo(t)=1$. The \emph{model TFT} $\TFT^{\mathrm m}_v$ is the TFT on the colour space, with the orthonormal basis $\tilde e_i$, defined by $\TFT^{\mathrm m}_{v;h,n}(\tilde e_i,\dots,\tilde e_i):=(\zeta^{4i}\varsigma)^{2h-2+n}$ and $\TFT^{\mathrm m}_{v;h,n}=0$ on mixed colours. Let $\Psi^{\mathrm m}_v$ be the map from the colour space to $\bbH$ with the lowered entries $\eta(H^b,\Psi^{\mathrm m}_v\tilde e_i):=\varsigma^{\nu_b-1}(\Psi_0)_{bi}$; it is an isometry. The \emph{model cohomological field theory} on $(\bbH,\eta)$ is
\[
\Omega^{\mathrm{mod}}_v:=(\Psi^{\mathrm m}_v)_*\bigl(\Mo(z/\varsigma^4).\TFT^{\mathrm m}_v\bigr),\qquad(A_*\Omega)_{h,n}:=\Omega_{h,n}\circ(A^{-1})^{\otimes n},
\]
with the unit-preserving action of \cite[Defs.~2.5, 2.13]{PPZ15}. Its ancestor correlators $\langle\cdots\rangle_h(v)$ and its dilaton-shifted potentials $\overline F_h(q)$, with the base point $q=-z1$, are those of \cref{def:set-graphsum}. The \emph{free energies} are $\Fm_h:=\langle\ \rangle_h(v)$ ($h\geq2$).
\label{def:set-modcohft}
\end{defn}



\begin{lem}[Homogeneity, dilaton equation; {\Lr{lem:set-modhom,lem:set-dilaton}}, abridged]\genusstatus{AI-proposed, unchecked}
\label{lem:set-dilaton}
\begin{enumerate}
\item Every ancestor correlator of the model is $\langle\prod_j\psi^{k_j}H^{b_j}\rangle_h(v)=C\,v^e$ with $C\in\bbQ$, $e:=2-2h+\sum_j\wt(k_j,b_j)$ and $\wt(k,b):=(4k+\nu_b-4)/5$, and $C=0$ unless $e\in\bbZ$. A correlator is \emph{regular} if $e\geq0$; then it lies in $\bbQ[v]$.
\item $\Omega^{\mathrm{mod}}_v$ has the flat unit $H^0=1$. For a cohomological field theory with flat unit $1$ and $2h-2+n>0$, the insertion of $\psi1$ multiplies a correlator by $2h-2+n$: $\langle\psi1\cdot X\rangle_h=(2h-2+n)\langle X\rangle_h$ for $X=\psi^{k_1}a_1\cdots\psi^{k_n}a_n$.
\item Let $\Omega$ be a cohomological field theory with flat unit $1$ on a free module of finite rank over a complete local $\bbQ$-algebra with maximal ideal $\mathfrak m$, and $\varkappa\in1+\mathfrak m$; or over $\bbQ(\zeta)\lser s$, and $\varkappa\in1+u\bbQ\fps u$. Then, for $2h-2+n>0$, $\dd^n\overline F_h(q/\varkappa)=\varkappa^{2h-2+n}\,\dd^n\overline F_h(q)$, the left side being defined coefficientwise in $q+z1$ as a convergent series in $1-\varkappa^{-1}$.
\end{enumerate}
\end{lem}

The dilaton equation in (ii) uses \cite[(1.7)]{AC96}; in the applications of (iii), $\varkappa$ is the dilaton factor $d\in\bbQ\fps u$ of the canonical gauge, $d(0)=1$.

%-------------------------------------------------------------------------------
\subsection{The outer graph sum and Property N(\texorpdfstring{$g$}{g})}
\label{sec:set-hat}

A conifold frame is fixed, and the rows, which lie in $\Nring\otimes\bbQ(\zeta)$ (\cref{cor:gau-rowsNN}), are expanded in it: $\Qr^{(i)}_a(z)$, and every series below, is a series in $z$ over $\bbQ(\zeta)\lser s$.

\begin{defn}[The outer graph sum; {\Lr{def:set-hat,def:set-outer}}, abridged]
\label{def:set-outer}
Let $M:=\Mo(\tau(u)\zh)$, taken exactly in $u$: with $\zh=z/s^4$ it is a power series in $z$ over $\bbQ(\zeta)\lser s$, with $M(0)=1$ and $M(-z)^TM(z)=1$. Let $\Qr$ be the matrix of the lowered rows of \cref{def:set-rows} expanded in the frame, and $R':=\Qr\,M^{-1}$, a series of maps from the colour space (orthonormal basis $\tilde e_i$) to $\bbH$. Its constant term $R'(0)=\Qr(0)$ is an isometry which maps $1_{\TFT}$ to $1$ \Lr{lem:mai-Psi}. The \emph{outer graph sum} has the vertices $(M.\TFT)_{g,n}$, the classes of the cohomological field theory $M.\TFT$ on the colour space \cite[Def.~2.13]{PPZ15}, where $\TFT$ is the TFT of \cref{def:set-frob} expanded in the frame, $\TFT_{g,n}(\tilde e_i,\dots,\tilde e_i)=(\zeta^{4i}s\varpi_0)^{2g-2+n}$, and the edges and leaves of $R'$: with $N(\psi,\psi'):=1-\bigl(R'^{-1}(\psi)\otimes R'^{-1}(\psi')\bigr)\eta^{-1}$, where $1=\sum_i\tilde e_i\otimes\tilde e_i$ and $\eta^{-1}\in\bbH\otimes\bbH$ is the inverse bivector of $(\bbH,\eta)$,
\[
E_{R'}(\psi,\psi'):=\frac{N(\psi,\psi')-\frac12N(-\psi',\psi')-\frac12N(\psi,-\psi)}{\psi+\psi'},\qquad q_{R'}(z):=R'^{-1}(z)(-z1),
\]
with the leaf insertion $q_{R'}+z1_{\TFT}\in z^2(\text{colour space})\fps z$. The numerator vanishes at $\psi'=-\psi$, so $E_{R'}$ is a power series in $\psi,\psi'$. If $R'$ is symplectic, $R'(-z)^T\eta^{-1}R'(z)=1$, then $N(-\psi',\psi')=N(\psi,-\psi)=0$, $E_{R'}$ is the edge of \cite[Def.~2.2]{PPZ15}, and with $X:=R'(0)^{-1}R'$ the outer graph sum is the leg-free part of the unit-preserving action $X.(M.\TFT)$ on the colour space. For a leg-free stable graph $\Gamma$ of genus $g\geq2$ the \emph{outer term} is $T_\Gamma:=\mathrm{Cont}_\Gamma[E_{R'},q_{R'};M.\TFT]\in\bbQ(\zeta)\lser s$. The \emph{vertex graph} is the stable graph with one vertex, of genus $g$, and no edge; its outer term contains all the leaf insertions of $R'$. The \emph{hatted term} $G_\Gamma$ is the same contraction formed with the hatted classes $\hat H^a:=s^{1-\nu_a}H^a$, the hatted rows $\hat\Qr^{(i)}_a(\zh):=s^{1-\nu_a}\Qr^{(i)}_a(s^4\zh)$, the hatted TFT $\hat\TFT_{g,n}:=s^{-(2g-2+n)}\TFT_{g,n}$, $\hat M(\zh):=\Mo(\tau\zh)$ and the unit $\hat1:=\hat H^0=s1$, in the variable $\zh$.
\end{defn}

\begin{defn}[Property N($g$); {\Lr{def:set-N}}]
\label{def:set-N}
Let $g\geq2$ and $c_g:=(-1)^{g-1}B_{2g}\,5^{1-g}/(2g(2g-2))=|B_{2g}|/(2g(2g-2)\,5^{g-1})$. \emph{Property N($g$)} holds if, for every conifold frame of \cref{def:set-frames}, the outer terms of \cref{def:set-outer}, with $M=\Mo(\tau(u)\zh)$ exactly in $u$, satisfy: (i) $T_\Gamma\in\bbQ\fps u$ for every leg-free stable graph $\Gamma$ of genus $g$ with at least one edge; (ii) $T_\Gamma\in c_g\,\tc^{2-2g}+\bbQ\fps u$ for the vertex graph. Here $c_2=\frac1{1200}$, $c_3=\frac1{25200}$, $c_4=\frac1{180000}$, $c_5=\frac1{660000}$.
\end{defn}


% Abridged version, Section 3, part 1 of 2: the cyclic form of PF4, the two cyclic systems, Theorem U. Continued in part 2.
\section{Cyclic systems, Theorem U and the canonical gauge}
\label{sec:gau}

This section follows \Lr{sec:gau}. Here $\Dq=-(1-u)\,\dd/\dd u$, a prime denotes $\dd/\dd u$, $z$ is a constant, $E_{ab}$ ($a,b\in\bbZ/5$) are the matrix units and $\mathbf e_a$ the standard basis vectors; solutions, gauges and rows are in lowered components.

%-------------------------------------------------------------------------------
\subsection{The cyclic form of the fourth-order operator}
\label{sec:gau-cyc}

Let $\bbK:=\bbQ\lser u$, let $\ell$ be a formal $\log u$ ($\Dq\ell=1-u^{-1}$), let $y_3=\varpi_1\ell+\rho_3$, $\rho_3\in u^3\bbQ\fps u$, be the logarithmic solution of $\PFfour$ (\cref{lem:set-frobbasis}), and $\mathcal S:=\operatorname{span}_\bbQ(f_0,f_1,f_2,y_3)$; thus $\mathcal S\cap\bbK=\operatorname{span}(f_0,f_1,f_2)$ is the space of power-series solutions. For $y_0,y_1\in\bbK$ with $y_0\neq0$ and $t:=y_1/y_0$ not constant, the \emph{cyclic operator} is
\[
\Mcyc(y_0,y_1):=\frac1{y_0}\circ\Dq\circ\frac1{\Dq t}\circ\Dq\circ\bigl(u\,y_0^2(\Dq t)^2\bigr)\circ\Dq\circ\frac1{\Dq t}\circ\Dq\circ\frac1{y_0}\ \in\bbK[\Dq].
\]

\begin{thm}[The cyclic form of $\PFfour$, (CYC); {\Lr{thm:gau-cyc}}, abridged]\genusstatus{AI-proposed, unchecked}
\label{thm:gau-cyc}
Let $y_0\in\mathcal S\cap\bbK$ be non-zero and not proportional to $\varpi_1$. Then, as operators in $\bbK[\Dq]$, $\PFfour=\Mcyc(y_0,\varpi_1)$. In particular $\PFfour=\Mcyc(\varpi_0,\varpi_1)$ in every conifold frame of \cref{def:set-frames}.
\end{thm}

\begin{sketch}[Idea of the proof {\Lr{sec:gau-cycproof}}]
$\PFfour$ is formally self-adjoint, so its concomitant $\Conc$ is a $\bbQ$-valued alternating form on $\mathcal S$. The automorphism $\ell\mapsto\ell+1$ maps $y_3$ to $y_3+\varpi_1$ (a formal Picard--Lefschetz transformation); hence $\Conc(f,\varpi_1)=0$ for every power-series solution $f$, and every plane in $\mathcal S\cap\bbK$ that contains $\varpi_1$ is isotropic. A transversality identity then gives the relations of special geometry: with $\de_t=(\Dq t)^{-1}\Dq$, the functions $u\,y_0^2(\Dq t)^3\,\de_t^2(y/y_0)$, $y\in\mathcal S$, lie in $\operatorname{span}_\bbQ(1,t)$. So $\Mcyc(y_0,\varpi_1)$ annihilates $\mathcal S$ and has the leading coefficient $u$ of $\PFfour$; the difference, of order at most $3$, vanishes.
\end{sketch}

\begin{cor}[The cyclic system is $\PFfive$; {\Lr{cor:gau-cycsys}}, for $y_1=\varpi_1$]\genusstatus{AI-proposed, unchecked}
\label{cor:gau-cycsys}
Let $y_0$ be as in \cref{thm:gau-cyc}, $t=\varpi_1/y_0$, and $\kappa=(\kappa_0,\dots,\kappa_4):=\bigl(\Dq t,\ L^5/(y_0^2(\Dq t)^2),\ \Dq t,\ y_0,\ y_0\bigr)$, $L^5:=u^{-1}$. Let $(\cR,\Dq)$ be a differential ring extending $(\bbK,\Dq)$ and $z\in\cR$ with $\Dq z=0$.
\begin{enumerate}
\item For $Y_0,\dots,Y_4\in\cR$ the cyclic system $z\Dq Y_a=\kappa_aY_{a+1}$ ($a\in\bbZ/5$) holds if and only if $y_0Y_0$ solves $\PFfive$ and $Y_{a+1}=z\Dq Y_a/\kappa_a$ for $a=0,\dots,3$.
\item If $z$ is invertible in $\cR$, then for every $w\in\cR$ the same holds with $z\Dq$ replaced by $z\Dq+w$ in the system, in the relations $Y_{a+1}=z\Dq Y_a/\kappa_a$ and in $\PFfive$.
\end{enumerate}
\end{cor}



%-------------------------------------------------------------------------------
\subsection{The quintic system and the pulled-back model}
\label{sec:gau-systems}

\begin{defn}[Cyclic systems with one simple pole; regular gauges; {\Lr{def:gau-onepole}}]
\label{def:gau-onepole}
Let $\Bbbk$ be a field of characteristic $0$. A \emph{cyclic system with one simple pole} over $\Bbbk\fps u$ is a system $z\de_uY=AY$ with $A=\sum_{a\in\bbZ/5}b_aE_{a,a+1}$ whose coefficients have the \emph{one-pole shape}: $b_a\in\Bbbk\fps u^\times$ for $a\neq1$ and $ub_1\in\Bbbk\fps u^\times$. Its \emph{residue} is $r:=(ub_1)(0)$ and its \emph{product} is $\Bprod:=\prod_ab_a$. It is \emph{self-dual} if $b_0=b_2$ and $b_3=b_4$, i.e.\ $\eta A=A^T\eta$. For a second such system $\tilde A=\sum_a\tilde b_aE_{a,a+1}$, the condition (EP) of \emph{equal products} is $\prod_ab_a=\prod_a\tilde b_a$. The \emph{gauge equation} from $\tilde A$ to $A$ is $zP'=AP-P\tilde A$ for $P=\sum_{k\geq0}z^kP_k$, $P_k\in\mathfrak{gl}_5(\Bbbk\fps u)$; then $z\de_u(P\tilde Y)=AP\tilde Y$ whenever $z\de_u\tilde Y=\tilde A\tilde Y$. $\Gau(A,\tilde A)$ denotes the set of all solutions; it is a $\Bbbk\fps z$-module. A \emph{regular gauge} is an element of $\Gau(A,\tilde A)$ with $P_0\in\mathrm{GL}_5(\Bbbk\fps u)$. $P$ is \emph{graded} if $P_k=\sum_bp_{k,b}E_{b+k,b}$ for every $k$.
\end{defn}

\begin{prop}[The quintic cyclic system; {\Lr{prop:gau-quintic}}, abridged]\genusstatus{AI-proposed, unchecked}
\label{prop:gau-quintic}
Let $(\alpha,\beta)\in\bbQ^2$ give a conifold frame, $\tc=\varpi_1/\varpi_0$, and let $\kappa:=\bigl(\Dq\tc,\ L^5/(\varpi_0^2(\Dq\tc)^2),\ \Dq\tc,\ \varpi_0,\ \varpi_0\bigr)$, $L^5=u^{-1}$, be the vector of \cref{cor:gau-cycsys} for $y_0=\varpi_0$.
\begin{enumerate}
\item $\kappa_a=\Cexp_{\alpha,\beta}(C_1\kappa^{\mathrm{LR}}_a)$ for $a\in\bbZ/5$. Hence the conifold expansion carries the \emph{forward relations} $(z\Dq+\zeta^iL)X_a=C_1\kappa^{\mathrm{LR}}_aX_{a+1}$, $a=0,\dots,3$, of elements $X_a$ of $\Nring\lser z\otimes\bbQ(\zeta)$ (cf.\ \cref{prop:set-cyclicLR}~(iii)) to the relations $(z\Dq+\zeta^iL)\Cexp(X_a)=\kappa_a\Cexp(X_{a+1})$, $L=s^{-1}$; by \cref{cor:gau-cycsys}~(ii) the \emph{closing relation}, $a=4$, is then equivalent to the equation $\PFfive$, with $z\Dq$ replaced by $z\Dq+\zeta^iL$, for $\varpi_0\Cexp(X_0)$; for the rows it is obtained in \cref{thm:gau-R0}.
\item The cyclic system $z\Dq Y_a=\kappa_aY_{a+1}$ is $z\de_uY=AY$ with $A=\sum_ab_aE_{a,a+1}$, $b_a:=-\kappa_a/(1-u)$,
\[
b=\Bigl(\tc',\ -\frac1{u\varpi_0^2(1-u)^3\tc'^2},\ \tc',\ -\frac{\varpi_0}{1-u},\ -\frac{\varpi_0}{1-u}\Bigr),
\]
a self-dual cyclic system with one simple pole over $\bbQ\fps u$, with residue $r=-1$, product $\prod_ab_a=-1/(u(1-u)^5)$, and $b_a(0)=1,1,-1,-1$ for $a=0,2,3,4$.
\end{enumerate}
\end{prop}

\begin{prop}[Equal products and $v(u)$; {\Lr{prop:gau-equalprod}}, abridged]\genusstatus{AI-proposed, unchecked}
\label{prop:gau-equalprod}
Let $v\in u\bbQ\fps u$ with $v'(0)\neq0$. The model system $z\,\dd Y_b/\dd v=b^{\mathrm m}_bY_{b+1}$ pulled back along $v=v(u)$ is $z\de_uY=\tilde AY$ with $\tilde b=\bigl(v',\ -v'/v,\ v',\ -v',\ -v'\bigr)$ and $\prod_a\tilde b_a=-v'^5/v$.
\begin{enumerate}
\item If $v=u+O(u^2)$, $\tilde A$ is a self-dual cyclic system with one simple pole over $\bbQ\fps u$, with residue $\tilde r=-1$ and $\tilde b_a(0)=1,1,-1,-1$ for $a=0,2,3,4$.
\item For $v=u+O(u^2)$, the products of $\tilde A$ and of the quintic system of \cref{prop:gau-quintic} are equal if and only if $v'^5/v=1/(u(1-u)^5)$, if and only if $v=u\,\tau(u)^{-5/4}$ (\cref{def:set-tau}).
\end{enumerate}
\end{prop}

So the product of the coefficients, which does not depend on the frame, fixes the natural variable $\tau$: this is the origin of the natural form.

%-------------------------------------------------------------------------------
\subsection{Theorem U}
\label{sec:gau-U}

Let $A$, $\tilde A$ be cyclic systems with one simple pole over $\Bbbk\fps u$ satisfying (EP). Put $d^U_0:=1$, $d^U_{a+1}:=d^U_a\tilde b_a/b_a$ (consistently around the cycle, by (EP)), $\Theta_U:=\diag(d^U_0,\dots,d^U_4)$, so that $A\Theta_U=\Theta_U\tilde A$, and $\mathfrak d:=(\det\Theta_U/\det\Theta_U(0))^{-1/5}\in1+u\Bbbk\fps u$.

\begin{thm}[Theorem U; {\Lr{thm:gau-U}}]\genusstatus{AI-proposed, unchecked}
\label{thm:gau-U}
Let $A$, $\tilde A$ be cyclic systems with one simple pole over $\Bbbk\fps u$ with equal products.
\begin{enumerate}
\item[(a)] (Existence.) There is $P^\natural\in\Gau(A,\tilde A)$ with $P^\natural_0=\mathfrak d\,\Theta_U$, a diagonal matrix with unit entries.
\item[(b)] (Uniqueness.) For every $P^\natural\in\Gau(A,\tilde A)$ with $P^\natural_0=\mathfrak d\,\Theta_U$, the map $\Bbbk\fps z\to\Gau(A,\tilde A)$, $\cn\mapsto P^\natural\cn(z)$, is bijective. Consequently every solution $P$ of the gauge equation with coefficients in $\mathfrak{gl}_5(\Bbbk\fps u)$ is $P^\natural\cn(z)$ for a unique $\cn\in\Bbbk\fps z$; its $z^0$-term is $P_0=\cn(0)\,\mathfrak d\,\Theta_U$, a diagonal matrix, and $P$ is a regular gauge if and only if $\cn(0)\neq0$. If $P$ is one regular gauge, the regular gauges are exactly the $P\cn(z)$ with $\cn\in\Bbbk\fps z^\times$. For the pair $(\tilde A,\tilde A)$: $\Gau(\tilde A,\tilde A)=\Bbbk\fps z\cdot1$.
\item[(c)] (Graded gauges.) Graded regular gauges exist; if $P$ is one, the graded regular gauges are exactly the $P\cn(z)$ with $\cn\in\Bbbk\fps{z^5}^\times$.
\item[(d)] (The field.) If all $b_a$, $\tilde b_a$ lie in $\bbQ\lser u$, then $P^\natural$ can be chosen in $\mathfrak{gl}_5(\bbQ\fps u)\fps z$; for every field $\Bbbk\supseteq\bbQ$ the solutions with coefficients in $\mathfrak{gl}_5(\Bbbk\fps u)$ are then $P^\natural\Bbbk\fps z$, and $P^\natural\cn$ has rational coefficients if and only if $\cn\in\bbQ\fps z$.
\end{enumerate}
\end{thm}

\begin{sketch}[Sketch of the proof {\Lr{lem:gau-comm,lem:gau-lift,lem:gau-scal,lem:gau-lead,lem:gau-degree,prop:gau-graded}}]
Write $P=\Theta_U\check P$. The gauge equation becomes $[\tilde A,\check P_0]=0$ and $[\tilde A,\check P_{k+1}]=\de_u\check P_k+\Theta_U^{-1}(\de_u\Theta_U)\,\check P_k$ for $k\geq0$. Over $\Bbbk\lser u$ the commutant of $\tilde A$ is $\Bbbk\lser u[\mathcal Y]$, $\mathcal Y:=\tilde A^{-1}$, and $\mathfrak{gl}_5$ is the direct sum of the commutant and of the image of $\operatorname{ad}\tilde A$, with a projection $\operatorname{pr}$ onto the commutant. So the step $k+1$ is solvable if and only if $\operatorname{pr}$ of the right side vanishes, and $\check P_{k+1}$ is then determined up to an element $w=\sum_{j=0}^4w_j\mathcal Y^j$ of the commutant, to be chosen so that the next step is solvable. This condition splits into five scalar equations $\de_uw_j+\chi^U_jw_j=-\phi_j$ with $\chi^U_j=\frac j{5u}+(\text{a power series})$: for $j=1,\dots,4$ the pivots $n+\frac j5$ do not vanish, and there is exactly one power-series solution; for $j=0$ there is one free constant. This constructs $P^\natural$; the same analysis shows that the lowest coefficient of any solution is a constant multiple of $\mathfrak d\,\Theta_U$, and peeling off these constants gives (b). For (c), the gauge equation is homogeneous for the $\bbZ/5$-degree $k+b-a$ of $z^kE_{ab}$, so the part of degree $0$ of $P^\natural$ is a solution with the same $z^0$-term. The construction is rational, which gives (d).
\end{sketch}

% Abridged version, Section 3, part 2 of 2: the residue relation, the vanishing solutions, (G), the canonical gauge, symplecticity; the normalisation series.

%-------------------------------------------------------------------------------
\subsection{The vanishing solutions, the gauge identity and symplecticity}
\label{sec:gau-van}

From now on, unless said otherwise, $A$ is the quintic system of \cref{prop:gau-quintic} in a conifold frame and $\tilde A$ is the model pulled back along $v=v(u)$; by \cref{prop:gau-equalprod}, \cref{thm:gau-U} applies to them over $\bbQ$.

\begin{lem}[The residue relation (R1) and the diagonal; {\Lr{lem:gau-R1,lem:gau-diag}}, abridged]\genusstatus{AI-proposed, unchecked}
\label{lem:gau-R1}
Let $A,\tilde A$ be cyclic systems with one simple pole over $\Bbbk\fps u$, with residues $r,\tilde r$, and let $P\in\Gau(A,\tilde A)$; put $P(0,z):=\sum_kz^kP_k(0)$. Then $rE_{12}P(0,z)=\tilde r\,P(0,z)E_{12}$: column $1$ of $P(0,z)$ is $P_{11}(0,z)\mathbf e_1$, row $2$ is $P_{22}(0,z)\mathbf e_2^T$, and $rP_{22}(0,z)=\tilde rP_{11}(0,z)$. For the quintic system in a conifold frame and the pulled-back model, and a regular gauge $P$ whose diagonal $z^0$-term $P_0=\diag(p_{0,0},\dots,p_{0,4})$ has $p_{0,1}(0)=1$, one has $p_{0,a}(0)=1$ for all $a$ and $p_{0,0}(u)=\tau(u)^{-1/4}/\varpi_0$.
\end{lem}

\begin{defn}[Vanishing solutions; {\Lr{def:gau-van}}]
\label{def:gau-van}
\begin{enumerate}
\item For $m\in\bbZ$ the formal residue at $p=0$, taken coefficientwise in $v$, is
\[
\IC_m(v,z):=\Res_{p=0}p^m\re^{-\Wv(p)/z}\dd p=\sum_{i\geq0,\ m+4i+1\geq0}\frac{(-1/(4z))^i}{i!}\,\frac{(-v/z)^{m+4i+1}}{(m+4i+1)!}\in\bbQ[z^{-1}]\fps v .
\]
The \emph{vanishing solution of the model} is $\YC:=(\IC_0,-\IC_{-1},-\IC_3,\IC_2,\IC_1)$, i.e.\ $\YC_b=\varepsilon_b\IC_{\nu_b}$.
\item For a cyclic system $A=\sum_ab_aE_{a,a+1}$ with one simple pole over $\bbQ\fps u$ put $y^{(0)}:=-\mathbf e_1$, $y^{(n+1)}_a:=\int_0^ub_a\,y^{(n)}_{a+1}\,\dd u$ and $Y^{\mathrm C}_A:=\sum_{n\geq0}z^{-n}y^{(n)}$. For the quintic system, $\YCQ:=Y^{\mathrm C}_A$ is the \emph{vanishing solution of the quintic}.
\end{enumerate}
\end{defn}

\begin{lem}[{\Lr{lem:gau-vanmod,lem:gau-vanshape,lem:gau-uniq,lem:gau-vanq}}, abridged]\genusstatus{AI-proposed, unchecked}
\label{lem:gau-vanq}
\begin{enumerate}
\item The residues $\IC_m$ satisfy (Mod1) and (Mod2); hence $\YC$ solves the model system and $\YC(0,z)=-\mathbf e_1$. With $m:=\lfloor n/5\rfloor$, $\coef{z^{-n}}\YC_b=\varepsilon_b(-1)^nv^{n-m}/(4^m\,m!\,(n-m)!)$ if $n\equiv\nu_b+1\bmod5$, and $0$ otherwise; so $\YC(v(u),z)\in\bbQ[z^{-1}]\fps u^5$ has no positive power of $z$.
\item In \cref{def:gau-van}~(ii) every integrand is a power series; no logarithm arises. $Y^{\mathrm C}_A\in\bbQ[z^{-1}]\fps u^5$ solves $z\de_uY=AY$, and $Y^{\mathrm C}_A(0)=-\mathbf e_1$. For the quintic system in a conifold frame, $\coef{z^{-1}}\YCQ=-\tc\,\mathbf e_0$.
\item Let $B=u^{-1}B_{-1}+\sum_{j\geq0}u^jB_j$ with $B_j\in\mathfrak{gl}_5(\Bbbk)$ independent of $z$. If $Y\in\Bbbk\lser z\fps u^5$ satisfies $z\de_uY=BY$ and $Y(0)=0$, then $Y=0$.
\end{enumerate}
\end{lem}

\begin{thm}[The gauge identity (G); {\Lr{thm:gau-G}}, abridged]\genusstatus{AI-proposed, unchecked}
\label{thm:gau-G}
For every $P\in\Gau(A,\tilde A)$, in particular for every regular gauge from the pulled-back model to the quintic system,
\[
P(u,z)\,\YC(v(u),z)=P_{11}(0,z)\,\YCQ(u,z)\qquad\text{in }\bbQ\lser z\fps u^5 .
\]
\begin{enumerate}
\item[(a)] For a graded regular gauge $P$, $P_{11}(0,z)$ is a unit of $\bbQ\fps{z^5}$, $\Pc:=P\,P_{11}(0,z)^{-1}$ is a graded regular gauge, and $\Pc$ is the unique regular gauge $P'$ with $P'\,\YC(v(u),z)=\YCQ$.
\item[(b)] For a regular gauge $P$, $P\,\YC(v(u),z)$ has no positive power of $z$ if and only if $P_{11}(0,z)$ is a constant; the regular gauges with this property are the $\mu\Pc$, $\mu\in\bbQ^\times$.
\end{enumerate}
\end{thm}

\begin{proof}[Proof of the identity]
By \cref{lem:gau-vanq}, $\YC(v(u),z)=Y^{\mathrm C}_{\tilde A}$. The product $P\,Y^{\mathrm C}_{\tilde A}$ lies in $\bbQ\lser z\fps u^5$ and solves $z\de_uY=AY$, by the gauge equation; its value at $u=0$ is $P(0,z)(-\mathbf e_1)=-P_{11}(0,z)\mathbf e_1$ by (R1). $P_{11}(0,z)Y^{\mathrm C}_A$ is a solution with the same value, and \cref{lem:gau-vanq}~(iii) gives the identity. Parts (a) and (b) follow from \cref{thm:gau-U}~(b), (c) and the coefficient of $u^0$.
\end{proof}

\begin{defn}[The canonical gauge; {\Lr{def:gau-can}}]
\label{def:gau-can}
In a conifold frame, the \emph{canonical gauge} $\Pc$ is the regular gauge from the pulled-back model to the quintic system with $(\Pc)_{11}(0,z)=1$. It exists and is unique. Equivalently, $\Pc$ is the unique regular gauge with $\Pc\YC(v(u))=\YCQ$. It is graded, rational, colour covariant (that is, $\zeta^{i(a-b)}(\Pc)_{ab}(u,\zeta^{-i}z)=(\Pc)_{ab}(u,z)$ for all $i$, which is equivalent to being graded), and $p_{0,a}(0)=1$ for its diagonal $z^0$-term. We write $\Pc=\sum_kz^k\sum_bp_{k,b}(u)E_{b+k,b}$.
\end{defn}

\begin{prop}[Symplecticity; {\Lr{prop:gau-symp}}]\genusstatus{AI-proposed, unchecked}
\label{prop:gau-symp}
Let $A,\tilde A$ be self-dual cyclic systems with one simple pole over $\Bbbk\fps u$ with (EP) and $r=\tilde r$, and let $P$ be a regular gauge from $\tilde A$ to $A$, $\cn(z):=P_{11}(0,z)$. Then $P(u,-z)^T\eta\,P(u,z)=\cn(z)\cn(-z)\,\eta$. In particular, for the two systems of this subsection, $\Pc$ is symplectic, $\Pc(u,-z)^T\eta\Pc(u,z)=\eta$; the symplectic regular gauges are the $\Pc\cn(z)$ with $\cn(z)\cn(-z)=1$, so symplecticity constrains only the even part of $\log(\cn/\cn(0))$. Among the symplectic regular gauges, only (G) singles out $\Pc$.
\end{prop}

\begin{proof}[Proof of the identity]
By self-duality, $T:=\eta^{-1}P(u,-z)^T\eta P(u,z)$ satisfies $zT'=\tilde AT-T\tilde A$, so $T=t(z)\,1$ by \cref{thm:gau-U}~(b) for the pair $(\tilde A,\tilde A)$; by (R1), $r=\tilde r$ and $\eta_{a1}=5\delta_{a2}$, the entry $(2,1)$ of $\eta T$ at $u=0$ gives $5t(z)=5\cn(-z)\cn(z)$.
\end{proof}

%-------------------------------------------------------------------------------
\subsection{The normalisation series}
\label{sec:gau-R0}

We compare $\Pc$ applied to the model columns with the rows of \cref{def:set-rows}. Here $\theta_u:=u\,\dd/\dd u$, $\theta_{\zh}:=\zh\,\dd/\dd\zh$, and a series $\Phi(\zh,u)=\sum_n\zh^n\mathrm{ph}_n(u)$ is identified with $\sum_nz^ns^{-4n}\mathrm{ph}_n$ through $\zh=z/s^4$.

\begin{lem}[HPF; {\Lr{lem:gau-HPF,lem:gau-poly}}, abridged]\genusstatus{AI-proposed, unchecked}
\label{lem:gau-HPF}
For $e\in\bbZ$ let $\mathscr O_e:=1+\zh(1-u)\bigl(\tfrac45\theta_{\zh}-\tfrac e5-\theta_u\bigr)$, and put $\Phi^{[0]}:=\Phi$, $\Phi^{[j]}:=\mathscr O_{-j}\Phi^{[j-1]}$, $j=1,\dots,5$.
\begin{enumerate}
\item $L\Phi$ satisfies $\PFfive^{(0)}(L\Phi)=0$, where $\PFfive^{(0)}$ is $\PFfive$ with $z\Dq$ replaced by $z\Dq+L$ (that is, $\re^{u_0/z}L\Phi$ solves $\PFfive$ formally), if and only if
\begin{equation}
\label{eq:gau-HPF}
\Phi^{[5]}-\Phi-(1-u)\Bigl[3\zh\Phi^{[4]}+\tfrac{17}5u\zh^2\Phi^{[3]}+\tfrac95u^2\zh^3\Phi^{[2]}+\tfrac{274}{625}u^3\zh^4\Phi^{[1]}+\tfrac{24}{625}u^4\zh^5\Phi\Bigr]=0 .
\end{equation}
\item The coefficient of $\zh^{n+1}$ of the left side of \eqref{eq:gau-HPF} is $(1-u)\bigl[(4n-5\theta_u)\mathrm{ph}_n+r_n\bigr]$, where $r_n$ depends only on $\mathrm{ph}_0,\dots,\mathrm{ph}_{n-1}$. So, if every $\mathrm{ph}_n$ lies in $\bbQ\fps u$, the coefficient $\coef{u^k}\mathrm{ph}_n$ is fixed by $\mathrm{ph}_0,\dots,\mathrm{ph}_{n-1}$ for $5k\neq4n$, while for $n=5m$ the coefficient $\fc_m:=\coef{u^{4m}}\mathrm{ph}_{5m}$ is free.
\item Let $\Phi$ be a solution of \eqref{eq:gau-HPF} with $\mathrm{ph}_0=1$ and every $\mathrm{ph}_n\in\bbQ\fps u$. Then every $\mathrm{ph}_n$ is a polynomial in $u$, and the solutions with $\mathrm{ph}_0=1$ and coefficients in $\bbQ\fps u$ are exactly the $\Phi\,\mathsf k(z)$, $\mathsf k\in1+z^5\bbQ\fps{z^5}$ (with $z^n=s^{4n}\zh^n$).
\end{enumerate}
\end{lem}

\begin{lem}[Formal WKB uniqueness and colour covariance; {\Lr{lem:gau-wkb,lem:gau-colmod,lem:gau-colrows}}, abridged]\genusstatus{AI-proposed, unchecked}
\label{lem:gau-colrows}
Let $(\cR,\Dq,L)$ be $(\Bbbk\fps q,\,q\,\dd/\dd q,\,(1-5^5q)^{-1/5})$ or $(\Bbbk\lser s,\Dq,s^{-1})$, $\Bbbk\supseteq\bbQ(\zeta)$ a field. A \emph{colour-$i$ WKB amplitude} of $\PFfive$ over $\cR$ is a series $\cA=\sum_{n\geq0}\cA_nz^n\in\cR\fps z$ with $\PFfive^{(i)}\cA=0$, where $\PFfive^{(i)}$ is $\PFfive$ with $z\Dq$ replaced by $z\Dq+\zeta^iL$.
\begin{enumerate}
\item If $\cA,\cA'$ are colour-$i$ WKB amplitudes with $\cA_0=\cA'_0=L$, then $\cA'=\cA\,\mathsf k(z)$ for a unique $\mathsf k\in1+z\Bbbk\fps z$; and $\cA(z)\mapsto\cA(\zeta^{-i}z)$ is a bijection from colour-$0$ to colour-$i$ WKB amplitudes that preserves $\cA_0$.
\item The model columns $\Qr^{\mathrm m,(i)}_b(v,z):=\varsigma^{\nu_b-1}\Rhat^{\mathrm m,(i)}_b(z/\varsigma^4)$ satisfy $\Qr^{\mathrm m,(i)}_b(u,z)=\zeta^{i(b+1)}\Qr^{\mathrm m,(0)}_b(u,\zeta^{-i}z)$ at $v=v(u)$.
\item For every $i$, as identities of formal series in $z$ with coefficients in $\bbQ(\zeta)\fps q$, $\Phi_i(\zh,u)=\Phi_0(\zeta^{-i}\zh,u)$ and $\Qr^{(i)}_a(z)=\zeta^{i(a+1)}\Qr^{(0)}_a(\zeta^{-i}z)$.
\end{enumerate}
\end{lem}

In (iii), $L\Phi_i$ and $L\Phi_0(\zeta^{-i}\zh,u)$ are colour-$i$ WKB amplitudes, so they differ by a constant series $\mathsf k(z)$, and (QRR) at $q=0$ gives $\mathsf k=1$; this is the only place in this section where the value of (QRR) enters.

\begin{thm}[The normalisation series $\Dn$; {\Lr{thm:gau-R0,cor:gau-colours}}, abridged]\genusstatus{AI-proposed, unchecked}
\label{thm:gau-R0}
Let $(\alpha,\beta)$ be a conifold frame, $\Pc$ the canonical gauge, and $\Qr^{\mathrm m,(i)}$ the model columns at $v=v(u)$.
\begin{enumerate}
\item For every $i$, $\Pc\Qr^{\mathrm m,(i)}$ is a formal WKB solution of the quintic system with exponential $\re^{u_i/z}$: its rows $X_a\in\bbQ(\zeta)\lser s\fps z$ satisfy $(z\Dq+\zeta^iL)X_a=\kappa_aX_{a+1}$ for $a\in\bbZ/5$, with $L=s^{-1}$.
\item $\Phican(\zh,u):=\varpi_0\,s\,\bigl(\Pc\Qr^{\mathrm m,(0)}\bigr)_0\big|_{z=s^4\zh}$ (row $0$, colour $0$) solves \eqref{eq:gau-HPF}, with coefficients in $\bbQ\fps u$ and $\mathrm{ph}^{\mathrm{can}}_0=1$. Hence every $\mathrm{ph}^{\mathrm{can}}_n$ is a polynomial in $u$ and $A^{\mathrm{can}}_n:=L^{4n+1}\mathrm{ph}^{\mathrm{can}}_n(L^{-5})$ lies in $\bbQ[L,L^{-1}]$.
\item Let $A^{\mathrm{act}}:=C_0\Qr^{(0)}_0=\sum_nA^{\mathrm{act}}_nz^n$ be the flat amplitude of the actual colour-$0$ row $0$. In $\bbQ(\zeta)\fps q$, with $L=(1-5^5q)^{-1/5}$, both $A^{\mathrm{can}}:=\sum_nA^{\mathrm{can}}_nz^n$ and $A^{\mathrm{act}}$ are colour-$0$ WKB amplitudes of $\PFfive$ with coefficient $L$ at $z^0$. Hence $A^{\mathrm{can}}=A^{\mathrm{act}}\,\Dn_0(z)$ for a unique $\Dn_0\in1+z\bbQ(\zeta)\fps z$, and every $A^{\mathrm{act}}_n$ is the image of a unique element of $\bbQ(\zeta)[L,L^{-1}]$.
\item In hatted form, $\Phiact(\zh,u)\,\Dn_0(s^4\zh)=\Phican(\zh,u)$, $\Phiact:=\Phi_0$. For every $i$, $\Dn_i(z):=\Dn_0(\zeta^{-i}z)$ is the unique series in $\bbQ(\zeta)\fps z$ with $\Qr^{(i)}\,\Dn_i(z)=\Pc\,\Qr^{\mathrm m,(i)}(v(u))$ in $\bbQ(\zeta)\lser s\fps z$, for the conifold expansions of the rows. Hence, with $\Dn:=\diag(\Dn_0,\dots,\Dn_4)$, $\Qr\,\Dn=\Pc\,\Qr^{\mathrm m}(v(u))$, and $\Dn_0=1$ implies $\Dn=1$.
\end{enumerate}
\end{thm}

\begin{sketch}
(i) The model columns solve the pulled-back model with the exponentials $\re^{U_i/z}$, and $U_i(v(u))=-\frac54\zeta^iu^{4/5}/\tau=u_i-u_i^{\mathrm c}$ (\cref{def:set-tau}): this is where the choice of $\tau$ enters. The gauge equation then gives the quintic system for $\Pc\Qr^{\mathrm m,(i)}$. (ii) By (i) and \cref{cor:gau-cycsys}~(ii), $L\Phican$ solves $\PFfive^{(0)}$; in the hatted variables the entries of $\Pc$ carry non-negative integer powers of $u$ only, and $\mathrm{ph}^{\mathrm{can}}_0=\varpi_0\tau^{1/4}p_{0,0}=1$ by \cref{lem:gau-R1}; then \cref{lem:gau-HPF}~(iii) applies. (iii) $\bbQ[L,L^{-1}]$ embeds in $\bbQ\fps q$ and in $\bbQ\lser s$ compatibly with $\Dq$; $A^{\mathrm{act}}$ is a WKB amplitude by \cref{cor:set-cyclicPF5}~(b), and \cref{lem:gau-colrows}~(i) gives $\Dn_0$. (iv) Row $0$ is (iii); the other rows follow from the forward relations, $\Dn_0$ being constant, and the other colours from the colour covariance of $\Pc$, of the model columns and of the rows.
\end{sketch}

\begin{cor}[The rows lie in the differential ring; the window $\Fg{g}$; {\Lr{cor:gau-rowsNN,cor:gau-colours,def:gau-Fg}}, abridged]\genusstatus{AI-proposed, unchecked}
\label{cor:gau-rowsNN}
For all $i,a,n$ the row coefficient $\coef{z^n}\Qr^{(i)}_a$ is the evaluation of an explicit $\logq$-free element of $\Nring\otimes\bbQ(\zeta)$: for colour $0$, row $0$ is $A^{\mathrm{act}}_n(L)\,c_0^{-1}$; the rows $a=1,\dots,4$ follow from the forward relations of \cref{prop:set-cyclicLR}~(iii); and colour $i$ is $\zeta^{i(a+1-n)}$ times colour $0$. These representatives are used throughout; in this sense the rows lie in $\Nring\otimes\bbQ(\zeta)$ and have conifold expansions, and the expanded rows of every colour form a WKB solution of the quintic system with exponential $\re^{u_i/z}$. For $g\geq2$, the \emph{window} $\Fg{g}$ is the statement $\Dn_0\equiv1\pmod{z^{3g-2}}$ in the conifold frame considered; it implies $\Dn\equiv1\pmod{z^{3g-2}}$, and it is the hypothesis of the comparison theorem (\cref{thm:flh-main}).
\label{def:gau-Fg}
\end{cor}


% Abridged version, Section 4, part 1 of 2: the vanishing tower, the slice, dressings, Lemmas E, T, F. Continued in part 2 (Theorem R).
\section{Transport and the reduction to the model}
\label{sec:tra}

This section states the transport lemmas: for every symplectic dressing $P$ of the model that annihilates the vanishing period, in particular for the canonical gauge $\Pc$, the edges lie in $\sKo\otimes\sKo$ and the leaf point lies in $\sK$. Theorem~R then reduces the regularity of the graph terms to a statement about the model. It follows \Lr{sec:tra,sec:red}. In this section $q$ denotes an element of a loop space, as in Givental's formalism.

%-------------------------------------------------------------------------------
\subsection{The vanishing tower and the slice}
\label{sec:tra-tower}

\emph{Settings} \Lr{conv:tra-forms}. The model is used with two rings of coefficients. In the \emph{graded setting} the coefficients lie in $\bbQ[v]$; the \emph{weight} of $v^mz^kH^b$ is $5m+4k+\nu_b+1$, $V_w$ is the $\bbQ$-span of the $v^mz^kH^b$ ($m,k\geq0$) of weight $w$, and $\cH^{\mathrm{gr}}_+:=\bigoplus_wV_w=\bbH\otimes\bbQ[v][z]$. In the \emph{$u$-adic setting} $v=v(u)$ is a non-zero element of $u\bbQ\fps u$, with $v'(0)\neq0$ from \cref{def:tra-dressing} on, and
\[
\cH^u:=\bbH\otimes\bbQ\lser z\fps u,\qquad\cH^u_+:=\bbH\otimes\bbQ\fps{u,z},\qquad\cH^u_-:=\bbH\otimes z^{-1}\bbQ[z^{-1}]\fps u ;
\]
the form $\Omega_{\mathrm s}(f,g):=\Res_z\,\eta(f(-z),g(z))\,\dd z$ is antisymmetric, and $\cH^u_\pm$ are isotropic and complementary. An insertion $\psi^kH^b$ is identified with the vector $z^kH^b$, with flat components $q_{(k,b)}:=\coef{z^k}q^b$. The \emph{vector form} of $P$ is $P^{\mathrm{vec}}:=\eta^{-1}P\eta$, and $J^{\mathrm C}:=\eta^{-1}\YC$ is the vector form of the vanishing solution. Infinite sums are sums of \emph{summable families}, finite modulo every power of $u$ in each coefficient \Lr{lem:tra-sum,lem:tra-conv}.

\begin{defn}[The vanishing tower and the slice; {\Lr{def:tra-tower,def:tra-slice}}]
\label{def:tra-tower}
For $q$ in $\cH^{\mathrm{gr}}_+$ or in $\cH^u_+$ and $i\geq0$ put $\Pi_i(q):=-\coef{z^i}\,\langle\YC(-z),q(z)\rangle$, where $\langle Y,n\rangle:=\sum_aY_an^a$; $\Pi_0$ is the \emph{vanishing covector}. The \emph{slice} is
\[
\sK:=\{q:\Pi_i(q)=0\ \text{for all }i\geq1\},\qquad\sKo:=\sK\cap\ker\Pi_0 .
\]
\end{defn}

By \cref{lem:gau-vanq}~(i), $\Pi_i(z^kH^b)=\Pi_0(z^{k-i}H^b)$ for $k\geq i$, $\Pi_i(z^kH^b)=0$ for $k<i$, and, with $r=\lfloor m/5\rfloor$,
\[
\Pi_0(z^mH^b)=-\varepsilon_b\,\frac{v^{m-r}}{4^r\,r!\,(m-r)!}\quad\text{if }b+m\equiv1\pmod5,\qquad\Pi_0(z^mH^b)=0\quad\text{otherwise};
\]
for example $\Pi_0(H)=1$, $\Pi_0(H^a)=0$ for $a\neq1$, and $\Pi_0(z1)=-v$. The tower consists of symplectic pairings with the vanishing solution: $\Omega_{\mathrm s}(z^{-k}J^{\mathrm C},q)=(-1)^{k+1}\Pi_{k-1}(q)$ for $k\geq1$.

\begin{lem}[Frames; {\Lr{lem:tra-frames,cor:tra-slice,lem:tra-twoslot}}, abridged]\genusstatus{AI-proposed, unchecked}
\label{lem:tra-frames}
Let $A$ be $\bbQ[v]$ or $\bbQ\fps u$, and $\cH^A_+$ the corresponding space $\cH^{\mathrm{gr}}_+$ or $\cH^u_+$.
\begin{enumerate}
\item $\cH^A_+=\sKo\oplus A\{z\}H$, where $A\{z\}$ is $\bbQ[v][z]$, respectively $\bbQ\fps{u,z}$; and $\sK=\sKo\oplus AH$. Moreover $-z1\in\sK$, and $\Pi_0(-z1+n+\mu H)=v+\mu$ for $n\in\sKo$ and $\mu\in A$.
\item For every $(k,a)$ with $a\neq1$ there is exactly one element of $\sKo$ of the form $z^kH^a+\chi(z)H$; it is the \emph{frame}
\[
\fr(k,a)=z^kH^a+\sum_{j<k}f_jz^jH,\qquad f_j=c_{k,a,j}\,v^{(4(k-j)+\nu_a+1)/5},\quad c_{k,a,j}\in\bbQ,
\]
where $c_{k,a,j}=0$ unless the exponent is an integer, which is then $\geq1$. For example $\fr(1,0)=z1+vH=:\se$, the Euler--dilaton vector. The frames of the $u$-adic setting are those of the graded setting at $v=v(u)$.
\item Every $n\in\sKo$ is the sum of the summable family $(n_{(k,a)}\fr(k,a))_{k\geq0,\,a\neq1}$, and $n\mapsto(n_{(k,a)})$ is injective.
\item For a bivector series $E\in\bbH\otimes\bbH\otimes\bbQ\fps{\psi,\psi'}\fps u$ the following are equivalent: $(\Pi_i\otimes1)E=0$ and $(1\otimes\Pi_i)E=0$ for all $i\geq0$; $E$ is the sum of a summable family $(E^{\alpha\beta}\,\fr_\alpha(\psi)\otimes\fr_\beta(\psi'))$ with $E^{\alpha\beta}\in\bbQ\fps u$. In this case we say that $E$ lies in $\sKo\otimes\sKo$.
\end{enumerate}
\end{lem}



%-------------------------------------------------------------------------------
\subsection{Dressings, edges and the leaf point}
\label{sec:tra-dress}

\begin{defn}[Dressings; {\Lr{def:tra-dressing,def:tra-edge}}, abridged]
\label{def:tra-dressing}
Let $v=v(u)\in u\bbQ\fps u$ with $v'(0)\neq0$. A \emph{dressing of the model at $v(u)$} is a matrix series $P=\sum_{k\geq0}z^kP_k(u)$ with $P_k\in\mathfrak{gl}_5(\bbQ\fps u)$, graded, i.e.\ $P_k=\sum_bp_{k,b}E_{b+k,b}$, and with $P_0$ invertible; it acts on lowered components. We consider the two conditions
\begin{description}
\item[(Sp)] $P(-z)^T\eta P(z)=\eta$;
\item[(A)] $Y':=P\,\YC(v(u))$ has no positive power of $z$.
\end{description}
Put $T_P:=-\coef{z^{-1}}Y'_0\in\bbQ\fps u$, the transported vanishing coordinate; $d:=(P_0)_{33}=p_{0,3}$, the dilaton factor, a unit of $\bbQ\fps u$; and $R:=P_0^{-1}P=1+O(z)$, with vector form $R^{\mathrm{vec}}$. For a dressing with (Sp), the \emph{edge} is $E_P:=E_{R^{\mathrm{vec}}}$ (\cref{def:set-graphsum}), the \emph{leaf point} is $q_T:=(P^{\mathrm{vec}})^{-1}(-z1)\in\cH^u_+$, and $\tilde q_T:=d\,q_T=-z1+dT_{\geq2}$, where $dT_{\geq2}=\ttr_{R^{\mathrm{vec}}}(z)\in z^2\cH^u_+$ is the translation of \cite[Def.~2.13]{PPZ15}.
\end{defn}



\begin{lem}[Lemma E; {\Lr{lem:tra-E}}]\genusstatus{AI-proposed, unchecked}
\label{lem:tra-E}
Let $R=1+O(z)$ lie in $\mathfrak{gl}_5(\bbQ\fps{u,z})$ with $R(-z)^T\eta R(z)=\eta$, and let $f\in\cH^u_-$, both in vector form. Then $Rf\in\cH^u$, the contraction of the covector $\Omega_{\mathrm s}(f,\cdot)$ into the first slot of $E_R$ converges $u$-adically, and
\[
(\Omega_{\mathrm s}(f,\cdot)\otimes1)\,E_R(\psi')=-R(\psi')^{-1}\,(Rf)_{\geq0}(\psi')\qquad\text{in }\bbH\otimes\bbQ\fps{u,\psi'} ,
\]
where $(\cdot)_{\geq0}$ is the part with non-negative powers of $z$. Hence $E_R$ annihilates $\Omega_{\mathrm s}(f,\cdot)$ in the first slot, equivalently (by symmetry) in the second, if and only if $Rf\in\cH^u_-$.
\end{lem}

\begin{sketch}
The contraction is the residue in $\psi$ of $\eta f(-\psi)$ against $E_R=(\psi+\psi')^{-1}\bigl(\eta^{-1}-R^{-1}(\psi)\eta^{-1}R^{-1}(\psi')^T\bigr)$, with $(\psi+\psi')^{-1}$ expanded in $\psi'/\psi$. The first term has no residue; in the second, symplecticity turns $R^{-1}(\psi)^T\eta f(-\psi)$ into $\eta\,(Rf)(-\psi)$, and the residue against the geometric series picks out its non-negative part.
\end{sketch}

\begin{cor}[Edges in $\sKo\otimes\sKo$; {\Lr{cor:tra-edges}}]\genusstatus{AI-proposed, unchecked}
\label{cor:tra-edges}
Let $P$ be a dressing with (Sp). Then (A) holds if and only if $(\Pi_0\otimes1)E_P=0$, and if and only if $(\Pi_i\otimes1)E_P=(1\otimes\Pi_i)E_P=0$ for all $i\geq0$. If (A) holds, $E_P$ lies in $\sKo\otimes\sKo$, in the sense of \cref{lem:tra-frames}~(iv), with coefficients $E_P^{\alpha\beta}\in\bbQ\fps u$.
\end{cor}

Indeed $\Pi_{k-1}=\pm\Omega_{\mathrm s}(z^{-k}J^{\mathrm C},\cdot)$, and by Lemma~E the edge annihilates it if and only if $R^{\mathrm{vec}}z^{-k}J^{\mathrm C}$, that is $z^{-k}Y'$ up to an invertible matrix independent of $z$, has no non-negative power of $z$.

\begin{lem}[Lemma T: the leaf point; Lemma F: finiteness; {\Lr{lem:tra-T,lem:tra-F}}, abridged]\genusstatus{AI-proposed, unchecked}
\label{lem:tra-T}
Let $P$ be a dressing with (Sp) and (A).
\begin{enumerate}
\item $\Pi_0(q_T)=T_P$ and $\Pi_i(q_T)=0$ for $i\geq1$; so $q_T\in\sK$.
\item $\tilde q_T=d\,q_T=-z1+dT_{\geq2}$ lies in $-z1+\sK=\sK$, with $\Pi_0(\tilde q_T)=d\,T_P=:v+\delta_P$. Moreover $dT_{\geq2}=y^\circ+\delta_PH$ with $y^\circ\in\sKo$, whose frame coordinates $y^\circ_{(k,a)}$ vanish for $k\leq1$.
\item $\delta_P\in u^2\bbQ\fps u$; more precisely $T_P-p_{0,0}v\in u^2\bbQ\fps u$ and $d\,p_{0,0}=1$.
\item For a dressing $P$ with (Sp) and a leg-free stable graph $\Gamma$, the contraction along $\Gamma$ of the edges $E_P$ with the vertex tensors $d^{2g_v-2+n_v}\dd^{n_v}\overline F_{g_v}(\tilde q_T)$ is a finite sum. If $d(0)=1$, each vertex tensor equals the Taylor expansion $\dd^{n_v}\overline F_{g_v}(q_T)$ at the base point $-z1$ in the direction $q_T+z1$.
\end{enumerate}
\end{lem}

Part (i) follows from the adjunction $\Omega_{\mathrm s}(f,(P^{\mathrm{vec}})^{-1}g)=\Omega_{\mathrm s}(P^{\mathrm{vec}}f,g)$ with $f=z^{-i-1}J^{\mathrm C}$, $g=-z1$, from (A) and from the grading.

\begin{prop}[The canonical gauge as a dressing; {\Lr{prop:tra-can}}]\genusstatus{AI-proposed, unchecked}
\label{prop:tra-can}
Let $(\alpha,\beta)\in\bbQ^2$ give a conifold frame and let $v(u)=u\tau(u)^{-5/4}$. The canonical gauge $\Pc$ is a dressing of the model at $v(u)$ with (Sp) and (A), with $Y'=\YCQ$, $T_{\Pc}=\tc$ and $d(0)=1$. Hence \cref{cor:tra-edges,lem:tra-T} apply to $\Pc$: its edges lie in $\sKo\otimes\sKo$, and $\tilde q_T\in-z1+\sK$ with $\Pi_0(\tilde q_T)=d\,\tc=v+\delta_{\Pc}$, $\delta_{\Pc}\in u^2\bbQ\fps u$.
\end{prop}

Here (Sp) is \cref{prop:gau-symp}, and (A) and $T_{\Pc}=\tc$ come from the gauge identity $\Pc\YC(v(u))=\YCQ$ (\cref{thm:gau-G}, \cref{lem:gau-vanq}~(ii)). So the gauge identity has two effects: it places the edges and the leaf of the canonical gauge in the slice, and it identifies the vanishing coordinate of the leaf with the flat coordinate $\tc$ of the frame.

% Abridged version, Section 4, part 2 of 2: the flat sum, the Model Lemma on K, Theorem R.

%-------------------------------------------------------------------------------
\subsection{The flat sum, the Model Lemma and Theorem R}
\label{sec:red}

\begin{defn}[The flat sum; {\Lr{def:red-flatsum}}]
\label{def:red-flatsum}
Let $P$ be a dressing with (Sp), with edge $E_P$, leaf point $q_T$ and $\tilde q_T=d\,q_T$. For a leg-free stable graph $\Gamma$ of genus $g\geq2$, $\mathrm{Cont}_\Gamma$ is $|\Aut\Gamma|^{-1}$ times the contraction along $\Gamma$, as in \cref{def:set-graphsum}, of $E_P$ on every edge with the vertex tensor $d^{\,2g_v-2+n_v}\,\dd^{n_v}\overline F_{g_v}(\tilde q_T)$ at every vertex, of genus $g_v$ and valence $n_v$; here $\overline F_h$ is the dilaton-shifted ancestor potential of the model. The \emph{flat sum} is $F^P_g:=\sum_\Gamma\mathrm{Cont}_\Gamma$. For the vertex graph,
\[
\mathrm{Cont}_{\mathrm{vertex}}=d^{\,2g-2}\,\overline F_g(\tilde q_T)=d^{\,2g-2}\sum_{m\geq0}\frac1{m!}\langle(dT_{\geq2})^m\rangle_g .
\]
For $P=\Pc$ we write $F_\Gamma:=\mathrm{Cont}_\Gamma$, the \emph{flat term} of $\Gamma$.
\end{defn}



\begin{defn}[The Model Lemma on $\sK$; {\Lr{def:red-ML}}]
\label{def:red-ML}
For $h\geq2$ let $\Fm_h(v):=\langle\ \rangle_h(v)$ be the genus-$h$ free energy of the model, a rational multiple of $v^{2-2h}$; put $\Fm_1(v):=-\frac1{12}\log v$ and $\Fm_0(v):=-\frac52v^2\log v$, of which only the derivatives of order $\geq1$, resp.\ $\geq3$, are used. For $h\geq0$, $\ML{h}$ is the following statement.
\begin{enumerate}
\item For $n\geq1$ with $2h-2+n>0$ and homogeneous $X_1,\dots,X_n\in\sK$ (graded setting, coefficients in $\bbQ[v]$),
\[
\langle X_1\cdots X_n\rangle_h(v)-(\Fm_h)^{(n)}(v)\prod_{i=1}^n\Pi_0(X_i)\ \in\ \bbQ[v] .
\]
\item For $h\geq2$, $\Fm_h=c_hv^{2-2h}$, with $c_h$ as in \cref{def:set-N}.
\end{enumerate}
In short, $\overline F_h|_{-z+\sK}=\Fm_h\circ\Pi_0+\mathrm{Reg}_h$ with a regular remainder $\mathrm{Reg}_h$, an identity of the Taylor coefficients at the base point of order $n>2-2h$.
\end{defn}



\begin{thm}[Theorem R; {\Lr{thm:red-main,cor:red-can}}]\genusstatus{AI-proposed, unchecked}
\label{thm:red-main}
Fix $g\geq2$ and a dressing $P$ of the model at $v(u)$, $v(0)=0$, $v'(0)\neq0$, with (Sp) and (A).
\begin{description}
\item[(a)] If $\ML{h}$ holds for $0\leq h\leq g-1$, then $\mathrm{Cont}_\Gamma\in\bbQ\fps u$ for every leg-free stable graph $\Gamma$ of genus $g$ with at least one edge.
\item[(b)] If $\ML{g}$ holds, then $\mathrm{Cont}_{\mathrm{vertex}}\in c_g\,T_P^{2-2g}+\bbQ\fps u$.
\end{description}
The gauge equation is not used. In particular, in a conifold frame: if $\ML{h}$ holds for $0\leq h\leq g-1$, every \emph{edge graph} (a graph with at least one edge) of genus $g$ of the flat sum of $\Pc.\Omega^{\mathrm{mod}}_{v(u)}$ is regular, $F_\Gamma\in\bbQ\fps u$; if moreover $\ML{g}$ holds, the vertex graph satisfies $F_{\mathrm{vertex}}\in c_g\,\tc^{2-2g}+\bbQ\fps u$.
\end{thm}

\begin{sketch}
\emph{Edges and leaves.} By \cref{cor:tra-edges}, $E_P=\sum E_P^{\alpha\beta}\fr_\alpha\otimes\fr_\beta$ with coefficients in $\bbQ\fps u$, so every half-edge insertion is a frame, a homogeneous element of $\sKo$ evaluated at $v(u)$. By \cref{lem:tra-T}, the leaf is $dT_{\geq2}=y^\circ+\delta_PH$ with $y^\circ\in\sKo$ and $\delta_P\in u^2\bbQ\fps u$.

\emph{Edge graphs.} Every vertex of a graph of genus $g$ with an edge carries at least one half-edge and has genus $h\leq g-1$ \Lr{lem:red-genus}. Its factor, with the frames $t_{\alpha_i}:=\fr_{\alpha_i}$ at its half-edges, is $d^{\,2h-2+n}\sum_m\frac1{m!}\langle t_{\alpha_1}\cdots t_{\alpha_n}(dT_{\geq2})^m\rangle_h$, a finite sum; expanded multilinearly, with $y^\circ$ expanded in frames, every term is a product of elements of $\bbQ\fps u$ and of a model correlator of frames and of $H$ that contains the frame $t_{\alpha_1}$. As $\Pi_0(t_{\alpha_1})=0$, the singular term of $\ML{h}$~(i) is absent, and the correlator lies in $\bbQ[v]$, hence in $\bbQ\fps u$ at $v=v(u)$. The contraction along $\Gamma$ is the sum of a summable family of such products, so it lies in $\bbQ\fps u$.

\emph{The vertex graph.} Here $\mathrm{Cont}_{\mathrm{vertex}}=d^{\,2g-2}\,\overline F_g(\tilde q_T)$ with $\tilde q_T\in-z1+\sK$ and $\Pi_0(\tilde q_T)=v+\delta_P$. Under $\ML{g}$ the Taylor expansion at the base point resums, $u$-adically because $\delta_P\in u^2\bbQ\fps u$ and $v'(0)\neq0$, to $\overline F_g(\tilde q_T)\in c_g\,(v+\delta_P)^{2-2g}+\bbQ\fps u$. Finally $v+\delta_P=dT_P$, so $d^{\,2g-2}c_g(dT_P)^{2-2g}=c_gT_P^{2-2g}$: the dilaton factor cancels. The last statement is the case $P=\Pc$, $T_{\Pc}=\tc$ (\cref{prop:tra-can}).
\end{sketch}




% Abridged version, Section 5, part 1 of 2: the model as topological recursion, the leg dictionary, the lattice theorem.
% Continued in part 2 (L-regularity, the Gap Theorem, Theorem V, the Model Lemma).
\section{The model: topological recursion on the model curve}
\label{sec:mod}

Topological recursion is applied here, and only here, to the curve $x=\Wv(p)$, $y=p$ on $\bbP^1$, $v\neq0$, and to its rescalings; nothing about the quintic enters. This section follows \Lr{sec:mod}; the inputs (Q) are those of \cref{sec:int-inputs}.

%-------------------------------------------------------------------------------
\subsection{The model as topological recursion}
\label{sec:mod-tr}

\emph{Conventions} \Lr{conv:mod-kernel,def:mod-curve,def:mod-pair,def:mod-xi,def:mod-ancvec}. For meromorphic functions $x,y$ on $\bbP^1$ such that $\dd x$ has only simple zeros, the \emph{branch points} $\varsigma_b$, and $\dd y\neq0$ there, put $\omega_{0,1}:=y\,\dd x$ and $\omega_{0,2}(p_1,p_2):=\dd p_1\,\dd p_2/(p_1-p_2)^2$; the correlators $\omega_{h,n}$, $2h-2+n>0$, are defined by the recursion of \cite[(2-26)--(2-29)]{Dunin-Barkowski:2015caa}, a sum of residues at the branch points \Lr{conv:mod-kernel}; for $h\geq2$ the free energy is defined by $(2-2h)\FTR_h:=\sum_b\Res_{\varsigma_b}\Phi_{y\dd x}\,\omega_{h,1}$, $\dd\Phi_{y\dd x}=y\,\dd x$. On the \emph{model curve} $x=\Wv(p)$, $y=p$ one has $\dd x=(p^5-v)\,\dd p/p^2$; the branch points are $\varsigma_b:=\zeta^{-b}\varsigma$, and $p=0$, $p=\infty$ are the poles of $x$. The \emph{leg class} is $\cL:=\bbQ[v][p]+\bbQ(v)[x]$; for legs $\phi_1,\dots,\phi_n\in\cL$ the \emph{residue pairing} is
\[
\langle\phi_1\cdots\phi_n\rangle^{\mathrm{TR}}_h:=\sum_{b_1,\dots,b_n}\Res_{p_1=\varsigma_{b_1}}\cdots\Res_{p_n=\varsigma_{b_n}}\prod_{l=1}^n\phi_l(p_l)\,\omega_{h,n}(p_1,\dots,p_n),
\]
and the \emph{log leg} $-\log p$ is paired by integration along a ray $\mathfrak r$ from $0$ to $\infty$ avoiding the branch points, $\langle\cdots,-\log p,\cdots\rangle^{\mathrm{TR}}_h:=\int_{p'\in\mathfrak r}\omega_{h,n}(\dots,p',\dots)$. Let $D_x:=\dd/\dd x$, and for $a\in\bbZ/5$ and $d\geq0$ let
\[
\xi_a:=\varepsilon'_a\,v^{\alpha_a}\frac{p^{m_a}}{p^5-v},\quad(m_a,\varepsilon'_a,\alpha_a)=(4,-1,0),\ (0,1,1),\ (1,1,0),\ (2,-1,0),\ (3,-1,0),\qquad\xi_{a,d}:=D_x^d\xi_a ;
\]
for a leg $\phi$ put $q_\phi(z):=\sum_{a,d}\langle\phi,\dd\xi_{a,d}\rangle H^az^d$, where $\langle\phi,w\rangle:=\sum_b\Res_{\varsigma_b}\phi\,w$.

\begin{thm}[The model as topological recursion; {\Lr{thm:mod-tr}}]\genusstatus{AI-proposed, unchecked}
\label{thm:mod-tr}
Let $v\neq0$, $n\geq1$ and $2h-2+n>0$. The correlators of the model curve are
\[
\omega_{h,n}(p_1,\dots,p_n)=(-5)^{h-1}\sum_{a_1,\dots,a_n}\sum_{d_1,\dots,d_n\geq0}\Big\langle\prod_{i=1}^n\psi^{d_i}H^{a_i}\Big\rangle_h(v)\,\prod_{i=1}^n\dd\xi_{a_i,d_i}(p_i),
\]
a finite sum, the brackets being the ancestor correlators of $\Omega^{\mathrm{mod}}_v$. Equivalently, for legs $\phi_1,\dots,\phi_n$ in $\cL$ or equal to $-\log p$, $\langle q_{\phi_1}\cdots q_{\phi_n}\rangle_h=(-5)^{1-h}\langle\phi_1\cdots\phi_n\rangle^{\mathrm{TR}}_h$.
\end{thm}

\begin{sketch}[Sketch of the proof {\Lr{lem:mod-laplace,lem:mod-scale,prop:mod-Rtr}}]
The quoted theorem \cite[Theorem~4.1]{DuninBarkowski:2012bw}, as recalled in \cite[\S2.6]{Dunin-Barkowski:2015caa}, says: for a compact curve with simple branch points, if the matrix $R$ defined by the Laplace transforms at the branch points of the forms $\dd\xi^i$ attached to them is symplectic, satisfies \cite[(2-32)]{Dunin-Barkowski:2015caa}, and the expansion of $\omega_{0,2}$ at the branch points is the edge of $R$, then the correlators are the graph sum of $R$ on a topological field theory, paired with the forms $\dd(D_x^d\xi^i)$. The three hypotheses are verified for the curves $x=\kappa\Wv$, $y=p$: the Laplace transforms reduce to the series of the forms $p^m\dd p$ at each branch point, which satisfy (Mod1) and (Mod2), hence the T-system of \cref{prop:set-modrows}; by its uniqueness, $R_{\mathrm{TR}}(\psi)=\Mo\bigl(\psi/(\kappa\varsigma^4)\bigr)$, which is symplectic (\cref{def:set-M0}). The theorem is applied at $\kappa=4$ and transferred to $\kappa=1$ by scaling; passing to the normalised canonical basis of $\TFT^{\mathrm m}_v$ produces the factor $(-5)^{h-1}$ (the metric of topological recursion is $-\frac15$ times that of the model), and passing to the flat basis by $\Psi^{\mathrm m}_v$ turns the forms $\xi^i$ into the $\xi_a$.
\end{sketch}

%-------------------------------------------------------------------------------
\subsection{The leg dictionary and the lattice theorem}
\label{sec:mod-legs}

Let $\phi^b_0:=p,\ \log p,\ v/p,\ p^3,\ p^2$ for $b=0,\dots,4$, and $\phi^b_{k+1}:=I(\phi^b_k)$, where $I(F)$ is the termwise antiderivative of $-F\,\dd x$ without constant term; so $D_x\phi^b_{k+1}=-\phi^b_k$, and $\phi^b_k$ is homogeneous of weight $\nu_b+1+4k$ for the weights $p:1$, $v:5$, $x:4$.

\begin{prop}[The leg dictionary; {\Lr{prop:mod-LD}}, abridged]\genusstatus{AI-proposed, unchecked}
\label{prop:mod-LD}
Let $\varrho:=(-1,1,-1,-3,-2)$. For all $a,b\in\bbZ/5$ and $k,d\geq0$, $\langle\phi^b_k,\dd\xi_{a,d}\rangle=-\varrho_b\,\delta_{ab}\,\delta_{dk}$. Hence $q_{\phi^b_k}=-\varrho_b\,z^kH^b$: the leg of the insertion $v^m\psi^kH^b$ is $-\varrho_b^{-1}v^m\phi^b_k$. In particular the unit $H^0$ has the leg $y=p$; $\Phi_{y\dd x}=p^5/5-v\log p$ has $q_{\Phi_{y\dd x}}=-z\,H^0$; $-\log p$ has $q_{-\log p}=H$; the Euler--dilaton vector $\se=z\,H^0+vH$ has the leg $-p^5/5$; and $q_{f(x)}=0$ for $f\in\bbQ(v)[x]$.
\end{prop}

Functions of $x$ pair to zero with every $\omega_{h,n}$, because the sum of $\omega_{h,n}$ over the five sheets of $x$ vanishes \Lr{lem:mod-xfun}. With the dilaton equation the dictionary gives $\Fm_h=(-5)^{1-h}\FTR_h$ for $h\geq2$ \Lr{cor:mod-trF}.

The \emph{leg map} is the $\bbQ$-linear map $\Lambda(v^m\psi^kH^b):=\varrho_b^{-1}v^m\phi^b_k$ on the free $\bbQ[v]$-module on the $\psi^kH^b$, graded by the weight spaces $V_w$; it takes values in $\mathsf M:=\bbQ[p,\bar p]\oplus\log p\cdot\bbQ[v,x]$, where $\bar p:=v/p$ and $\bbQ[p,\bar p]=\bbQ[v][p]+\bbQ[v][x]$ is the ring of the Laurent polynomials whose terms $v^mp^j$ have $j\geq-m$. The element of $\bbQ[v,x]$ multiplying $\log p$ is the \emph{$\log p$-coefficient}.

\begin{prop}[Log coefficients of the legs; {\Lr{prop:mod-dii}}]\genusstatus{AI-proposed, unchecked}
\label{prop:mod-dii}
For all $b$ and $k$, the $\log p$-coefficient of $\phi^b_k$ is $\varrho_b\sum_{i=0}^k\Pi_i(\psi^kH^b)\,(-x)^i/i!$. Hence, for $X\in V_w$, the $\log p$-coefficient of $\Lambda(X)$ is $\sum_{i\geq0}\Pi_i(X)(-x)^i/i!$, and $\Lambda(X)$ is log-free if and only if $X\in\sKo$.
\end{prop}

\begin{sketch}
The $\log p$-coefficients of the $\phi^b_k$ are generated by the residues $\Res_{p=0}\Wv^n\,\dd\phi^b_0/n!$, and these are $\varrho_b\,\Pi_0(\psi^nH^b)$, since $\IC_m=\Res_{p=0}p^m\re^{-\Wv/z}\dd p$ \Lr{lem:mod-genfn,lem:mod-resid}.
\end{sketch}

Thus the vanishing tower, defined by the residue at the pole $p=0$ of the model, measures exactly the logarithms of the legs.

\begin{thm}[The lattice theorem; {\Lr{thm:mod-lattice,cor:mod-Ksplit}}, abridged]\genusstatus{AI-proposed, unchecked}
\label{thm:mod-lattice}
For every $w\geq0$:
\begin{enumerate}
\item[(a)] for $X\in\sKo\cap V_w$, $\Lambda(X)$ lies in $\bbQ[p,\bar p]=\bbQ[v][p]+\bbQ[v][x]$, exactly and not only modulo functions of $x$;
\item[(b)] $\Lambda\colon\sKo\cap V_w\to(\cL/\bbQ(v)[x])_w$ is bijective;
\item[(c)] $\sK\cap V_w=(\sKo\cap V_w)\oplus\bbQ v^{w/5}H$ if $5\mid w$, and $\sK\cap V_w=\sKo\cap V_w$ otherwise; for $X\in\sK\cap V_w$, $\Lambda(X)=\Lambda(X-\Pi_0(X)H)+\Pi_0(X)\log p$ with $\Lambda(X-\Pi_0(X)H)\in\bbQ[p,\bar p]$.
\end{enumerate}

\end{thm}

Part (a) is \cref{prop:mod-dii}; the bijectivity is proved modulo $v$ and lifted by a graded Nakayama argument \Lr{lem:mod-modv}. So an insertion from $\sK$ has a leg that is a polynomial in $p$ with coefficients in $\bbQ[v]$, plus a function of $x$, which pairs to zero, plus $\Pi_0(X)$ times the log leg.

% Abridged version, Section 5, part 2 of 2: L-regularity, the Gap Theorem, Theorem V, the Model Lemma on K.

%-------------------------------------------------------------------------------
\subsection{L-regularity, the Gap Theorem and Theorem V}
\label{sec:mod-Lreg}

The next two statements use the closed formula \cite[Theorem~7.1, (7.9), (7.10)]{ABDKS25} for the differentials of topological recursion on a rational curve $(x=x(z),y=z)$; its assumptions hold for $x=\Wv$, $y=p$, $v\neq0$.

\begin{thm}[L-regularity; {\Lr{lem:mod-abdks,thm:mod-Lreg}}]\genusstatus{AI-proposed, unchecked}
\label{thm:mod-Lreg}
For $2h-2+n>0$ there are an integer $K\geq0$ and $N_{h,n}\in\bbQ[v][p_1,\dots,p_n]$ with $\omega_{h,n}=N_{h,n}(p;v)\prod_i\dd p_i\big/\prod_i(p_i^5-v)^K$. Consequently, for legs $\phi_1,\dots,\phi_n\in\bbQ[v][p]$ the residue pairing $\langle\phi_1\cdots\phi_n\rangle^{\mathrm{TR}}_h$ lies in $\bbQ[v]$; for legs in $\cL=\bbQ[v][p]+\bbQ(v)[x]$ it is the restriction of an element of $\bbQ[v]$, and the components in $\bbQ(v)[x]$ pair to zero.
\end{thm}

\begin{sketch}
The formula of \cite{ABDKS25} writes $\omega_{h,n}$ as a finite sum of terms in $\bbQ[v][p_i,p_i^{-1},(p_i^5-v)^{-1},(p_i-p_j)^{-1}]\prod_i\dd p_i$; there are no poles at $p_i=0$ or on the diagonals, so these denominators cancel in $\bbQ[v][p]$, without negative powers of $v$. For polynomial legs the sum of the residues at the branch points is minus the residue at $p=\infty$, and $\Res_{p=\infty}p^a\,\dd p/(p^5-v)^K$ is an integer multiple of a non-negative power of $v$.
\end{sketch}

\begin{thm}[The Gap Theorem; {\Lr{prop:mod-period,thm:mod-gap}}]\genusstatus{AI-proposed, unchecked}
\label{thm:mod-gap}
On the model curve, for every $h\geq1$, $\omega_{h,1}=\dd\Phi_h$ with $\Phi_h$ rational and holomorphic at $p=0$ and $p=\infty$, and $\int_\infty^0\omega_{h,1}=\frac{B_{2h}}{2h}\,v^{1-2h}$ along any path avoiding the branch points. For every $h\geq2$,
\[
\FTR_h=\frac{B_{2h}}{2h(2h-2)}\,v^{2-2h},\qquad\text{hence}\qquad\Fm_h=c_hv^{2-2h},\quad c_h=(-5)^{1-h}\frac{B_{2h}}{2h(2h-2)}=\frac{(-1)^{h-1}B_{2h}}{2h(2h-2)\,5^{h-1}} .
\]
\end{thm}

\begin{sketch}
By \cite[(7.10)]{ABDKS25}, $\omega_{h,1}$ is the differential of an explicit rational function $\Phi_h$, a finite sum of terms $D_x^{r-1}(a_r/x')$. It vanishes at $p=\infty$ by a count of degrees. At $p=0$ only the leading parts of the factors contribute, those of the pole $v/p$ of $\Wv$, and the value is $v^{1-2h}$ times a universal constant, which is $B_{2h}/(2h)$. The free energy follows from $(2-2h)\FTR_h=\langle\Phi_{y\dd x}\rangle^{\mathrm{TR}}_h$ with $\Phi_{y\dd x}=p^5/5-v\log p$, whose polynomial part pairs to zero.
\end{sketch}

For logarithmic legs the tool is the variational formula of topological recursion \cite[Theorem~5.1, (5-14)]{Eynard:2007kz}: for points $J$ off the branch points, $n\geq1$ and $h+n>1$,
\[
\partial_v|_x\,\omega_{h,n}(J)=\int_{q\in\mathfrak r}\omega_{h,n+1}(J,q),
\]
where $\partial_v|_x$ is the derivative in $v$ at fixed $x$ \Lr{thm:mod-var}. With $\partial_v|_x=\partial_v|_p-\mathrm{Lie}_{\cX}$, $\cX:=\frac p{p^5-v}\frac{\dd}{\dd p}$, and an integration by parts, it becomes a recursion on the number of log legs.

\begin{thm}[Theorem V; {\Lr{thm:mod-V}}]\genusstatus{AI-proposed, unchecked}
\label{thm:mod-V}
Let $h\geq0$, $a,k\geq0$, $n=a+k\geq1$, $2h-2+n>0$, and $\phi_1,\dots,\phi_a\in\bbQ[v][p]$.
\begin{enumerate}
\item If $a\geq1$, then $\langle\phi_1\cdots\phi_a,(-\log p)^k\rangle^{\mathrm{TR}}_h\in\bbQ[v]$.
\item If $a=0$, then $\langle(-\log p)^k\rangle^{\mathrm{TR}}_h=(\FTR_h)^{(k)}(v)$, with $\FTR_h$ of \cref{thm:mod-gap} for $h\geq2$, $\FTR_1:=-\frac1{12}\log v$ (here $k\geq1$) and $\FTR_0:=\frac12v^2\log v$ (here $k\geq3$).
\item The recursion (V1) holds:
\[
\begin{aligned}
\langle\phi_1\cdots\phi_a,(-\log p)^{k+1}\rangle^{\mathrm{TR}}_h={}&\partial_v\langle\phi_1\cdots\phi_a,(-\log p)^k\rangle^{\mathrm{TR}}_h-\sum_{i=1}^a\langle\cdots,\partial_v\phi_i,\cdots,(-\log p)^k\rangle^{\mathrm{TR}}_h\\
&+\sum_{i=1}^a\langle\cdots,(\cX\dd\phi_i)_+,\cdots,(-\log p)^k\rangle^{\mathrm{TR}}_h ,
\end{aligned}
\]
where $(\cX\dd\phi)_+\in\bbQ[v][p]$ is the quotient in the division of $p\,\phi'(p)$ by $p^5-v$.
\end{enumerate}
\end{thm}

\begin{sketch}
(iii) is the variational formula, differentiated under the pairings. For (i), induct on $k$: the right side of (V1) consists of a $v$-derivative of, and pairings with, $a$ polynomial legs and $k$ log legs, and the induction starts with \cref{thm:mod-Lreg} or, in genus $0$, with the explicit $\omega_{0,3}$. For (ii), (V1) with $a=0$ reads $\langle(-\log p)^{k+1}\rangle^{\mathrm{TR}}_h=\partial_v\langle(-\log p)^k\rangle^{\mathrm{TR}}_h$, so it suffices to treat the first admissible $k$: for $h\geq2$ this is \cref{thm:mod-gap}, and for $h=0,1$ it follows from $\omega_{0,3}$ and from a local formula for $\omega_{1,1}$.
\end{sketch}

%-------------------------------------------------------------------------------
\subsection{The Model Lemma on \texorpdfstring{$\sK$}{K}}
\label{sec:mod-MLK}

\begin{thm}[The Model Lemma on $\sK$; {\Lr{thm:mod-MLK}}]\genusstatus{AI-proposed, unchecked}
\label{thm:mod-MLK}
$\ML{h}$ holds for every $h\geq0$.
\end{thm}

\begin{proof}[Proof, given the results of this section]
Part (ii) of $\ML{h}$ is \cref{thm:mod-gap}. For part (i), let $X_i\in\sK\cap V_{w_i}$, $i=1,\dots,n$. By \cref{thm:mod-lattice}~(c), $X_i=Y_i+\Pi_0(X_i)H$ with $Y_i\in\sKo\cap V_{w_i}$ and $\Pi_0(X_i)\in\bbQ v^{w_i/5}$, a monomial with a non-negative exponent. By \cref{prop:mod-LD}, $X_i=q_{\ell_i}$ for the leg $\ell_i:=-\Lambda(Y_i)+\Pi_0(X_i)\,(-\log p)$, and by \cref{thm:mod-lattice}~(a), $-\Lambda(Y_i)=\phi_i+f_i(x)$ with $\phi_i\in\bbQ[v][p]$ and $f_i\in\bbQ[v][x]$. By \cref{thm:mod-tr}, multilinearity and the vanishing of the pairings with functions of $x$,
\[
\langle X_1\cdots X_n\rangle_h=(-5)^{1-h}\sum_{S\subseteq\{1,\dots,n\}}\prod_{i\notin S}\Pi_0(X_i)\,\big\langle(\phi_i)_{i\in S},(-\log p)^{n-|S|}\big\rangle^{\mathrm{TR}}_h .
\]
Every term with $S\neq\emptyset$ lies in $\bbQ[v]$, by \cref{thm:mod-Lreg} if $S=\{1,\dots,n\}$ and by \cref{thm:mod-V}~(i) otherwise. By \cref{thm:mod-V}~(ii) and $\Fm_h=(-5)^{1-h}\FTR_h$ (for every $h\geq0$, with $\FTR_0$, $\FTR_1$ as in \cref{thm:mod-V}), the term $S=\emptyset$ is $(\Fm_h)^{(n)}(v)\prod_i\Pi_0(X_i)$.
\end{proof}




% Abridged version, Section 6: the comparison of the flat and the outer graph sums.
\section{Flat and outer graph sums}
\label{sec:flh}

This section compares the outer terms $T_\Gamma$ of the natural form (\cref{def:set-outer}) with the flat terms $F_\Gamma$ (\cref{def:red-flatsum}).
The argument is graph by graph; it uses no total potential and no topological recursion. This section follows \Lr{sec:flh}. A conifold frame is fixed, $M:=\Mo(\tau(u)\zh)$ is read, exactly in $u$, as a power series in $z$ over $\bbQ(\zeta)\lser s$, and $R':=\Qr M^{-1}$.

Both sides are instances of one definition \Lr{def:flh-outer}. Let $\cA$ be $\bbQ\fps u$ or $\bbQ(\zeta)\lser s$, let $(\mathbb V,\eta,1)$ and $(\mathbb V',\eta',1')$ be free $\cA$-modules of the same finite rank with perfect pairings and distinguished elements, $\Omega$ a cohomological field theory with flat unit $1$ on $\mathbb V$, and $O(z)=\sum_{k\geq0}O_kz^k$, $O_k\colon\mathbb V\to\mathbb V'$, a symplectic series, $O(-z)^*O(z)=1$, with
\[
\text{(unit)}\qquad O_0^{-1}1'=\varkappa^{-1}1\quad\text{for a unit }\varkappa\text{ of }\cA .
\]
Its edge and leaf point are $E_O(\psi,\psi'):=\bigl(\eta^{-1}-(O^{-1}(\psi)\otimes O^{-1}(\psi'))\eta'^{-1}\bigr)/(\psi+\psi')$ and $q_O(z):=O^{-1}(z)(-z1')$, and $\mathrm{Cont}_\Gamma[E_O,q_O;\Omega]$ is formed as in \cref{def:set-graphsum}, with the potentials of $\Omega$ shifted by its own unit. For a unit $\varkappa$, $\varkappa.\Omega$ is the theory $(\varkappa.\Omega)_{g,n}:=\varkappa^{2g-2+n}\Omega_{g,n}$, with flat unit $\varkappa^{-1}1$. The canonical gauge $\Pc$ satisfies (unit) with $\varkappa=d$, and its edge and leaf point are those of \cref{def:tra-dressing}; the outer matrix $R'$ satisfies (unit) with $\varkappa=1$, a condition on its constant term only; $R'$ is not assumed to be symplectic, and $T_\Gamma$ is formed with the edge of \cref{def:set-outer}.

\begin{lem}[Four lemmas; {\Lr{lem:flh-frame,lem:flh-dil,lem:flh-comm,lem:flh-hat}}]\genusstatus{AI-proposed, unchecked}
\label{lem:flh-hat}
\begin{enumerate}
\item (Frame change.) Let $A\colon(\widetilde{\mathbb V},\tilde\eta)\to(\mathbb V,\eta)$ be an isometry independent of $z$ and $\widetilde\Omega$ a cohomological field theory with flat unit on $\widetilde{\mathbb V}$. For every leg-free stable graph $\Gamma$ and all $E$ and $q$, $\mathrm{Cont}_\Gamma[E,q;A_*\widetilde\Omega]=\mathrm{Cont}_\Gamma\bigl[(A^{-1}\otimes A^{-1})E,\;A^{-1}q;\;\widetilde\Omega\bigr]$.
\item (TFT dilaton invariance.) Let $\varkappa\in1+u\bbQ\fps u$ and $q\in\mathbb V\fps z$ with $q+z\varkappa^{-1}1\in z^2\mathbb V\fps z$. Then for every leg-free stable graph $\Gamma$, $\mathrm{Cont}_\Gamma[E,q;\varkappa.\Omega]=\mathrm{Cont}_\Gamma[E,q;\Omega]$.
\item (The $R$-action commutes with the rescaling.) For a unit $\varkappa$ of $\cA$ and an $\End(\mathbb V)$-valued series $R(z)=1+O(z)$ with $R(-z)^*R(z)=1$, $R.(\varkappa.\Omega)=\varkappa.(R.\Omega)$.
\item (Hatting.) Let $g\geq2$. For every leg-free stable graph $\Gamma$ of genus $g$, the hatted term and the outer term of \cref{def:set-outer} satisfy $G_\Gamma=s^{10(g-1)}\,T_\Gamma=u^{2g-2}\,T_\Gamma$.
\end{enumerate}
\end{lem}

Part (ii) is the dilaton equation in the form of \cref{lem:set-dilaton}~(iii), and (iv) is a count of powers of $s$.

\begin{lem}[The inputs of the comparison; {\Lr{lem:flh-inputs}}]\genusstatus{AI-proposed, unchecked}
\label{lem:flh-inputs}
Let $\Qr^{\mathrm m}$ be the matrix of the model columns at $v(u)$, $\Psi^{\mathrm m}:=\Qr^{\mathrm m}|_{z=0}$, $\Psi_{\mathrm Q}:=\Qr|_{z=0}$, $P_0:=\Pc|_{z=0}$ and $d:=(P_0)_{33}$, and identify the colour bases of the formal quintic and of the model through the index $i$.
\begin{enumerate}
\item $\Qr^{\mathrm m}=\Psi^{\mathrm m}M$ and $\Omega^{\mathrm{mod}}_{v(u)}=(\Psi^{\mathrm m})_*\bigl(M.\TFT^{\mathrm m}_{v(u)}\bigr)$; $\Psi^{\mathrm m}$ is an isometry, independent of $z$, from the colour space to $(\bbH,\eta)$, and it maps the unit of $\TFT^{\mathrm m}_{v(u)}$ to $H^0$.
\item $\Qr\Dn=\Pc\Qr^{\mathrm m}$ with $\Dn=\diag(\Dn_i)$ and $\Dn_i(z)=\Dn_0(\zeta^{-i}z)$. Under $\Fg{g}$, $R'\equiv\Pc\Psi^{\mathrm m}\bmod z^{3g-2}$.
\item $\Pc$ is symplectic; so $E_{\Pc}$ is a power series and $\Pc\Psi^{\mathrm m}$ is symplectic. Moreover
\[
R'(-z)^T\eta^{-1}R'(z)=M(z)\bigl(\Dn(z)\Dn(-z)\bigr)^{-1}M(z)^{-1};
\]
hence $R'$ is symplectic, $R'(-z)^T\eta^{-1}R'(z)=1$, if and only if $\Dn_0(z)\Dn_0(-z)=1$, and under $\Fg{g}$ it is symplectic modulo $z^{3g-2}$.
\item $\Pc$ satisfies (unit) with $\varkappa=d$, where $d\in\bbQ\fps u$ and $d(0)=1$; $R'$ satisfies (unit) with $\varkappa=1$; $\Psi_{\mathrm Q}=P_0\Psi^{\mathrm m}$; and $\TFT=d.\TFT^{\mathrm m}_{v(u)}$ on the identified colour spaces.
\end{enumerate}
\end{lem}

Here (ii) is \cref{thm:gau-R0}~(iv), since $R'=\Pc\Psi^{\mathrm m}\,M\Dn^{-1}M^{-1}$, and (iii) is \cref{prop:gau-symp}. The symplecticity of $R$ at large radius is not used: this is why (iii) carries the factor $\Dn$.

\begin{thm}[The comparison theorem; {\Lr{thm:flh-main}}]\genusstatus{AI-proposed, unchecked}
\label{thm:flh-main}
Let $g\geq2$ and let $\Gamma$ be a leg-free stable graph of genus $g$. Let $F_\Gamma$ be its term in the flat sum of $\Pc.\Omega^{\mathrm{mod}}_{v(u)}$, $T_\Gamma$ its outer term and $G_\Gamma$ its hatted term (with $M=\Mo(\tau(u)\zh)$ exactly in $u$). Assume $\Fg{g}$. Then
\[
T_\Gamma=F_\Gamma,\qquad\text{hence}\qquad G_\Gamma=u^{2g-2}F_\Gamma .
\]
\end{thm}

\begin{sketch}
\emph{Only $R'$ through $z^{3g-3}$ enters.} The classes of $M.\TFT$ have non-negative degree and every leaf insertion has degree at least $2$, so the dimension of the moduli spaces bounds the degrees: a monomial $\psi^a\psi'^b$ of an edge occurs only with $a+b\leq3g-4$, and a leaf insertion only up to degree $3g-2$. These involve $R'^{-1}$ only through $z^{3g-3}$, so $T_\Gamma$ depends on $R'$ only modulo $z^{3g-2}$. Under $\Fg{g}$ we may therefore replace $R'$ by the symplectic series $R'':=\Pc\Psi^{\mathrm m}$ (\cref{lem:flh-inputs}~(ii), (iii)), whose edge and leaf point are those of $\Pc$ transported by the isometry $\Psi^{\mathrm m}$.

\emph{The chain.} By \cref{lem:flh-inputs}~(iv), (i) and \cref{lem:flh-hat}~(iii), (ii), (i),
\begin{align*}
T_\Gamma&=\mathrm{Cont}_\Gamma[E_{R''},q_{R''};M.(d.\TFT^{\mathrm m})]=\mathrm{Cont}_\Gamma[E_{R''},q_{R''};d.(M.\TFT^{\mathrm m})]=\mathrm{Cont}_\Gamma[E_{R''},q_{R''};M.\TFT^{\mathrm m}]\\
&=\mathrm{Cont}_\Gamma[E_{\Pc},q_{\Pc};(\Psi^{\mathrm m})_*(M.\TFT^{\mathrm m})]=\mathrm{Cont}_\Gamma[E_{\Pc},q_{\Pc};\Omega^{\mathrm{mod}}_{v(u)}] .
\end{align*}
The last contraction is formed with the leaf point $q_T$ and with the Taylor expansions of the model potentials at $-z1$; since $d(0)=1$ these converge and equal the vertex tensors $d^{2g_v-2+n_v}\dd^{n_v}\overline F_{g_v}(dq_T)$ of \cref{def:red-flatsum} (\cref{lem:tra-T}~(iv)). So it is $F_\Gamma$.
\end{sketch}

At genus $g$ the outer terms involve the normalisation series only through $z^{3g-3}$: $\Fg{g}$ is all of $\Dn_0=1$ that genus $g$ sees. If $\Dn=\cn(z)1$ were a scalar with $\cn(z)\cn(-z)=1$, the proof would identify $T_\Gamma$ with the term of the flat sum of the symplectic gauge $\Pc\cn(z)^{-1}$, a different graph decomposition, which satisfies (A) only if $\cn=1$ (\cref{thm:gau-G}~(b)) \Lr{rmk:flh-fail,rmk:mai-stored}.

% Abridged version, Section 7, part 1 of 2: formal stationary phase, the full-Stirling series, Theorem P(a). Continued in part 2.
\section{The connection problem}
\label{sec:con}

The canonical gauge is normalised at the conifold point $u=0$, the rows of the formal quintic at $q=0$; the normalisation series $\Dn_0$ of \cref{thm:gau-R0} measures the difference. This section reduces $\Dn_0=1$ to one identity at $q=0$. The bridge is the saddle series $c(z;L)$ of a Mellin--Barnes summand, with hatted form $\PhiSP$, which is defined for all $q$ at once: Theorem~P shows that $\PhiSP$ is $\Phican$, and the glue gives $\Dn_0(z)=\PhiSP(z,1)/h(z)$. We work in colour~$0$. This section follows \Lr{sec:con}.

%-------------------------------------------------------------------------------
\subsection{Formal stationary phase}
\label{sec:con-fsp}

Stationary phase is used as an algebraic operation on Taylor series \Lr{sec:fsp}; nothing is evaluated at $z\neq0$. Let $A$ be a commutative $\bbQ$-algebra and $c\in A^n$. A \emph{phase} $F\in A\fps{x-c}$ has the \emph{non-degenerate critical point} $c$ if $\nabla F(c)=0$ and $\mathsf H:=\Hess F(c)\in GL_n(A)$. Write $F(c+t)=F(c)+\frac12t^{\mathrm T}\mathsf Ht+F_{\geq3}(t)$ and $C:=-\mathsf H^{-1}$. For an amplitude $a\in A\fps{x-c}\lser z$ the \emph{normalised stationary phase} is
\[
\SPn_c[F;a]:=E_{zC}\bigl[\exp\bigl(F_{\geq3}(t)/z\bigr)\,a(c+t)\bigr]\in A\lser z,\qquad E_{zC}[t^\alpha z^j]:=\langle t^\alpha\rangle_C\,z^{j+|\alpha|/2},
\]
where $\langle t^\alpha\rangle_C$ is the Gaussian moment with covariance $C$, the sum over the perfect matchings of the products of the entries of $C$ \Lr{def:fsp-ring,def:fsp-wick}. If $a=\sum_{m\geq0}a_mz^m$, its constant term is $a_0(c)$; if this is a unit, $\SPh_c[F;a]:=\SPn_c[F;a]/a_0(c)\in1+zA\fps z$.  The coefficients of $\SPn_c[F;a]$ are universal polynomials over $\bbQ$ in the entries of $C$ and in finitely many Taylor coefficients of $F$, of order at least $3$, and of $a$ \Lr{lem:fsp-univ}; in particular they do not see constants added to $F$. By \Lr{sec:fsp}, the stationary phase vanishes on exact forms $\dd(\re^{F/z}\iota)$; it is invariant under formal changes of variables and under translations by multiples of $z$; it can be iterated variable by variable (a formal Fubini theorem); it commutes with derivatives in a parameter. The \emph{formal Stirling form} is the expression $\Gst(Y):=\exp\bigl((Y-\frac12)\log Y-Y+\frac12\log2\pi+\sum_{k\geq1}\bk_kY^{1-2k}\bigr)$, with $\Gst(1+w):=w\,\Gst(w)$; it satisfies the functional equation, and it is the stationary phase of Euler's integral: for a unit $a$ and an integer $M$,
\[
\SPn_a\bigl[a\log x-x;\ x^{M-1}\bigr]=a^{M-1}\,\mathrm{Poch}_M(z/a)\,\exp\Bigl(\sum_{k\geq1}\bk_k\,(z/a)^{2k-1}\Bigr)
\]
\Lr{lem:fsp-stirfe,lem:fsp-stirling}, where $\mathrm{Poch}_M(w)$ is the polynomial or rational function with $\mathrm{Poch}_{m+1}(w)=(1+mw)\mathrm{Poch}_m(w)$ and $\mathrm{Poch}_0=\mathrm{Poch}_1=1$.

%-------------------------------------------------------------------------------
\subsection{The full-Stirling saddle series and \texorpdfstring{$\PhiSP$}{the full-Stirling WKB function}}
\label{sec:con-form}

\begin{defn}[The full-Stirling saddle series; {\Lr{def:gam-c,prop:gam-rational}}]
\label{def:gam-c}
Let $\Cstar:=\bbC\setminus\{0,1,\zeta,\zeta^2,\zeta^3,\zeta^4\}$ and $q(L):=(1-L^{-5})/5^5$. For $L\in\Cstar$ fix an open disc $\bbD_L$ around $L$ containing neither $0$ nor any $\zeta^j$, holomorphic branches of $\log(5\sigma)$ and of $\log(\zeta^j-\sigma)$ on $\bbD_L$, and put $\lambda:=-5\log(5L)+\sum_{j=0}^4\log(\zeta^j-L)$, so that $\re^\lambda=-q(L)$. On $\bbD_L$ define
\begin{align*}
\mathsf S_0(\sigma)&:=\lambda\sigma+5\sigma\log(5\sigma)-5\sigma+\sum_{j=0}^4\bigl[(\zeta^j-\sigma)\log(\zeta^j-\sigma)-(\zeta^j-\sigma)\bigr],\\
\mathsf S_1(\sigma)&:=\tfrac12\log(5\sigma)-\tfrac12\sum_{j=0}^4\log(\zeta^j-\sigma),\qquad\mathsf S_{2k}(\sigma):=\bk_k\Bigl[(5\sigma)^{1-2k}+\sum_{j=0}^4(\zeta^j-\sigma)^{1-2k}\Bigr]\quad(k\geq1),
\end{align*}
and $A_{\mathrm{St}}(\sigma,z):=\exp\bigl(\mathsf S_1(\sigma)+\sum_{k\geq1}z^{2k-1}\mathsf S_{2k}(\sigma)\bigr)$.
Then $\mathsf S_0'(L)=0$ and $\mathsf S_0''=5/(\sigma(1-\sigma^5))$, which does not vanish at $L$. The \emph{full-Stirling saddle series} is
\[
c(z;L):=\SPh_L[\mathsf S_0;A_{\mathrm{St}}](z)=\sum_{n\geq0}c_n(L)z^n,\qquad c_0(L)=1 .
\]
It does not depend on the disc nor on the branches of the logarithms, and every $c_n$ is the value at $L$ of an element of $\AL=\bbQ[L,L^{-1},(L^5-1)^{-1}]$.
\end{defn}

Up to a factor $(2\pi)^3z^2$, $\exp(\mathsf S_0/z)A_{\mathrm{St}}$ is the product $\re^{\lambda\sigma/z}\,\Gst(1+5\sigma/z)\prod_j\Gst((\zeta^j-\sigma)/z)$: this is the summand of \cref{thm:int-gamma} with all six Gamma factors replaced by $\Gst$. 

By the reflection formula for $\Gst$, the same series is the saddle series of the colour-$0$ summand $q^{\sigma/z}\,\Gamma(1+5\sigma/z)/\prod_j\Gamma(1+(\sigma-\zeta^j)/z)$ \Lr{lem:con-form}~(d). In the variable $x=1/\sigma$ this summand, with its Stirling expansions, is $\re^{-\cW/z}\fa\,\dd x$ up to a factor independent of $x$, where, with $\Xi=-\log(1-u)$,
\begin{gather*}
\cW(x):=\frac{\Xi}x+\sum_{m\geq1}\frac{x^{5m-1}}{m(5m-1)},\\
\log\fa(x,z)=-\tfrac12\log(1-x^5)+\sum_{k\geq1}\bk_k(xz)^{2k-1}\Bigl[5^{1-2k}-5\sum_{P\geq0}\binom{2k-2+5P}{5P}x^{5P}\Bigr].
\end{gather*}
In the hatted variables $x=sX$, $z=s^4\zh$ one has $\cW/z=\hat\cW/\zh$ and $\hat\fa(X;\zh):=\fa(sX,s^4\zh)\in\bbQ[X,\zh]\fps u$, with
\[
\hat\cW(X;u)=\sum_{m\geq1}\frac{u^{m-1}}m\Bigl[\frac1X+\frac{X^{5m-1}}{5m-1}\Bigr],\qquad\hat\cW'(1)=0,\qquad\hat\cW''(1)=\frac5{1-u},\qquad\hat\cW(1)=\frac5{4\,\tau(u)} :
\]
the point $X=1$ is critical for every $u$, and by \cref{lem:set-canon} the critical exponential is $\re^{-\hat\cW(1)/\zh}=\re^{(u_0-u_0^{\mathrm c})/z}$ \Lr{lem:con-form}.

\begin{defn}[The full-Stirling WKB function and the residue; {\Lr{def:con-PhiSP}}]
\label{def:con-PhiSP}
The \emph{full-Stirling WKB function} is the normalised stationary-phase series, with constant term $1$, of $\re^{-\hat\cW/\zh}\hat\fa\,\dd X$ at $X=1$:
\[
\PhiSP(\zh,u):=\SPh_{X=1}[-\hat\cW;\hat\fa](\zh)=\sum_{n\geq0}\mathrm{ph}^{\mathrm{SP}}_n(u)\zh^n,\qquad\mathrm{ph}^{\mathrm{SP}}_0=1,\quad\mathrm{ph}^{\mathrm{SP}}_n\in\bbQ\fps u .
\]
The \emph{residue} of the flat $1$-form is $\cV:=\Res_{x=0}\re^{-\cW/z}\fa\,\dd x:=\sum_{k\geq1}\frac{(-\Xi/z)^k}{k!}\,[x^{k-1}]\bigl(\re^{-\cG/z}\fa\bigr)\in z^{-1}\bbQ[z^{-1}]\fps u$, where $\cG:=\cW-\Xi/x$.
\end{defn}

So one $1$-form carries two functionals, the stationary phase at the critical point and the residue at the pole; for the model form $p^m\re^{-\Wv/z}\dd p$ these are the model rows and the vanishing solution $\IC_m$.

%-------------------------------------------------------------------------------
\subsection{Theorem P}
\label{sec:con-P}

The ring $\AL$ carries the derivation $\Dq$ with $\Dq L=(L^6-L)/5$, for which $\Dq q=q$; let $\jmath\colon\AL\to\bbQ\lser s$ be the injective ring homomorphism with $\jmath(L)=s^{-1}$, and write $\PFfive^L$ for $\PFfive$ with $z\Dq$ replaced by $z\Dq+L$.

\begin{thm}[Theorem P(a); {\Lr{thm:con-Pa}}]\genusstatus{AI-proposed, unchecked}
\label{thm:con-Pa}
Let $c(z;L)=\sum_nc_n(L)z^n$ be the full-Stirling saddle series, with $c_n\in\AL$.
\begin{enumerate}
\item $\PFfive^{L}\bigl(L\,c(z;L)\bigr)=0$ in $\AL\fps z$.
\item In $\bbQ\lser s\fps z$ one has $\jmath(c)(z)=\PhiSP(s^{-4}z,s^5)$; that is, $\jmath(c_n)=s^{-4n}\mathrm{ph}^{\mathrm{SP}}_n$ for every $n$.
\item $\PhiSP$ solves HPF, \eqref{eq:gau-HPF}, in $\bbQ\fps u\fps{\zh}$.
\item Every $\mathrm{ph}^{\mathrm{SP}}_n$ is a polynomial in $u$; hence $c_n(L)=L^{4n}\mathrm{ph}^{\mathrm{SP}}_n(L^{-5})$ in $\AL$, and $c_n\in\bbQ[L,L^{-1}]$.
\end{enumerate}
The proof uses no input from outside [L].
\end{thm}

\begin{sketch}
(i) The series $c$ depends on $q$ only through the factor $(-q)^{\sigma/z}$ of the summand, so $z\Dq$ acts under the stationary phase as multiplication by $\sigma$: for $c\langle\chi\rangle:=\SPn_L[\mathsf S_0;\chi A_{\mathrm{St}}/A_{\mathrm{St}}(L,0)]$, $\chi\in\bbQ[\sigma,z]$, one has $(z\Dq+L)(L\,c\langle\chi\rangle)=L\,c\langle\sigma\chi\rangle$ (the derivative in a parameter, with the moving critical point $L$). Hence $\PFfive^L(L\,c)=L\bigl(c\langle\sigma^5-1\rangle-c\langle q\prod_{k=1}^5(5\sigma+kz)\rangle\bigr)$. The two terms are equal, because the translation $\sigma\mapsto\sigma+z$ leaves the stationary phase invariant and, by the functional equation of $\Gst$, turns the amplitude $(\sigma^5-1)A_{\mathrm{St}}$ into $q\prod_k(5\sigma+kz)A_{\mathrm{St}}$. (ii) is the change of variables $\sigma=L/X$, under which the phase $\mathsf S_0$ becomes $-s^4(\hat\cW-\hat\cW(1))$ and the normalised amplitude becomes $\hat\fa/\hat\fa(1;0)$. (iii) follows from (i) and (ii) by the chain rule of \cref{lem:gau-HPF}, and (iv) from \cref{lem:gau-HPF}~(iii).
\end{sketch}

% Abridged version, Section 7, part 2 of 2: Theorem P(b) and the glue.

Theorem~P(b) identifies $\PhiSP$ with $\Phican$. Both solve HPF with constant term $1$ and coefficients in $\bbQ\fps u$, so by \cref{lem:gau-HPF}~(iii) they differ by a series $\mathsf s(z)\in1+z^5\bbQ\fps{z^5}$, which has to be shown to be $1$. The equation does not determine it; it is determined by a second functional on the same $1$-form, the residue at the pole, compared with the gauge identity (G).

\begin{lem}[The uniformising map and the reduction; {\Lr{lem:con-cfu,lem:con-red}}, abridged]\genusstatus{AI-proposed, unchecked}
\label{lem:con-cfu}
There are unique $Y_k\in\bbQ[T]$ ($k\geq1$) and $V\in1+u\bbQ\fps u$ such that $Y:=X\bigl(1+\sum_{k\geq1}u^kY_k(X^5)\bigr)$ satisfies
\[
\hat\cW(X;u)=V/Y+Y^4/4=W_V(Y)\qquad\text{in }\bbQ[X,X^{-1}]\fps u.
\]
Moreover $V=\tau^{-5/4}$, and the point $X=1$ goes to the critical point $Y_c:=Y(1)=V^{1/5}=\tau^{-1/4}$ of $W_V$. Let $\fa_Y(Y;\zh):=\hat\fa(X(Y);\zh)\,X'(Y)\in\bbQ[Y,\zh]\fps u$ be the hatted amplitude in the variable $Y$. There are $\pi_m\in\bbQ[\zh]\fps u$ ($m=-1,\dots,3$) and $\mathcal E\in Y\bbQ[Y,\zh]\fps u$ with
\[
\fa_Y\,\re^{-W_V/\zh}\,\dd Y=\sum_{m=-1}^{3}\pi_m\,Y^m\re^{-W_V/\zh}\,\dd Y+\dd\bigl(\mathcal E\,\re^{-W_V/\zh}\bigr),
\]
and $\pi_m$ contains only powers $\zh^k$ with $k\equiv m\bmod5$.
\end{lem}

Here $V$ is identified by the critical value $\hat\cW(1)=\frac54\tau^{-1}=\frac54V^{4/5}$; in the flat variables $p=sY$ the parameter of the model potential is again $uV=v(u)$. The reduction uses (Mod2), in the form of the exact forms $\zh\,\dd(Y^{n-3}\re^{-W_V/\zh})$, to lower the degree in $Y$ to the five exponents $-1,\dots,3$, which are the charges $\nu_b$.

\begin{lem}[The two functionals on the reduced form; {\Lr{lem:con-S2,lem:con-id,prop:con-res}}, abridged]\genusstatus{AI-proposed, unchecked}
\label{lem:con-id}
Work in a conifold frame.
\begin{enumerate}
\item (The model identity.) For $W_1(\tilde p):=\tilde p^4/4+1/\tilde p$ and every $b\in\bbZ/5$, $\SPn_{\tilde p=1}[-W_1;\tilde p^{\nu_b}](t)=\varepsilon_b\Fs_b(t)$.
\item $\PhiSP=\sum_b\varepsilon_b\,\pi_{\nu_b}\,\tau^{(1-\nu_b)/4}\Fs_b(\tau\zh)$.
\item Let $P$ be a graded regular gauge with coefficients in $\bbQ\fps u$ from the pulled-back model to the quintic system, and let $\Phi_P$ be $\varpi_0$ times the colour-$0$ entry of the hatted row $0$ of $P\,\Qr^{\mathrm m}(v(u))$; for $P=\Pc$ this is $\Phican$. Then $\Phi_P/\varpi_0=\sum_b\tau^{(1-\nu_b)/4}\,G_b\,\Fs_b(\tau\zh)$ with $G_b:=s^{\nu_b}P_{0b}(u,s^4\zh)$, and such an expansion, with coefficients that contain only powers $\zh^k$, $k\equiv-b\bmod5$, is unique. Hence, if $\Phi_P=\PhiSP$, then $G_b=\varepsilon_b\pi_{\nu_b}/\varpi_0$ for every $b$.
\item (The residue.) If $P$ is as in (iii) and $\Phi_P=\PhiSP$, then $\cV=\varpi_0\,\bigl(P\,\YC(v(u),z)\bigr)_0$ in $\bbQ\lser z\fps u$.
\end{enumerate}
\end{lem}

\begin{sketch}
(i) The stationary phases $J_m$ of $p^m\re^{-\Wv/z}\dd p$ at the critical point $p=\varsigma$ satisfy (Mod1), by the derivative in the parameter $v$, and (Mod2), by the vanishing on exact forms; so they solve the model system, and the uniqueness of the T-system (\cref{prop:set-modrows}) identifies them. (ii) The stationary phase is invariant under the uniformising map and vanishes on the exact part, so \cref{lem:con-cfu} and (i), rescaled to $Y_c$, give the formula. (iv) The formal residue at the pole is also invariant under the uniformising map and vanishes on exact forms; so the residue of the hatted form is $\sum_m\pi_m\Res_Y[Y^m\re^{-W_V/\zh}\dd Y]$, and the model residues are the $\IC_m(v(u),z)$ up to powers of $s$. The same coefficients $\pi_{\nu_b}$, which are the entries of $P$ by (iii), now multiply the entries of the vanishing solution: the sum is $\varpi_0\sum_bP_{0b}\YC_b$.
\end{sketch}

\begin{thm}[Theorem P(b); {\Lr{thm:con-Pb}}]\genusstatus{AI-proposed, unchecked}
\label{thm:con-Pb}
In every conifold frame of \cref{def:set-frames}, $\PhiSP=\Phican$ in $\bbQ\fps u\fps{\zh}$. The proof does not use (QRR).
\end{thm}

\begin{proof}[Proof, given the lemmas]
Let $\PhiSP=\Phican\,\mathsf s(z)$ with $\mathsf s\in1+z^5\bbQ\fps{z^5}$, and put $P^{\mathrm{SP}}:=\mathsf s(z)\Pc$. By \cref{thm:gau-U} it is a graded regular gauge, and $\Phi_{P^{\mathrm{SP}}}=\mathsf s(z)\Phican=\PhiSP$. The gauge identity $\Pc\YC(v(u),z)=\YCQ$ (\cref{thm:gau-G}), multiplied by $\mathsf s(z)$, gives $(P^{\mathrm{SP}}\YC)_0=\mathsf s(z)(\YCQ)_0$, and \cref{lem:con-id}~(iv) for $P=P^{\mathrm{SP}}$ gives
\[
\cV=\varpi_0\,\mathsf s(z)\,(\YCQ)_0\qquad\text{in }\bbQ\lser z\fps u .
\]
Compare the coefficients of $u^1$. On the left only $k=1$ contributes in \cref{def:con-PhiSP}, because $\Xi=u+O(u^2)$, and $[x^0](\re^{-\cG/z}\fa)=1$; so $[u^1]\cV=-1/z$. On the right, $[u^0](\YCQ)_0=0$ and $[u^1](\YCQ)_0=-1/z$, since only $y^{(1)}_0=-\tc$ contributes at $u^1$ and $\tc=u+O(u^2)$ (\cref{lem:gau-vanq}, \Lr{lem:gau-vanq}); with $\varpi_0(0)=1$ the right side has the $u^1$-coefficient $-\mathsf s(z)/z$. Hence $\mathsf s=1$.
\end{proof}

So the freedom $\cn(z^5)$ left by the differential equation is fixed on both sides by the same normalisation, the residue at the pole: on the side of the gauge by (G), on the side of the Mellin--Barnes summand by the explicit residue $\cV$.

%-------------------------------------------------------------------------------
\subsection{The glue}
\label{sec:con-glue}

For a colour-$0$ WKB function with polynomial coefficients write $\Phi(z,1):=\sum_n\mathrm{ph}_n(1)z^n$, its value at $u=1$, where $\zh=z$.

\begin{thm}[The glue; {\Lr{thm:con-glue}}]\genusstatus{AI-proposed, unchecked}
\label{thm:con-glue}
In $\bbQ(\zeta)\fps z$, $h(z)\,\Dn_0(z)=\Phican(z,1)$; hence $\Dn_0\in1+z\bbQ\fps z$, and, by \cref{thm:con-Pb}, $\Dn_0(z)=\PhiSP(z,1)/h(z)$.
\end{thm}

\begin{proof}
By \cref{thm:gau-R0}~(ii), (iii), $A^{\mathrm{can}}=A^{\mathrm{act}}\,\Dn_0$ is an identity in $\bbQ(\zeta)[L,L^{-1}]\fps z$, with $A^{\mathrm{can}}_n=L^{4n+1}\mathrm{ph}^{\mathrm{can}}_n(L^{-5})$. Evaluation at $L=1$, that is $u=1$, $q=0$, $\zh=z$, is a ring homomorphism $\bbQ(\zeta)[L,L^{-1}]\to\bbQ(\zeta)$; it gives $A^{\mathrm{act}}(z)|_{L=1}\,\Dn_0(z)=\Phican(z,1)$. As $(1-5^5q)^{-1/5}=1+O(q)$, $A^{\mathrm{act}}_n|_{L=1}$ is the coefficient of $q^0$ of $A^{\mathrm{act}}_n$, which is $h_n$ by (QRR) (\cref{thm:set-qrr}). As $h\in1+z\bbQ\fps z$ is invertible, $\Dn_0=\Phican(z,1)/h(z)$ has rational coefficients.
\end{proof}

\begin{cor}[{\Lr{cor:con-A2}}]\genusstatus{AI-proposed, unchecked}
\label{cor:con-A2}
$\Dn=1$ if and only if $\PhiSP(z,1)=h(z)$. By \cref{thm:con-Pa}~(iv), $\PhiSP(z,1)=c(z;1)$, the value at $L=1$ of the Laurent polynomials $c_n$; so the condition reads $c(z;1)=h(z)$.
\end{cor}

The glue is the place where the value of (QRR) at colour $0$ enters. The coefficient of $z^n$ of $\Dn_0$ enters $\Fg{g}$ first at $g=\lceil(n+3)/3\rceil$ \Lr{rmk:con-genus}.


% Abridged version, Section 8: local Mellin inversion, Theorem Gamma, Theorem C.
\section{The Gamma identity and the normalisation}
\label{sec:gam}

This section proves $c(z;1)=h(z)$, which by \cref{cor:con-A2} is equivalent to $\Dn=1$. The proof is local: near a point at which exactly one Gamma factor of the Mellin--Barnes summand degenerates, only that factor is written as an Euler integral; after a change of chart the two-variable stationary phase depends holomorphically on a parameter that vanishes at the point, where a Gaussian computation evaluates it exactly. This section follows \Lr{sec:gam}.

%-------------------------------------------------------------------------------
\subsection{Local Mellin inversion}
\label{sec:gam-lmi}

\begin{lem}[The delta lemma with a source; {\Lr{lem:gam-deltasrc}}]\genusstatus{AI-proposed, unchecked}
\label{lem:gam-deltasrc}
Let $a_*\in\bbC^\times$ and $c\in\bbC$, let $P\in s^2\bbC\fps s$, and let $A\in\bbC\fps s\fps z$. In two variables $(s,m)$ consider the phase $F(s,m):=P(s)+a_*sm$, of exactly this form, and the amplitude $\re^{cm}A(s,z)$. Then $(0,0)$ is a non-degenerate critical point of $F$, and, exactly,
\[
\SPn_{(0,0)}\bigl[F;\,\re^{cm}A\bigr]=\exp\bigl(P(\delta)/z\bigr)\,A(\delta,z),\qquad\delta:=-cz/a_*,
\]
where both factors are re-expanded in $z$. For $c=0$ the series is $A(0,z)$.
\end{lem}

The covariance of this phase has no $ss$-entry, so in the Wick expansion every copy of $s$ is paired with a copy of $m$. (In this section $s$ is a local variable, not $u^{1/5}$.)

\begin{defn}[LMI data; {\Lr{def:gam-lmidata}}]
\label{def:gam-lmidata}
LMI data consist of a point $s_*\in\bbC$, a number $a_*\in\bbC^\times$, an integer $M_*$, an open disc $\bbD$ around $s_*$, a holomorphic function $\varphi$ on $\bbD$, and a series $\cF(\sigma,z)=\sum_{m\geq0}\cF_m(\sigma)z^m$ of holomorphic functions on $\bbD$, such that $\cF_0$ has no zero on $\bbD$ and $a_*+(\sigma-s_*)\varphi''(\sigma)\neq0$ for every $\sigma\in\bbD$. Put $a(\sigma):=a_*(\sigma-s_*)$, which is exactly linear in $\sigma$, and $U_*:=\bbD\setminus\{s_*\}$. For $L\in U_*$ choose a branch of $\log a$ near $L$ and put
\[
\lambda(L):=-\varphi'(L)-a_*\log a(L),\qquad S(\sigma):=\lambda(L)\,\sigma+\varphi(\sigma)+a(\sigma)\log a(\sigma)-a(\sigma),
\]
\[
A(\sigma,z):=a(\sigma)^{M_*-1/2}\,\mathrm{Poch}_{M_*}\bigl(z/a(\sigma)\bigr)\exp\Bigl(\sum_{k\geq1}\bk_k\bigl(z/a(\sigma)\bigr)^{2k-1}\Bigr)\,\cF(\sigma,z),
\]
that is, the expansion of $z^{M_*+a/z}\,\Gst(M_*+a/z)$ without its factor $(2\pi z)^{1/2}$ and without its phase factor, times $\cF$. Let $c^{\cF}(z;L)=\sum_{n\geq0}c^{\cF}_n(L)z^n:=\SPh_L\bigl[S;\,A\bigr]$. Finally put $\delta:=-M_*z/a_*$ and $\tilde\varphi(x):=\varphi(s_*+x)-\varphi(s_*)-\varphi'(s_*)x$.
\end{defn}

\begin{thm}[Theorem LMI: local Mellin inversion; {\Lr{thm:gam-lmi}}]\genusstatus{AI-proposed, unchecked}
\label{thm:gam-lmi}
Let LMI data be given.
\begin{enumerate}
\item For $L\in U_*$ every coefficient $c^{\cF}_n(L)$ depends only on $L$, not on the branch of $\log a$ or on $\lambda(L)$, and $c^{\cF}_n$ is holomorphic on $U_*$.
\item Every $c^{\cF}_n$ extends holomorphically to $L=s_*$.
\item With these extensions,
\[
c^{\cF}(z;s_*)=\exp\bigl(\tilde\varphi(\delta)/z\bigr)\,\frac{\cF(s_*+\delta,z)}{\cF(s_*,0)},\qquad\delta=-\frac{M_*z}{a_*},
\]
where both factors are re-expanded in $z$.
\end{enumerate}
\end{thm}

\begin{sketch}
\emph{Two variables.} By the formal Fubini theorem and Stirling's series as the stationary phase of Euler's integral (\cref{sec:con-fsp}), $c^{\cF}(z;L)$ is the normalised stationary phase, at the point $(L,\log a(L))$, of the two-variable integrand with phase $\lambda\sigma+\varphi(\sigma)+a(\sigma)\chi-\re^\chi$ and amplitude $\re^{M_*\chi}\cF(\sigma,z)$: the inner integral over $\chi$ is Euler's integral for $\Gamma(M_*+a/z)$ in the variable $\re^\chi$.

\emph{The chart.} Translate $\chi=\mu-\lambda/a_*$. Because $a$ is exactly linear, the phase becomes a constant plus $\Psi_{\mathsf Q}(\sigma,\mu):=\varphi(\sigma)+a(\sigma)\mu-\mathsf Q\,\re^\mu$, with the chart parameter $\mathsf Q(L):=\re^{-\lambda(L)/a_*}=a(L)\,\re^{\varphi'(L)/a_*}$, a single-valued holomorphic function on $\bbD$ with a simple zero at $s_*$. For every $L\in\bbD$, including $L=s_*$, the point $(L,-\varphi'(L)/a_*)$ is a non-degenerate critical point of $\Psi_{\mathsf Q(L)}$, and the Taylor data there are holomorphic in $L$; by the universality of the coefficients, $c^{\cF}_n$ is holomorphic on $\bbD$. This is (ii): the critical point, which in the variable $\sigma$ alone runs into the singularity of the degenerate factor, survives in the chart.

\emph{The value.} At $L=s_*$, $\mathsf Q=0$ and the phase is $\varphi(s_*)+\tilde\varphi(s)+a_*sm$ in the variables $s=\sigma-s_*$, $m=\mu-\mu_*$, exactly: the exponential term has disappeared. \Cref{lem:gam-deltasrc} with $P=\tilde\varphi$, $c=M_*$ and $A=\cF(s_*+s,z)$ gives (iii).
\end{sketch}

%-------------------------------------------------------------------------------
\subsection{Theorem Gamma and the value of \texorpdfstring{$\Dn$}{D}}
\label{sec:gam-thm}

\begin{thm}[Theorem Gamma; {\Lr{thm:gam-main,thm:gam-ginf}}, abridged]\genusstatus{AI-proposed, unchecked}
\label{thm:gam-main}
Let $c(z;L)$ be the saddle series of \cref{def:gam-c}.
\begin{enumerate}
\item Every $c_n$ lies in $\AL$.
\item No $c_n$ has a pole at $L=1$, that is at $q=0$, and
\[
c(z;1)=h(z),\qquad\text{that is, }c_n(1)=h_n\ \text{for every }n\geq0.\tag{$\Gamma$}
\]
\item $c_n(\zeta L)=\zeta^{-n}c_n(L)$ for every $n$. Hence no $c_n$ has a pole at a fifth root of unity, $c_n$ lies in $\bbQ[L,L^{-1}]$, and $c(z;\zeta^i)=h_{\bw_i}(z)=h(\zeta^{-i}z)$.
\item No $c_n$ has a pole at $L=0$, that is at $q=\infty$, and $c(z;0)=\Ginf(z)$, the series of \cref{thm:int-gamma}~(b), which is the series displayed on the right of (37) of \cite[Lemma~7.7]{GJR18} at $m=5$, $k=1$ and equivariant parameter $1$ in the notation there, read as a series in $z$; the comparison with \cite{GJR18} enters no proof. Consequently $c_n\in\bbQ[L]$.
\end{enumerate}
Only (ii) is used in the proof of the main theorem.
\end{thm}

\begin{proof}[Proof of (ii), given \cref{thm:gam-lmi}]
Apply \cref{thm:gam-lmi} with $s_*=1$, $a_*=-1$, so that $a(\sigma)=1-\sigma$, and $M_*=0$, with
\[
\varphi(\sigma):=5\sigma\log(5\sigma)-5\sigma+\sum_{j=1}^4\bigl[(\zeta^j-\sigma)\log(\zeta^j-\sigma)-(\zeta^j-\sigma)\bigr],
\]
\[
\cF(\sigma,z):=(5\sigma)^{1/2}\prod_{j=1}^4(\zeta^j-\sigma)^{-1/2}\exp\Bigl(\sum_{k\geq1}\bk_kz^{2k-1}\Bigl[(5\sigma)^{1-2k}+\sum_{j=1}^4(\zeta^j-\sigma)^{1-2k}\Bigr]\Bigr),
\]
the phase and the amplitude of the five factors of the summand that do not degenerate at $\sigma=1$; on a small disc $\bbD$ around $1$, with fixed branches, the conditions of \cref{def:gam-lmidata} hold. The term $a\log a-a$ is the term $j=0$ of $\mathsf S_0$, and, as $\mathrm{Poch}_0=1$, the Stirling form of $\Gamma(a/z)$ contributes the factor $j=0$ of $A_{\mathrm{St}}$. Hence $S=\mathsf S_0$, $A=A_{\mathrm{St}}$, $\lambda(L)$ is the logarithm of \cref{def:gam-c}, and $c^{\cF}(z;L)=c(z;L)$ for $L\in U_*$. By \cref{thm:gam-lmi}~(ii) every $c_n$ extends holomorphically to $L=1$; since $c_n$ is rational by (i), it has no pole at $L=1$. By \cref{thm:gam-lmi}~(iii) with $\delta=0$, and since the first factor of $\cF$ does not involve $z$,
\[
c(z;1)=\frac{\cF(1,z)}{\cF(1,0)}=\exp\Bigl(\sum_{k\geq1}\bk_k\Bigl[5^{1-2k}+\sum_{j=1}^4(\zeta^j-1)^{1-2k}\Bigr]z^{2k-1}\Bigr)=h(z),
\]
because $1-2k$ is odd and $\bw_0=\{1-\zeta^j\}_{j=1,\dots,4}\cup\{-5\}$.
\end{proof}

So at $q=0$ the saddle series is the product of the Stirling series of the five factors that stay regular, evaluated at $\sigma=1$: these are the quantum Riemann--Roch factors of the four tangent weights $1-\zeta^j$ and of the twist weight $-5$. Part (iii) is the rotation $\sigma\mapsto\zeta\sigma$ of the Taylor data, and (iv) is \cref{thm:gam-lmi} at $s_*=0$ with $a(\sigma)=5\sigma$ and $M_*=1$, where only the twist factor degenerates and $\delta=-z/5\neq0$ \Lr{thm:gam-ginf}; with \cref{thm:con-Pb} it identifies the free coefficients of \cref{lem:gau-HPF}, $1+\sum_{m\geq1}\fc^{\mathrm{can}}_mz^{5m}=\Ginf(z)$ \Lr{cor:gam-free}.

\begin{thm}[Theorem C: the normalisation is canonical; {\Lr{thm:gam-D,prop:gam-hatflat}}]\genusstatus{AI-proposed, unchecked}
\label{thm:gam-D}
In every conifold frame, $\Dn_0=1$ in $\bbQ\fps z$, and $\Dn=1$. Consequently $\Fg{g}$ holds for every $g\geq2$. Moreover $\PhiSP(\zh,1)=h(\zh)$.
\end{thm}

\begin{proof}
By \cref{thm:con-glue}, with \cref{thm:con-Pb}, $\Dn_0(z)=\PhiSP(z,1)/h(z)$ exactly in $\bbQ\fps z$. By \cref{thm:con-Pa}~(iv), $c_n(L)=L^{4n}\mathrm{ph}^{\mathrm{SP}}_n(L^{-5})$, so $\mathrm{ph}^{\mathrm{SP}}_n(1)=c_n(1)=h_n$ by \cref{thm:gam-main}~(ii), and $\PhiSP(z,1)=h(z)$. Hence $\Dn_0=1$, and $\Dn_i(z)=\Dn_0(\zeta^{-i}z)=1$ for every colour $i$ (\cref{thm:gau-R0}~(iv)).
\end{proof}

The value of (QRR) enters this proof through the glue, at colour $0$, and through the colour covariance of the rows, that is (QRR) for the colours $i\neq0$; the rows and their WKB form rest on (M), (G0) and the flatness part of (QRR), and the $R$-matrix identified by (T) is used only through $R|_{z=0}=1$ and (QRR). With \cref{thm:gau-R0}~(iv) it gives \cref{thm:int-D}: $\Qr=\Pc\Qr^{\mathrm m}(v(u))$, the actual gauge is the canonical gauge \Lr{cor:mai-actual}. This is not circular, as the rows are normalised at $u=1$ and the canonical gauge at $u=0$.

The mechanism is that of \cite[\S10]{MR1901075}; we are not aware of a proof of the identity $(\Gamma)$ in the literature \Lr{rmk:gam-prior}.

% Abridged version, Section 9: the main theorem, Corollary B, complements.
\section{The main theorem}
\label{sec:mai}

A conifold frame of \cref{def:set-frames} is fixed, and $M=\Mo(\tau(u)\zh)$ is taken exactly in $u$. This section follows \Lr{sec:mai,sec:cmp}.

\begin{thm}[Theorem A; {\Lr{thm:mai-main}}]\genusstatus{AI-proposed, unchecked}
\label{thm:mai-main}
The statement of \cref{thm:int-main} holds: for every $g\geq2$, every conifold frame and every leg-free stable graph $\Gamma$ of genus $g$, the outer term $T_\Gamma$ of \cref{def:set-outer} lies in $\bbQ\fps u$ if $\Gamma$ has at least one edge, and in $c_g\,\tc^{2-2g}+\bbQ\fps u$ if $\Gamma$ is the vertex graph. The hatted terms are $G_\Gamma=u^{2g-2}T_\Gamma$.
\end{thm}

\begin{proof}[Proof, given the results of \cref{sec:gau,sec:tra,sec:mod,sec:flh,sec:con,sec:gam}]
Fix $g\geq2$ and the frame. By \cref{thm:gam-D}, $\Dn_0=1$ in $\bbQ\fps z$; in particular $\Fg{g}$ holds. By the comparison theorem, \cref{thm:flh-main}, $T_\Gamma=F_\Gamma$ for every leg-free stable graph $\Gamma$ of genus $g$, where $F_\Gamma$ is the term of $\Gamma$ in the flat sum of $\Pc.\Omega^{\mathrm{mod}}_{v(u)}$. By \cref{thm:mod-MLK} the Model Lemma $\ML{h}$ holds for every $h\geq0$. So \cref{thm:red-main} applies to $\Pc$: since $\ML{h}$ holds for $h\leq g-1$, $F_\Gamma$ lies in $\bbQ\fps u$ for every $\Gamma$ with an edge; since $\ML{g}$ holds as well, the vertex graph gives $F_\Gamma\in c_g\tc^{2-2g}+\bbQ\fps u$, the leaf point of $\Pc$ having the vanishing coordinate $T_{\Pc}=\tc$ (\cref{prop:tra-can}). The last sentence is \cref{lem:flh-hat}~(iv).
\end{proof}

A priori $T_\Gamma\in\bbQ(\zeta)\lser s$. Theorem~A, and Corollary~B below, hold modulo (Q), (QRR), (M), (G0), (T) and \cite[(1.7)]{AC96} (\cref{sec:int-inputs}).

\emph{The conifold expansion of $F_g$.} To expand $F_g$ at the conifold point we use its graph-sum representative in $\Nring$. In $\Nring\otimes\bbQ(\zeta)$ the constant term $\Psi_{\mathrm Q}:=\Qr|_{z=0}$ has the entries $(\Psi_{\mathrm Q})_{aj}=\zeta^{j(a+1)}\mathfrak m_a$ with $(\mathfrak m_0,\dots,\mathfrak m_4)=\bigl(L/C_0,\,L^2/(C_0C_1),\,C_0C_1/L^2,\,C_0/L,\,1\bigr)$, and $\Psi_{\mathrm Q}^T\eta^{-1}\Psi_{\mathrm Q}=1$ \Lr{lem:mai-Psi}.

\begin{defn}[The genus-$g$ sum in the ring and its conifold expansion; {\Lr{def:mai-Fcf}}]
\label{def:mai-Fcf}
Put $\tilde R:=\Psi_{\mathrm Q}^{-1}\Qr=\Psi_{\mathrm Q}^T\eta^{-1}\Qr$, the $R$-matrix in Teleman's colour-colour form, a series over $\Nring\otimes\bbQ(\zeta)$ with $\tilde R(0)=1$. For a matrix series $X(z)$ with $X(0)=1$ let $N_X(\psi,\psi'):=1-X^{-1}(\psi)X^{-1}(\psi')^T$ and
\[
E_X(\psi,\psi'):=\frac{N_X(\psi,\psi')-\frac12N_X(-\psi',\psi')-\frac12N_X(\psi,-\psi)}{\psi+\psi'} ,
\]
a power series for every $X$, which is the edge of \cite[Def.~2.2]{PPZ15} when $X(-z)^TX(z)=1$. With the leaf point $q_{\tilde R}=\tilde R^{-1}(z)(-z1_{\TFT})$, define
\[
F_g^{\Nring}:=\sum_\Gamma\mathrm{Cont}_\Gamma\bigl[E_{\tilde R},q_{\tilde R};\TFT\bigr]\in\Nring\otimes\bbQ(\zeta),
\]
the sum over the leg-free stable graphs of genus $g$. The \emph{conifold expansion} of $F_g$ in the frame $(\alpha,\beta)$ is $F_g^{\mathrm{cf}}:=\Cexp_{\alpha,\beta}\bigl(F_g^{\Nring}\bigr)\in\bbQ\lser s$.
\end{defn}

The definition is justified by three facts \Lr{def:mai-Fcf}: $F_g^{\Nring}$ is a polynomial over $\bbQ$ in finitely many row coefficients, in $\Delta_i^{\pm1/2}$ and in $\zeta$; it lies in $\Nring$, because the Galois group of $\bbQ(\zeta)$ permutes the colours; and $\ev(F_g^{\Nring})=F_g$, because under $\ev$ the matrix $\tilde R$ becomes the symplectic $R$-matrix of \cref{def:set-rows}, so that the graph sum is the potential by (T), with the normalisation (QRR). The symplecticity of $\tilde R$ is not known in $\Nring\otimes\bbQ(\zeta)$.

\begin{cor}[Corollary B: the conifold gap of the formal quintic; {\Lr{cor:mai-gap}}]\genusstatus{AI-proposed, unchecked}
\label{cor:mai-gap}
Let $g\geq2$. In every conifold frame of \cref{def:set-frames}, $F_g^{\mathrm{cf}}\in c_g\,\tc^{2-2g}+\bbQ\fps u$.
\end{cor}

\begin{sketch}
The conifold expansion is a ring homomorphism, and the edge $E_X$ commutes with ring homomorphisms; so $F_g^{\mathrm{cf}}$ is the graph sum of $\Cexp(\tilde R)$ on $\Cexp(\TFT)$. Since $\Dn=1$ (\cref{thm:gam-D}), $R'$ is symplectic (\cref{lem:flh-inputs}~(iii)), hence so are $\Cexp(\Qr)=R'M$ and $\Cexp(\tilde R)$: after the conifold expansion the graph sum is the integral over $\overline{\mathcal M}_g$ of the genus-$g$ class of the unit-preserving action $\Cexp(\tilde R).\Cexp(\TFT)$. Now $X:=\Cexp(\tilde R)M^{-1}$ is symplectic with $X(0)=1$, and by the group property of the action \cite[Prop.~2.14]{PPZ15}, $\Cexp(\tilde R).\Cexp(\TFT)=X.\bigl(M.\Cexp(\TFT)\bigr)$. The genus-$g$ part without legs of the right side is the outer graph sum, $\sum_\Gamma T_\Gamma$ (\cref{def:set-outer}). \Cref{thm:mai-main} gives the claim, the sum over $\Gamma$ being finite.
\end{sketch}

 

\emph{Complements} \Lr{sec:cmp}. After $\Dn=1$ the hatted rows factorise as $\hat\Qr(u,\zh)=\Ghat(u,\zh)\,\Rhat^{\mathrm m}(\tau(u)\zh)$ with a hatted gauge $\Ghat$ regular at $u=0$, so that the conifold symbol $\hat\Qr(0,\zh)$ of the formal quintic is that of the model \Lr{cor:mai-natfact}; and the flat half of the proof applies, suitably normalised, to every self-dual one-pole system over $\bbQ\fps u$ whose product is that of a pulled-back model \Lr{cor:mai-univ}.

% Abridged version, Section 10: the finite computations that accompanied the work, and what remains conditional.
\section{Evidence and conditionality}
\label{sec:evi}
\noindent\genusstatus{numerically supported} The finite computations recorded in this section are exact or modular computations; they support the statements at the genera, orders and primes given, and are not proofs. The code and outputs are in \texttt{code/} of this project.


\emph{Finite computations.} Nothing in this paragraph and the next two is a step of a proof: the proofs of [L] are complete without it. We summarise the tables of \Lr{sec:evi}; a zero test modulo a prime is probabilistic. The outer terms $T_\Gamma=u^{2-2g}G_\Gamma(u)$, with $M=\Mo(\tau(u)\zh)$, were computed for every leg-free stable graph of genus $2$ to $5$, modulo $2^{89}-1$ and $2^{61}-1$ with rational reconstruction, in the frames $(\alpha,\beta)=(0,0)$ and $(1/3,-2/7)$. There are $7$, $42$, $379$ and $4555$ stable graphs; all $6$, $41$, $378$ and $4554$ edge graphs are regular (at genus $2$, $3$ and $5$ in the second frame modulo one prime), and the vertex graph, respectively the total, has the polar part $c_g\tc^{2-2g}$ with $c_g=\frac1{1200},\frac1{25200},\frac1{180000},\frac1{660000}$, in the frame $(0,0)$ and, for $g=3,4,5$, also in the second frame. At genus $4$ and $5$ the outer sum was obtained by regrouping the plain graph sum along the edges of $M$, a regrouping checked against genuine vertex tensors only at genus $2$ and $3$. As controls: with the bare symbol $\Mo(\zh)$ all $149$ edge graphs of genus $4$ with at most five edges are singular, and with $\tau$ replaced by $\tau+u^5$ exactly the three one-edge graphs of genus $4$ fail. The value $c_g$ and the gap of the total do not depend on $M$; only the regularity of the edge graphs tests the factorisation \Lr{sec:evi-graphs}.

The identity $\Dn_0=1$ was tested order by order at the conifold point up to $z^{204}$ (that is, $g\leq69$) exactly and up to $z^{1004}$ modulo two primes; the identity $(\Gamma)$, $c_n(1)=h_n$, for $n\leq60$ exactly and up to $z^{500}$ modulo two primes; $\PhiSP=\Phican$ up to $\zh^{80}$ and $u^{75}$ modulo primes; and $c(z;0)=\Ginf$ up to $z^{80}$ exactly \Lr{sec:evi-normal}.

For the model, \cref{thm:mod-tr}, the values of $\FTR_h$ for $h\leq5$, the lattice theorem and the Model Lemma were checked on finite ranges; off $\sK$ the law $\langle X\rangle_h=\Pi_0(X)(\Fm_h)'$ for singular one-point correlators, which holds on $\sK$ by $\ML{h}$, fails in an example \Lr{cor:mod-onepoint,sec:evi-model}.

\emph{What remains conditional.} Theorems~A and~C and Corollary~B are conditional on the inputs of \cref{sec:int-inputs}, and on nothing else: the quoted theorems (Q) on topological recursion, used for the model curve only, and on the action of $R$-matrices; the flatness of the $R$-matrix of the formal quintic and its value at the large-radius point (QRR); the genus-zero mirror theorem (M) and the genus-zero axioms (G0) of the twisted theory; Teleman's reconstruction with the axioms of the twisted theory in all genera (T); and the formula \cite[(1.7)]{AC96}. They are quoted with their statement numbers and are read with the conventions fixed in [L], where the sign conventions of the sources are compared \Lr{sec:sgn}. Theorem~Gamma and Theorem~U use no quoted input. Every other step is proved in [L]; in this abridged version the proofs are sketched, except where a proof is given in full from the stated lemmas. The limits of the results are those of \cref{sec:int-notclaimed}; in particular the conifold expansion is a formal change of frame, with no analytic continuation, and nothing is claimed for the quintic threefold itself.


\bibliographystyle{amsalpha}
\bibliography{refs}

\end{document}
